\section*{Appendice}

\paragraph*{Configurazione CI (\texttt{.github/workflows/ci.yml}).}
\label{cmd:ci-config}
\begin{lstlisting}
name: Java CI
on:
  push:
    branches: [ release-4.1.1 ]
jobs:
  build:
    runs-on: ubuntu-latest
    steps:
      - uses: actions/checkout@v4
      - name: Setup Java JDKs
        uses: actions/setup-java@v4
        with:
          distribution: 'temurin'
          java-version: |
            8
            17
      - name: Build with Maven
        run: mvn -B install
      - name: Test LoginDialog (GWTTestCase)
        run: mvn -f openjpa-examples/opentrader/pom.xml -B test
        env:
          JAVA8_HOME: ${{ env.JAVA_HOME_8_X64 }}
\end{lstlisting}

\begin{table}[H]
\centering
\begin{tabularx}{\textwidth}{X l}
\toprule
Metodo & Interfaccia / classe \\
\midrule
\texttt{find(Object oid, boolean validate, FindCallbacks call)} & \texttt{StoreContext} \\
\texttt{lock(Object pc, int level, int timeout, OpCallbacks call)} & \texttt{Broker} \\
\texttt{detachAll(Collection objs, OpCallbacks call)} & \texttt{Broker} \\
\texttt{attachAll(Collection objs, boolean copyNew, OpCallbacks call)} & \texttt{Broker} \\
\texttt{newInstance(Class cls)} & \texttt{Broker} \\
\texttt{isDetached(Object obj, boolean find)} & \texttt{BrokerImpl} \\
\texttt{persist(Object pc, Object id, OpCallbacks call)} & \texttt{Broker} \\
\bottomrule
\end{tabularx}
\caption{Metodi di \texttt{BrokerImpl} selezionati per Category Partition.}
\label{tab:brokerimpl-metodi}
\end{table}

\begin{table}[H]
\centering
\begin{tabularx}{\textwidth}{>{\raggedright\arraybackslash}p{0.38\textwidth} >{\raggedright\arraybackslash}X}
\toprule
Parametro & Classi di equivalenza \\
\midrule
\texttt{oid} & \{null\} -- \{non valido\} -- \{valido\} \\
\texttt{cache}* & \{presente\} -- \{non presente\} \\
\texttt{store}* & \{presente\} -- \{non presente\} \\
\midrule
\texttt{validate} & \{true\} -- \{false\} \\
\midrule
\texttt{call} & \{null\} -- \{non null\} \\
\texttt{call.processArgument}** & \{ritorna un id non-null\} -- \{ritorna null\} -- \{lancia
\texttt{OpenJPAException}\} \\
\texttt{call.processReturn}** & \{ritorna un oggetto non-null\} -- \{ritorna null\} \\
\bottomrule
\end{tabularx}
\caption{Classi di equivalenza per \texttt{find(Object oid, boolean validate, FindCallbacks call)}.
* richiede un'identità effettiva (\texttt{oid} valido con \texttt{call} null, oppure un id
sostituito con successo da \texttt{processArgument}). ** richiede \texttt{call} non null.}
\label{tab:find-classi}
\end{table}

\begin{table}[H]
\centering
\begin{tabularx}{\textwidth}{>{\raggedright\arraybackslash}p{0.38\textwidth} >{\raggedright\arraybackslash}X}
\toprule
Parametro & Classi di equivalenza \\
\midrule
\texttt{pc} & \{null\} -- \{non gestito\} -- \{gestito, non ancora persistito\} -- \{gestito,
presente nello store\} \\
\texttt{level} & \{valore noto\} -- \{valore non noto\} \\
\texttt{timeout} & \{altro negativo ($<-1$)\} -- \{-1\} -- \{0\} -- \{positivo\} \\
\midrule
\texttt{call} & \{null\} -- \{non null\} \\
\texttt{call.processArgument}* & \{ritorna un codice azione\} -- \{lancia \texttt{OpenJPAException}\} \\
\bottomrule
\end{tabularx}
\caption{Classi di equivalenza per \texttt{lock(Object pc, int level, int timeout,
OpCallbacks call)}. * richiede \texttt{call} non null.}
\label{tab:lock-classi}
\end{table}

\begin{table}[H]
\centering
\begin{tabularx}{\textwidth}{>{\raggedright\arraybackslash}p{0.38\textwidth} >{\raggedright\arraybackslash}X}
\toprule
Parametro & Classi di equivalenza (corrette) \\
\midrule
\texttt{pc} & \{null\} -- \{non gestito\} -- \{gestito, non ancora persistito\} -- \{gestito,
presente nello store\} \\
\texttt{level} & \{valore noto\} -- \{valore non noto\} \\
\texttt{timeout} & \{altro negativo ($<-1$)\} -- \{-1\} -- \{0\} -- \{positivo\} \\
\midrule
\texttt{call} & \{null\} -- \{non null\} \\
\texttt{call.processArgument}* & \{ritorna \texttt{ACT\_RUN}\} -- \{ritorna \texttt{ACT\_NONE}/
\texttt{ACT\_CASCADE}\} -- \{lancia \texttt{OpenJPAException}\} \\
\bottomrule
\end{tabularx}
\caption{Classi di equivalenza corrette per \texttt{lock(Object pc, int level, int timeout,
OpCallbacks call)} (confronto con Tabella~\ref{tab:lock-classi}). * richiede \texttt{call} non null.}
\label{tab:lock-classi-corrette}
\end{table}

\begin{table}[H]
\centering
\begin{tabularx}{\textwidth}{>{\raggedright\arraybackslash}p{0.38\textwidth} >{\raggedright\arraybackslash}X}
\toprule
Parametro & Classi di equivalenza \\
\midrule
\texttt{objs} & \{null\} -- \{vuota\} -- \{un elemento gestito\} -- \{un elemento non gestito\}
-- \{elementi misti (un gestito e un non gestito)\} \\
\midrule
\texttt{call} & \{null\} -- \{non null\} \\
\texttt{call.processArgument}* & \{ritorna un codice azione\} -- \{lancia \texttt{OpenJPAException}\} \\
\bottomrule
\end{tabularx}
\caption{Classi di equivalenza per \texttt{detachAll(Collection objs, OpCallbacks call)}.
* richiede \texttt{call} non null.}
\label{tab:detachall-classi}
\end{table}

\begin{table}[H]
\centering
\begin{tabularx}{\textwidth}{>{\raggedright\arraybackslash}p{0.38\textwidth} >{\raggedright\arraybackslash}X}
\toprule
Parametro & Classi di equivalenza (corrette) \\
\midrule
\texttt{objs} & \{null\} -- \{vuota\} -- \{un elemento gestito\} -- \{un elemento non gestito\}
-- \{elementi misti (un gestito e un non gestito)\} \\
\midrule
\texttt{call} & \{null\} -- \{non null\} \\
\texttt{call.processArgument}* & \{ritorna \texttt{ACT\_RUN}\} -- \{ritorna \texttt{ACT\_NONE}/
\texttt{ACT\_CASCADE}\} -- \{lancia \texttt{OpenJPAException}\} \\
\bottomrule
\end{tabularx}
\caption{Classi di equivalenza corrette per \texttt{detachAll(Collection objs, OpCallbacks
call)} (confronto con Tabella~\ref{tab:detachall-classi}). * richiede \texttt{call} non null.}
\label{tab:detachall-classi-corrette}
\end{table}

\begin{table}[H]
\centering
\begin{tabularx}{\textwidth}{>{\raggedright\arraybackslash}p{0.42\textwidth} >{\raggedright\arraybackslash}X}
\toprule
Parametro & Classi di equivalenza \\
\midrule
\texttt{objs} & \{null\} -- \{vuota\} -- \{all new\} -- \{all old\} -- \{all attached\} --
\{misti: validi eterogenei\} -- \{misti: almeno 1 invalido\} \\
\texttt{copyNew} & \{true\} -- \{false\} \\
\midrule
\texttt{call} & \{null\} -- \{non null\} \\
\texttt{call.processArgument}* & \{ritorna un codice azione\} -- \{lancia \texttt{OpenJPAException}\} \\
\bottomrule
\end{tabularx}
\caption{Classi di equivalenza per \texttt{attachAll(Collection objs, boolean copyNew,
OpCallbacks call)}. * richiede \texttt{call} non null. Definizioni precise delle classi di
\texttt{objs} in Sezione~\ref{sec:cp}.}
\label{tab:attachall-classi}
\end{table}

\begin{table}[H]
\centering
\begin{tabularx}{\textwidth}{>{\raggedright\arraybackslash}p{0.42\textwidth} >{\raggedright\arraybackslash}X}
\toprule
Parametro & Classi di equivalenza (corrette) \\
\midrule
\texttt{objs} & \{null\} -- \{vuota\} -- \{all new\} -- \{all old\} -- \{all attached\}* --
\{misti: validi eterogenei\} -- \{misti: almeno 1 invalido\} \\
\texttt{copyNew} & \{true\} -- \{false\} \\
\midrule
\texttt{call} & \{null\} -- \{non null\} \\
\texttt{call.processArgument}** & \{ritorna \texttt{ACT\_RUN}\} -- \{ritorna \texttt{ACT\_NONE}\} --
\{ritorna \texttt{ACT\_CASCADE}\} -- \{lancia \texttt{OpenJPAException}\} \\
\bottomrule
\end{tabularx}
\caption{Classi di equivalenza corrette per \texttt{attachAll(Collection objs, boolean copyNew,
OpCallbacks call)} (confronto con Tabella~\ref{tab:attachall-classi}). * \{all attached\} resta
una classe distinta, ma il suo oracolo osservato coincide con quello di \{all new\}. ** richiede
\texttt{call} non null.}
\label{tab:attachall-classi-corrette}
\end{table}

\begin{table}[H]
\centering
\begin{tabularx}{\textwidth}{>{\raggedright\arraybackslash}p{0.38\textwidth} >{\raggedright\arraybackslash}X}
\toprule
Parametro & Classi di equivalenza \\
\midrule
\texttt{cls} & \{null\} -- \{interfaccia con metodi JavaBeans\} -- \{interfaccia senza metodi
JavaBeans\} -- \{classe astratta con metodi JavaBeans\} -- \{classe astratta senza JavaBeans\}
-- \{tipo managed concreto\} -- \{tipo non managed concreto\} \\
\bottomrule
\end{tabularx}
\caption{Classi di equivalenza per \texttt{newInstance(Class cls)}.}
\label{tab:newinstance-classi}
\end{table}

\begin{table}[H]
\centering
\begin{tabularx}{\textwidth}{>{\raggedright\arraybackslash}p{0.38\textwidth} >{\raggedright\arraybackslash}X}
\toprule
Parametro & Classi di equivalenza \\
\midrule
\texttt{obj} & \{null\} -- \{gestito\} -- \{non gestito\} \\
\texttt{find} & \{true\} -- \{false\} \\
\texttt{store}* & \{presente\} -- \{non presente\} \\
\bottomrule
\end{tabularx}
\caption{Classi di equivalenza per \texttt{isDetached(Object obj, boolean find)}. * non si
applica solo quando \texttt{obj}=\{null\}.}
\label{tab:isdetached-classi}
\end{table}

\begin{table}[H]
\centering
\begin{tabularx}{\textwidth}{>{\raggedright\arraybackslash}p{0.38\textwidth} >{\raggedright\arraybackslash}X}
\toprule
Parametro & Classi di equivalenza \\
\midrule
\texttt{pc} & \{null\} -- \{non gestito\} -- \{gestito\} \\
\texttt{id} & \{null\} -- \{non null\} \\
\texttt{cascade}* & \{presente\} -- \{assente\} \\
\midrule
\texttt{call} & \{null\} -- \{non null\} \\
\texttt{call.processArgument}** & \{ritorna un codice azione\} -- \{lancia \texttt{OpenJPAException}\} \\
\bottomrule
\end{tabularx}
\caption{Classi di equivalenza per \texttt{persist(Object pc, Object id, OpCallbacks call)}.
* non si applica per \texttt{pc}=\{null\}. ** richiede \texttt{call} non null.}
\label{tab:persist-classi}
\end{table}

\begin{table}[H]
\centering
\begin{tabularx}{\textwidth}{>{\raggedright\arraybackslash}p{0.38\textwidth} >{\raggedright\arraybackslash}X}
\toprule
Parametro & Classi di equivalenza (corrette) \\
\midrule
\texttt{pc} & \{null\} -- \{non gestito\} -- \{gestito\} \\
\texttt{id} & \{null\} -- \{non null\} \\
\texttt{cascade}* & \{presente\} -- \{assente\} \\
\midrule
\texttt{call} & \{null\} -- \{non null\} \\
\texttt{call.processArgument}** & \{ritorna \texttt{ACT\_RUN}\} -- \{ritorna \texttt{ACT\_NONE}/
\texttt{ACT\_CASCADE}\} -- \{lancia \texttt{OpenJPAException}\} \\
\bottomrule
\end{tabularx}
\caption{Classi di equivalenza corrette per \texttt{persist(Object pc, Object id, OpCallbacks
call)} (confronto con Tabella~\ref{tab:persist-classi}). * non si applica per \texttt{pc}=\{null\}.
** richiede \texttt{call} non null.}
\label{tab:persist-classi-corrette}
\end{table}

\begin{table}[H]
\centering
\begin{tabularx}{\textwidth}{>{\raggedright\arraybackslash}p{0.38\textwidth} >{\raggedright\arraybackslash}X}
\toprule
Parametro & Classi di equivalenza \\
\midrule
\texttt{session} & \{null\} -- \{valida\} \\
\bottomrule
\end{tabularx}
\caption{Classi di equivalenza per il costruttore \texttt{LoginDialog(OpenTrader session)}.}
\label{tab:logindialog-costruttore-classi}
\end{table}

\begin{table}[H]
\centering
\begin{tabularx}{\textwidth}{>{\raggedright\arraybackslash}p{0.38\textwidth} >{\raggedright\arraybackslash}X}
\toprule
Parametro & Classi di equivalenza \\
\midrule
\texttt{left} & \{negativo\} -- \{zero\} -- \{positivo\} \\
\texttt{top} & \{negativo\} -- \{zero\} -- \{positivo\} \\
\bottomrule
\end{tabularx}
\caption{Classi di equivalenza per \texttt{setPopupPosition(int left, int top)} (ereditato da
\texttt{PopupPanel}).}
\label{tab:logindialog-posizione-classi}
\end{table}

\begin{table}[H]
\centering
\begin{tabularx}{\textwidth}{>{\raggedright\arraybackslash}p{0.38\textwidth} >{\raggedright\arraybackslash}X}
\toprule
Parametro & Classi di equivalenza \\
\midrule
\texttt{title} & \{null\} -- \{vuota\} -- \{valida\} \\
\bottomrule
\end{tabularx}
\caption{Classi di equivalenza per \texttt{setTitle(String title)} (ereditato da
\texttt{PopupPanel}).}
\label{tab:logindialog-title-classi}
\end{table}

\begin{table}[H]
\centering
\begin{tabularx}{\textwidth}{>{\raggedright\arraybackslash}p{0.38\textwidth} >{\raggedright\arraybackslash}X}
\toprule
Proprietà & Classi di equivalenza \\
\midrule
\texttt{Modal} & \{impostato a \texttt{true}\} -- \{impostato a \texttt{false}\} \\
\texttt{AnimationEnabled} & \{impostato a \texttt{true}\} -- \{impostato a \texttt{false}\} \\
\texttt{AutoHideEnabled} & \{impostato a \texttt{true}\} -- \{impostato a \texttt{false}\} \\
\texttt{GlassEnabled} & \{impostato a \texttt{true}\} -- \{impostato a \texttt{false}\} \\
\texttt{Visible} & \{impostato a \texttt{true}\} -- \{impostato a \texttt{false}\} \\
\bottomrule
\end{tabularx}
\caption{Classi di equivalenza per le proprietà booleane ereditate da \texttt{PopupPanel}
(\texttt{setModal}/\texttt{isModal}, \texttt{setAnimationEnabled}/\texttt{isAnimationEnabled},
\texttt{setAutoHideEnabled}/\texttt{isAutoHideEnabled}, \texttt{setGlassEnabled}/
\texttt{isGlassEnabled}, \texttt{setVisible}/\texttt{isVisible}).}
\label{tab:logindialog-boolprop-classi}
\end{table}

\begin{table}[H]
\centering
\begin{tabularx}{\textwidth}{>{\raggedright\arraybackslash}p{0.38\textwidth} >{\raggedright\arraybackslash}X}
\toprule
Parametro & Classi di equivalenza \\
\midrule
\texttt{stato\_precedente} & \{non mostrato\} -- \{mostrato\} \\
\texttt{operazione} & \{\texttt{show()}\} -- \{\texttt{hide()}\} -- \{\texttt{hide(true)}\} --
\{\texttt{center()}\} \\
\bottomrule
\end{tabularx}
\caption{Classi di equivalenza per il ciclo di vita \texttt{show}/\texttt{hide}/\texttt{center}/
\texttt{isShowing} (ereditato da \texttt{PopupPanel}).}
\label{tab:logindialog-lifecycle-classi}
\end{table}

\begin{table}[H]
\centering
\begin{tabularx}{\textwidth}{l >{\raggedright\arraybackslash}p{0.28\textwidth} >{\raggedright\arraybackslash}X}
\toprule
Metodo & Parametro & Classi di equivalenza \\
\midrule
\texttt{onFailure} & \texttt{caught} & \{istanza qualunque di \texttt{Throwable}\} \\
\midrule
\texttt{onSuccess} & \texttt{uri}, formato & \{null\} -- \{vuota\} -- \{non valida\} --
\{valida\} \\
\texttt{onSuccess} & \texttt{uri}, raggiungibilità* & \{corretta\} -- \{non corretta\} \\
\bottomrule
\end{tabularx}
\caption{Classi di equivalenza per \texttt{GetServerURI} (\texttt{AsyncCallback<String>}).
* richiede formato di \texttt{uri}=\{valida\}.}
\label{tab:logindialog-getserveruri-classi}
\end{table}

\begin{table}[H]
\centering
\begin{tabularx}{\textwidth}{l >{\raggedright\arraybackslash}p{0.28\textwidth} >{\raggedright\arraybackslash}X}
\toprule
Metodo & Parametro & Classi di equivalenza \\
\midrule
\texttt{onFailure} & \texttt{caught} & \{istanza qualunque di \texttt{Throwable}\} \\
\midrule
\texttt{onSuccess} & \texttt{stocks} & \{null\} -- \{vuota\} -- \{con elementi\} \\
\bottomrule
\end{tabularx}
\caption{Classi di equivalenza per \texttt{InitializeStocks}
(\texttt{AsyncCallback<List<Stock>>}).}
\label{tab:logindialog-initializestocks-classi}
\end{table}

% Tabelle esaustive combinazioni+oracolo, tutte in lavorazione come CSV, da inserire
% qui solo dopo revisione:
% find       -> report/data/find_combinazioni.csv        (84 righe)
% lock       -> report/data/lock_combinazioni.csv         (96 righe)
% detachAll  -> report/data/detachAll_combinazioni.csv    (15 righe)
% attachAll  -> report/data/attachAll_combinazioni.csv    (42 righe)
% newInstance-> report/data/newInstance_combinazioni.csv   (7 righe)
% isDetached -> report/data/isDetached_combinazioni.csv    (10 righe)
% persist    -> report/data/persist_combinazioni.csv      (30 righe)

\paragraph*{Prompt Z1.}
\label{prompt:z1}
\begin{lstlisting}
As a professional software tester who writes Java test methods, generate a complete JUnit 4 test
file (BrokerImplTest.java) that covers 10 methods of your choice in the following class named
BrokerImpl, using at most 20 @Test-annotated methods in total. Your output file must start with
###Test START## and finish with ###Test END##. The source code is attached.
\end{lstlisting}

\paragraph*{Risposta Z1.}
\label{risposta:z1}
\begin{lstlisting}
This is a standalone test file, so I'll present it directly rather than running it
through the file-creation pipeline (it's a code deliverable meant to sit alongside
BrokerImpl.java in the OpenJPA source tree, not something to execute here).


Note on scope: BrokerImpl is a very large, heavily-wired OpenJPA
kernel class -- most of its methods require a fully initialized
AbstractBrokerFactory/OpenJPAConfiguration/DelegatingStoreManager (via initialize(...))
to avoid NPEs. Rather than mock that entire stack,
I selected 10 methods whose default code paths are safe to exercise on a bare new BrokerImpl():

setAuthentication (+ getConnectionUserName/getConnectionPassword)
getMultithreaded/setMultithreaded
setAutoDetach(int)/getAutoDetach (including the DETACH_NONE exclusivity check)
setDetachState/getDetachState
setDetachedNew/isDetachedNew
setRetainState/getRetainState
putUserObject/getUserObject
isActive
getRollbackOnly (+ setRollbackOnly exception path)
assertOpen (open and closed-via-reflection paths)

18 @Test methods total. The _closed field is flipped via reflection since
there's no public setter, in order to exercise assertOpen()'s exception branch.


###Test START##
package org.apache.openjpa.kernel;

import static org.junit.Assert.assertEquals;
import static org.junit.Assert.assertFalse;
import static org.junit.Assert.assertNull;
import static org.junit.Assert.assertTrue;
import static org.junit.Assert.fail;

import java.lang.reflect.Field;

import org.apache.openjpa.util.InvalidStateException;
import org.apache.openjpa.util.NoTransactionException;
import org.apache.openjpa.util.UserException;
import org.junit.Before;
import org.junit.Test;

/**
 * Unit tests for {@link BrokerImpl}.
 *
 * These tests exercise BrokerImpl methods that can be safely invoked on a
 * freshly constructed, un-initialized broker (i.e. without calling
 * {@code initialize(...)}, which requires a fully wired
 * AbstractBrokerFactory / OpenJPAConfiguration / DelegatingStoreManager).
 * Methods were selected so that their default code paths do not touch any
 * of the transient collaborator fields (_conf, _store, _runtime, _cache,
 * etc.) that remain null before initialization.
 */
public class BrokerImplTest {

    private BrokerImpl broker;

    @Before
    public void setUp() {
        broker = new BrokerImpl();
    }

    /**
     * Helper to flip the private "_closed" field via reflection, used to
     * exercise the exceptional branch of assertOpen().
     */
    private void setClosed(boolean closed) throws Exception {
        Field f = BrokerImpl.class.getDeclaredField("_closed");
        f.setAccessible(true);
        f.set(broker, closed);
    }

    // ------------------------------------------------------------------
    // setAuthentication(String, String) / getConnectionUserName() /
    // getConnectionPassword()
    // ------------------------------------------------------------------

    @Test
    public void testSetAuthenticationStoresUserAndPassword() {
        broker.setAuthentication("scott", "tiger");
        assertEquals("scott", broker.getConnectionUserName());
        assertEquals("tiger", broker.getConnectionPassword());
    }

    @Test
    public void testSetAuthenticationWithNullsClearsCredentials() {
        broker.setAuthentication("scott", "tiger");
        broker.setAuthentication(null, null);
        assertNull(broker.getConnectionUserName());
        assertNull(broker.getConnectionPassword());
    }

    // ------------------------------------------------------------------
    // getMultithreaded() / setMultithreaded(boolean)
    // ------------------------------------------------------------------

    @Test
    public void testGetMultithreadedDefaultsToFalse() {
        assertFalse(broker.getMultithreaded());
    }

    @Test
    public void testSetMultithreadedTrueUpdatesFlag() {
        broker.setMultithreaded(true);
        assertTrue(broker.getMultithreaded());
    }

    @Test
    public void testSetMultithreadedFalseAfterTrueUpdatesFlag() {
        broker.setMultithreaded(true);
        broker.setMultithreaded(false);
        assertFalse(broker.getMultithreaded());
    }

    // ------------------------------------------------------------------
    // setAutoDetach(int) / getAutoDetach()
    // ------------------------------------------------------------------

    @Test
    public void testSetAutoDetachValidFlagIsStored() {
        broker.setAutoDetach(AutoDetach.DETACH_CLOSE);
        assertEquals(AutoDetach.DETACH_CLOSE, broker.getAutoDetach());
    }

    @Test
    public void testSetAutoDetachDetachNoneCombinedWithOtherFlagThrows() {
        try {
            broker.setAutoDetach(AutoDetach.DETACH_NONE | AutoDetach.DETACH_CLOSE);
            fail("Expected a UserException when DETACH_NONE is combined with another flag");
        } catch (UserException expected) {
            // expected: DETACH_NONE must be exclusive
        }
    }

    // ------------------------------------------------------------------
    // setDetachState(int) / getDetachState()
    // ------------------------------------------------------------------

    @Test
    public void testGetDetachStateDefaultValue() {
        // Default value assigned at field declaration (DETACH_LOADED).
        assertEquals(DetachState.DETACH_LOADED, broker.getDetachState());
    }

    @Test
    public void testSetDetachStateUpdatesValue() {
        broker.setDetachState(DetachState.DETACH_ALL);
        assertEquals(DetachState.DETACH_ALL, broker.getDetachState());
    }

    // ------------------------------------------------------------------
    // setDetachedNew(boolean) / isDetachedNew()
    // ------------------------------------------------------------------

    @Test
    public void testIsDetachedNewDefaultsToTrue() {
        assertTrue(broker.isDetachedNew());
    }

    @Test
    public void testSetDetachedNewFalseUpdatesValue() {
        broker.setDetachedNew(false);
        assertFalse(broker.isDetachedNew());
    }

    // ------------------------------------------------------------------
    // setRetainState(boolean) / getRetainState()
    // ------------------------------------------------------------------

    @Test
    public void testGetRetainStateDefaultsToFalse() {
        assertFalse(broker.getRetainState());
    }

    @Test
    public void testSetRetainStateTrueUpdatesValue() {
        broker.setRetainState(true);
        assertTrue(broker.getRetainState());
    }

    // ------------------------------------------------------------------
    // putUserObject(Object, Object) / getUserObject(Object)
    // ------------------------------------------------------------------

    @Test
    public void testPutAndGetUserObjectRoundTrips() {
        Object key = "myKey";
        Object value = "myValue";
        Object previous = broker.putUserObject(key, value);
        assertNull(previous);
        assertEquals(value, broker.getUserObject(key));
    }

    @Test
    public void testGetUserObjectForUnknownKeyReturnsNull() {
        assertNull(broker.getUserObject("doesNotExist"));
    }

    @Test
    public void testPutUserObjectWithNullValueRemovesExistingEntry() {
        Object key = "removeMe";
        broker.putUserObject(key, "value");
        Object removed = broker.putUserObject(key, null);
        assertEquals("value", removed);
        assertNull(broker.getUserObject(key));
    }

    // ------------------------------------------------------------------
    // isActive()
    // ------------------------------------------------------------------

    @Test
    public void testIsActiveDefaultsToFalseOnNewBroker() {
        assertFalse(broker.isActive());
    }

    // ------------------------------------------------------------------
    // getRollbackOnly()
    // ------------------------------------------------------------------

    @Test
    public void testGetRollbackOnlyDefaultsToFalseWhenNoTransactionActive() {
        assertFalse(broker.getRollbackOnly());
    }

    @Test
    public void testSetRollbackOnlyWithoutActiveTransactionThrows() {
        try {
            broker.setRollbackOnly();
            fail("Expected NoTransactionException when there is no active transaction");
        } catch (NoTransactionException expected) {
            // expected: cannot mark rollback-only without an active transaction
        }
    }

    // ------------------------------------------------------------------
    // assertOpen()
    // ------------------------------------------------------------------

    @Test
    public void testAssertOpenDoesNotThrowWhenBrokerIsOpen() {
        // Should simply return without throwing on a fresh, un-closed broker.
        broker.assertOpen();
    }

    @Test
    public void testAssertOpenThrowsWhenBrokerIsClosed() throws Exception {
        setClosed(true);
        try {
            broker.assertOpen();
            fail("Expected InvalidStateException when broker has been closed");
        } catch (InvalidStateException expected) {
            // expected: operations are not allowed on a closed broker
        }
    }
}
###Test END##
\end{lstlisting}

\paragraph*{Prompt Z2.}
\label{prompt:z2}
\begin{lstlisting}
As a professional software tester who writes Java test methods, generate a complete JUnit 4 test
file (BrokerImplTest.java) that covers 10 methods of your choice in the following class named
BrokerImpl, using at most 20 @Test-annotated methods in total. Focus specifically on functional,
black-box test cases based on the documented and expected behavior of each method (including
edge cases and exception handling), not on internal implementation details. Your output file
must start with ###Test START## and finish with ###Test END##. The source code is attached.
\end{lstlisting}

\paragraph*{Risposta Z2.}
\label{risposta:z2}
\begin{lstlisting}
This test suite covers 10+ public methods of BrokerImpl
(setAuthentication, getConnectionUserName/getConnectionPassword, setConnectionFactoryName/getConnectionFactoryName,
setConnectionFactory2Name/getConnectionFactory2Name, getConnectionFactory, getConnectionFactory2,
setAutoDetach(int), setAutoDetach(int, boolean), putUserObject/getUserObject,
setSuppressBatchOLELogging/getSuppressBatchOLELogging, and clone()) with 20 @Test methods.
I deliberately chose methods whose documented behavior is observable on a freshly constructed
BrokerImpl without needing the full initialize(...) wiring (factory, config, store manager, etc.),
since that setup would require extensive internal mocking rather than true black-box testing.
Reflection is used in exactly one place (the clone() tests), solely to arrange the pre-condition
described in the javadoc ("clone fails once initialize() has run") -- not to assert on internals.

###Test START##
/*
 * Licensed to the Apache Software Foundation (ASF) under one
 * or more contributor license agreements.  See the NOTICE file
 * distributed with this work for additional information
 * regarding copyright ownership.  The ASF licenses this file
 * to you under the Apache License, Version 2.0 (the
 * "License"); you may not use this file except in compliance
 * with the License.  You may obtain a copy of the License at
 *
 * http://www.apache.org/licenses/LICENSE-2.0
 *
 * Unless required by applicable law or agreed to in writing,
 * software distributed under the License is distributed on an
 * "AS IS" BASIS, WITHOUT WARRANTIES OR CONDITIONS OF ANY
 * KIND, either express or implied.  See the License for the
 * specific language governing permissions and limitations
 * under the License.
 */
package org.apache.openjpa.kernel;

import static org.junit.Assert.assertEquals;
import static org.junit.Assert.assertFalse;
import static org.junit.Assert.assertNotNull;
import static org.junit.Assert.assertNotSame;
import static org.junit.Assert.assertNull;
import static org.junit.Assert.assertTrue;
import static org.junit.Assert.fail;

import java.lang.reflect.Field;

import org.apache.openjpa.util.UserException;
import org.junit.Before;
import org.junit.Test;

/**
 * Black-box functional tests for {@link BrokerImpl}.
 *
 * These tests exercise BrokerImpl methods that can be observed through their
 * public contract on a freshly constructed, un-initialized broker (i.e.
 * without invoking {@code initialize(...)}, which requires a fully wired
 * {@code AbstractBrokerFactory}/{@code OpenJPAConfiguration}). Only methods
 * whose documented behavior does not depend on that heavyweight setup are
 * covered here. Reflection is used strictly for test *arrangement* (forcing
 * a private flag that is otherwise only set deep inside initialize()), never
 * to inspect or assert on internal state.
 */
public class BrokerImplTest {

    private BrokerImpl broker;

    @Before
    public void setUp() {
        broker = new BrokerImpl();
    }

    // ------------------------------------------------------------------
    // setAuthentication(String, String)
    // ------------------------------------------------------------------

    @Test
    public void testSetAuthentication_StoresUserAndPassword() {
        broker.setAuthentication("alice", "secret");

        assertEquals("alice", broker.getConnectionUserName());
        assertEquals("secret", broker.getConnectionPassword());
    }

    @Test
    public void testSetAuthentication_AllowsNullValues() {
        broker.setAuthentication("alice", "secret");

        broker.setAuthentication(null, null);

        assertNull(broker.getConnectionUserName());
        assertNull(broker.getConnectionPassword());
    }

    // ------------------------------------------------------------------
    // setConnectionFactoryName(String) / getConnectionFactoryName()
    // ------------------------------------------------------------------

    @Test
    public void testSetConnectionFactoryName_TrimsWhitespace() {
        broker.setConnectionFactoryName("  jdbc/MyDS  ");

        assertEquals("jdbc/MyDS", broker.getConnectionFactoryName());
    }

    @Test
    public void testSetConnectionFactoryName_BlankBecomesNull() {
        broker.setConnectionFactoryName("    ");

        assertNull(broker.getConnectionFactoryName());
    }

    @Test
    public void testSetConnectionFactoryName_NullRemainsNull() {
        broker.setConnectionFactoryName(null);

        assertNull(broker.getConnectionFactoryName());
    }

    // ------------------------------------------------------------------
    // setConnectionFactory2Name(String) / getConnectionFactory2Name()
    // ------------------------------------------------------------------

    @Test
    public void testSetConnectionFactory2Name_TrimsWhitespace() {
        broker.setConnectionFactory2Name("  jdbc/MyDS2  ");

        assertEquals("jdbc/MyDS2", broker.getConnectionFactory2Name());
    }

    @Test
    public void testSetConnectionFactory2Name_BlankBecomesNull() {
        broker.setConnectionFactory2Name("   ");

        assertNull(broker.getConnectionFactory2Name());
    }

    // ------------------------------------------------------------------
    // getConnectionFactory() / getConnectionFactory2()
    // ------------------------------------------------------------------

    @Test
    public void testGetConnectionFactory_ReturnsNullWhenNameBlank() {
        // A brand new broker has no connection factory name configured.
        assertNull(broker.getConnectionFactoryName());

        assertNull(broker.getConnectionFactory());
    }

    @Test
    public void testGetConnectionFactory2_ReturnsNullWhenNameBlank() {
        assertNull(broker.getConnectionFactory2Name());

        assertNull(broker.getConnectionFactory2());
    }

    // ------------------------------------------------------------------
    // setAutoDetach(int) / getAutoDetach()
    // ------------------------------------------------------------------

    @Test
    public void testSetAutoDetach_SingleValidFlag() {
        broker.setAutoDetach(AutoDetach.DETACH_COMMIT);

        assertEquals(AutoDetach.DETACH_COMMIT, broker.getAutoDetach());
    }

    @Test
    public void testSetAutoDetach_DetachNoneAlone_Succeeds() {
        broker.setAutoDetach(AutoDetach.DETACH_NONE);

        assertEquals(AutoDetach.DETACH_NONE, broker.getAutoDetach());
    }

    @Test
    public void testSetAutoDetach_DetachNoneCombinedWithOtherFlag_ThrowsUserException() {
        try {
            broker.setAutoDetach(AutoDetach.DETACH_NONE | AutoDetach.DETACH_COMMIT);
            fail("Expected a UserException because DETACH_NONE must be exclusive");
        } catch (UserException expected) {
            // documented behavior: DETACH_NONE cannot be combined with other flags
        }
    }

    // ------------------------------------------------------------------
    // setAutoDetach(int, boolean)
    // ------------------------------------------------------------------

    @Test
    public void testSetAutoDetachOverload_AddsFlagWhenOnTrue() {
        broker.setAutoDetach(AutoDetach.DETACH_COMMIT, true);

        assertTrue((broker.getAutoDetach() & AutoDetach.DETACH_COMMIT) != 0);
    }

    @Test
    public void testSetAutoDetachOverload_RemovesFlagWhenOnFalse() {
        broker.setAutoDetach(AutoDetach.DETACH_COMMIT, true);
        assertTrue((broker.getAutoDetach() & AutoDetach.DETACH_COMMIT) != 0);

        broker.setAutoDetach(AutoDetach.DETACH_COMMIT, false);

        assertFalse((broker.getAutoDetach() & AutoDetach.DETACH_COMMIT) != 0);
    }

    // ------------------------------------------------------------------
    // putUserObject(Object, Object) / getUserObject(Object)
    // ------------------------------------------------------------------

    @Test
    public void testPutUserObject_StoreRetrieveAndOverwrite() {
        Object previous = broker.putUserObject("key1", "value1");
        assertNull(previous);
        assertEquals("value1", broker.getUserObject("key1"));

        Object overwritten = broker.putUserObject("key1", "value2");
        assertEquals("value1", overwritten);
        assertEquals("value2", broker.getUserObject("key1"));
    }

    @Test
    public void testGetUserObject_ReturnsNullForMissingKey() {
        assertNull(broker.getUserObject("does-not-exist"));
    }

    @Test
    public void testPutUserObject_NullValueRemovesExistingEntry() {
        broker.putUserObject("key1", "value1");
        assertEquals("value1", broker.getUserObject("key1"));

        Object removed = broker.putUserObject("key1", null);

        assertEquals("value1", removed);
        assertNull(broker.getUserObject("key1"));
    }

    // ------------------------------------------------------------------
    // setSuppressBatchOLELogging(boolean) / getSuppressBatchOLELogging()
    // ------------------------------------------------------------------

    @Test
    public void testSuppressBatchOLELogging_SetAndGet() {
        broker.setSuppressBatchOLELogging(true);
        assertTrue(broker.getSuppressBatchOLELogging());

        broker.setSuppressBatchOLELogging(false);
        assertFalse(broker.getSuppressBatchOLELogging());
    }

    // ------------------------------------------------------------------
    // clone()
    // ------------------------------------------------------------------

    @Test
    public void testClone_SucceedsWhenNotInitialized() throws Exception {
        Object copy = broker.clone();

        assertNotNull(copy);
        assertTrue(copy instanceof BrokerImpl);
        assertNotSame(broker, copy);
    }

    @Test
    public void testClone_ThrowsWhenAlreadyInitialized() throws Exception {
        // Force the internal flag that initialize(...) sets, without running
        // the full (heavyweight) initialization sequence, so we can observe
        // the documented post-initialize cloning restriction.
        setInitializeWasInvoked(broker, true);

        try {
            broker.clone();
            fail("Expected CloneNotSupportedException once the broker has been initialized");
        } catch (CloneNotSupportedException expected) {
            // documented behavior: an initialized broker cannot be cloned
        }
    }

    private static void setInitializeWasInvoked(BrokerImpl broker, boolean value) throws Exception {
        Field field = BrokerImpl.class.getDeclaredField("_initializeWasInvoked");
        field.setAccessible(true);
        field.setBoolean(broker, value);
    }
}
###Test END##
\end{lstlisting}

\paragraph*{Esempio few-shot 1.}
\label{prompt:fewshot1}
\begin{lstlisting}
public class Calculator {
    public int add(int a, int b) { return a + b; }
}

public class CalculatorTest {
    @Test
    public void testAdd() {
        Calculator calc = new Calculator();
        assertEquals(5, calc.add(2, 3));
    }
}
\end{lstlisting}

\paragraph*{Esempio few-shot 2.}
\label{prompt:fewshot2}
\begin{lstlisting}
public class Calculator {
    public int divide(int a, int b) {
        if (b == 0) {
            throw new ArithmeticException("division by zero");
        }
        return a / b;
    }
}

public class CalculatorTest {
    @Test(expected = ArithmeticException.class)
    public void testDivideByZero() {
        Calculator calc = new Calculator();
        calc.divide(10, 0);
    }
}
\end{lstlisting}

\paragraph*{Prompt F1.}
\label{prompt:f1}
\begin{lstlisting}
As a professional software tester who writes Java test methods, consider the attached example.
Now, generate JUnit 4 test cases that cover 10 methods of your choice in the following class
named BrokerImpl, using at most 20 @Test-annotated methods in total. Your output file must start
with ###Test START## and finish with ###Test END##. The source code is attached.
\end{lstlisting}

\paragraph*{Risposta F1.}
\label{risposta:f1}
\begin{lstlisting}
I reviewed the attached BrokerImpl class and wrote JUnit 4 tests for 10 methods that operate
on internal broker state without requiring the heavyweight initialize(...) wiring
(since lock()/unlock() are no-ops until multithreading is enabled, and assertOpen()
only checks the default-false _closed flag).
Methods covered: setAuthentication, setMultithreaded/getMultithreaded, setIgnoreChanges,
setRetainState, setAutoClear, setAutoDetach, setDetachedNew, putUserObject/getUserObject,
clone(), setConnectionFactoryName/getConnectionFactory.

18 @Test methods across 10 target methods, self-contained via new BrokerImpl()
(no mocking framework needed), with the required
###Test START## / ###Test END## markers wrapping the file content.

###Test START##
package org.apache.openjpa.kernel;

import static org.junit.Assert.assertEquals;
import static org.junit.Assert.assertFalse;
import static org.junit.Assert.assertNotNull;
import static org.junit.Assert.assertNotSame;
import static org.junit.Assert.assertNull;
import static org.junit.Assert.assertTrue;

import org.junit.Before;
import org.junit.Test;

/**
 * JUnit 4 tests for {@link BrokerImpl}.
 *
 * These tests exercise BrokerImpl methods that operate purely on the
 * broker's internal state and do not require the heavyweight
 * initialize(...) call (which needs a fully wired AbstractBrokerFactory,
 * DelegatingStoreManager and OpenJPAConfiguration). lock()/unlock() are
 * no-ops until setMultithreaded(true) is invoked, and assertOpen() only
 * checks the (initially false) _closed flag, so a plain "new BrokerImpl()"
 * is sufficient to exercise this behavior.
 */
public class BrokerImplTest {

    private BrokerImpl broker;

    @Before
    public void setUp() {
        broker = new BrokerImpl();
    }

    // ------------------------------------------------------------------
    // setAuthentication / getConnectionUserName / getConnectionPassword
    // ------------------------------------------------------------------

    @Test
    public void testSetAuthenticationStoresUserAndPassword() {
        broker.setAuthentication("scott", "tiger");
        assertEquals("scott", broker.getConnectionUserName());
        assertEquals("tiger", broker.getConnectionPassword());
    }

    @Test
    public void testConnectionUserNameAndPasswordDefaultToNull() {
        assertNull(broker.getConnectionUserName());
        assertNull(broker.getConnectionPassword());
    }

    // ------------------------------------------------------------------
    // setMultithreaded / getMultithreaded
    // ------------------------------------------------------------------

    @Test
    public void testSetMultithreadedTrueIsReflectedByGetter() {
        assertFalse(broker.getMultithreaded());
        broker.setMultithreaded(true);
        assertTrue(broker.getMultithreaded());
    }

    @Test
    public void testSetMultithreadedFalseAfterTrueDoesNotThrow() {
        broker.setMultithreaded(true);
        broker.setMultithreaded(false);
        assertFalse(broker.getMultithreaded());
        // lock()/unlock() must remain safe no-ops once the internal lock is cleared
        broker.lock();
        broker.unlock();
    }

    // ------------------------------------------------------------------
    // setIgnoreChanges / getIgnoreChanges
    // ------------------------------------------------------------------

    @Test
    public void testSetIgnoreChangesUpdatesState() {
        assertFalse(broker.getIgnoreChanges());
        broker.setIgnoreChanges(true);
        assertTrue(broker.getIgnoreChanges());
        broker.setIgnoreChanges(false);
        assertFalse(broker.getIgnoreChanges());
    }

    // ------------------------------------------------------------------
    // setRetainState / getRetainState
    // ------------------------------------------------------------------

    @Test
    public void testSetRetainStateUpdatesState() {
        assertFalse(broker.getRetainState());
        broker.setRetainState(true);
        assertTrue(broker.getRetainState());
    }

    // ------------------------------------------------------------------
    // setAutoClear / getAutoClear
    // ------------------------------------------------------------------

    @Test
    public void testSetAutoClearUpdatesState() {
        assertEquals(Broker.CLEAR_DATASTORE, broker.getAutoClear());
        broker.setAutoClear(Broker.CLEAR_ALL);
        assertEquals(Broker.CLEAR_ALL, broker.getAutoClear());
    }

    // ------------------------------------------------------------------
    // setAutoDetach(int) / getAutoDetach
    // ------------------------------------------------------------------

    @Test
    public void testSetAutoDetachStoresFlags() {
        assertEquals(0, broker.getAutoDetach());
        broker.setAutoDetach(AutoDetach.DETACH_CLOSE);
        assertEquals(AutoDetach.DETACH_CLOSE, broker.getAutoDetach());
    }

    @Test
    public void testSetAutoDetachOverwritesPreviousValue() {
        broker.setAutoDetach(AutoDetach.DETACH_CLOSE);
        broker.setAutoDetach(AutoDetach.DETACH_COMMIT);
        assertEquals(AutoDetach.DETACH_COMMIT, broker.getAutoDetach());
    }

    // ------------------------------------------------------------------
    // setDetachedNew / isDetachedNew
    // ------------------------------------------------------------------

    @Test
    public void testSetDetachedNewDefaultsToTrue() {
        assertTrue(broker.isDetachedNew());
    }

    @Test
    public void testSetDetachedNewCanBeSetFalse() {
        broker.setDetachedNew(false);
        assertFalse(broker.isDetachedNew());
    }

    // ------------------------------------------------------------------
    // putUserObject / getUserObject
    // ------------------------------------------------------------------

    @Test
    public void testGetUserObjectReturnsNullWhenNotSet() {
        assertNull(broker.getUserObject("missing-key"));
    }

    @Test
    public void testPutAndGetUserObjectRoundTrips() {
        Object previous = broker.putUserObject("key1", "value1");
        assertNull(previous);
        assertEquals("value1", broker.getUserObject("key1"));
    }

    @Test
    public void testPutUserObjectReturnsPreviousValue() {
        broker.putUserObject("key1", "value1");
        Object previous = broker.putUserObject("key1", "value2");
        assertEquals("value1", previous);
        assertEquals("value2", broker.getUserObject("key1"));
    }

    @Test
    public void testPutUserObjectWithNullValueRemovesEntry() {
        broker.putUserObject("key1", "value1");
        Object removed = broker.putUserObject("key1", null);
        assertEquals("value1", removed);
        assertNull(broker.getUserObject("key1"));
    }

    // ------------------------------------------------------------------
    // clone()
    // ------------------------------------------------------------------

    @Test
    public void testCloneReturnsDistinctEqualStateInstance() throws CloneNotSupportedException {
        broker.setAuthentication("scott", "tiger");
        Object clone = broker.clone();
        assertNotNull(clone);
        assertNotSame(broker, clone);
        assertTrue(clone instanceof BrokerImpl);
        assertEquals("scott", ((BrokerImpl) clone).getConnectionUserName());
    }

    // ------------------------------------------------------------------
    // setConnectionFactoryName / getConnectionFactoryName / getConnectionFactory
    // ------------------------------------------------------------------

    @Test
    public void testConnectionFactoryNameDefaultsToEmptyString() {
        assertEquals("", broker.getConnectionFactoryName());
    }

    @Test
    public void testSetConnectionFactoryNameTrimsBlankToNull() {
        broker.setConnectionFactoryName("   ");
        assertNull(broker.getConnectionFactoryName());
    }

    @Test
    public void testSetConnectionFactoryNameStoresTrimmedValue() {
        broker.setConnectionFactoryName("  myConnectionFactory  ");
        assertEquals("myConnectionFactory", broker.getConnectionFactoryName());
    }

    @Test
    public void testGetConnectionFactoryReturnsNullWhenNameBlank() {
        // default _connectionFactoryName is "" (blank), so no JNDI lookup should occur
        assertNull(broker.getConnectionFactory());
    }
}
###Test END##
\end{lstlisting}

\paragraph*{Prompt F2.}
\label{prompt:f2}
\begin{lstlisting}
As a professional software tester who writes Java test methods, consider the two attached
examples of well-written JUnit 4 test cases, one with a boundary value and one that expects an
exception. Now, generate JUnit 4 test cases that cover 10 methods of your choice in the
following class named BrokerImpl, using at most 20 @Test-annotated methods in total. Your output
file must start with ###Test START## and finish with ###Test END##. The source code is attached.
\end{lstlisting}

\paragraph*{Risposta F2.}
\label{risposta:f2}
\begin{lstlisting}
20 JUnit 4 tests covering 10 methods of BrokerImpl, all runnable against a
freshly-constructed instance (no mocking needed):

Simple getters/setters (ConnectionFactoryName, ConnectionFactory2Name,
Multithreaded, AutoClear, DetachState, AllowReferenceToSiblingContext,
PostLoadOnMerge) -- normal-value and boundary-value cases (blank-string
trimming, Integer.MIN_VALUE/MAX_VALUE, negative/zero, default-state checks).

getCachePreparedQuery/setCachePreparedQuery -- boundary case where false
short-circuits safely, plus an exception case where the default true forces
a null-dereference on the never-initialized _conf field.

getCacheFinderQuery/setCacheFinderQuery -- an exception test on the getter,
and a test that documents a real bug in the source: setCacheFinderQuery
mistakenly writes to _cachePreparedQuery instead of _cacheFinderQuery.

isCached -- boundary case (empty list returns true without touching the cache)
and exception case (non-empty list null-dereferences the uninitialized _cache).

Note: several tests expect NullPointerException because they exercise a BrokerImpl
that was never passed through initialize(...) -- that's intentional, since it's the
only way to reach that code without standing up the full OpenJPA configuration/store
stack, and it still verifies real, observable behavior.

###Test START##
package org.apache.openjpa.kernel;

import static org.junit.Assert.assertEquals;
import static org.junit.Assert.assertFalse;
import static org.junit.Assert.assertNull;
import static org.junit.Assert.assertTrue;

import java.util.ArrayList;
import java.util.Collections;
import java.util.List;

import org.junit.Before;
import org.junit.Test;

/**
 * JUnit 4 tests for {@link BrokerImpl}.
 *
 * These tests exercise BrokerImpl in its freshly-constructed state
 * (i.e. before {@code initialize(...)} has been called), which is
 * sufficient to cover the simple property accessors below as well as
 * a few methods whose behavior depends on collaborators (such as
 * {@code _conf} or {@code _cache}) that remain {@code null} until the
 * broker is initialized. Where a method depends on such an
 * uninitialized collaborator, the test documents the resulting
 * exception rather than mocking the collaborator.
 */
public class BrokerImplTest {

    private BrokerImpl broker;

    @Before
    public void setUp() {
        broker = new BrokerImpl();
    }

    // ------------------------------------------------------------------
    // getConnectionFactoryName / setConnectionFactoryName
    // ------------------------------------------------------------------

    @Test
    public void testSetConnectionFactoryName_NormalValueIsStored() {
        broker.setConnectionFactoryName("jdbc/MyFactory");
        assertEquals("jdbc/MyFactory", broker.getConnectionFactoryName());
    }

    @Test
    public void testSetConnectionFactoryName_BlankValueTrimmedToNull() {
        // boundary: a whitespace-only string is trimmed to null
        broker.setConnectionFactoryName("   ");
        assertNull(broker.getConnectionFactoryName());
    }

    // ------------------------------------------------------------------
    // getConnectionFactory2Name / setConnectionFactory2Name
    // ------------------------------------------------------------------

    @Test
    public void testSetConnectionFactory2Name_NormalValueIsStored() {
        broker.setConnectionFactory2Name("jdbc/MyFactory2");
        assertEquals("jdbc/MyFactory2", broker.getConnectionFactory2Name());
    }

    @Test
    public void testSetConnectionFactory2Name_BlankValueTrimmedToNull() {
        // boundary: empty string is trimmed to null
        broker.setConnectionFactory2Name("");
        assertNull(broker.getConnectionFactory2Name());
    }

    // ------------------------------------------------------------------
    // getMultithreaded / setMultithreaded
    // ------------------------------------------------------------------

    @Test
    public void testGetMultithreaded_DefaultIsFalse() {
        // boundary: default (un-set) state
        assertFalse(broker.getMultithreaded());
    }

    @Test
    public void testSetMultithreaded_TrueIsReflectedByGetter() {
        broker.setMultithreaded(true);
        assertTrue(broker.getMultithreaded());
    }

    // ------------------------------------------------------------------
    // getAutoClear / setAutoClear
    // ------------------------------------------------------------------

    @Test
    public void testSetAutoClear_BoundaryMinValue() {
        broker.setAutoClear(Integer.MIN_VALUE);
        assertEquals(Integer.MIN_VALUE, broker.getAutoClear());
    }

    @Test
    public void testSetAutoClear_BoundaryMaxValue() {
        broker.setAutoClear(Integer.MAX_VALUE);
        assertEquals(Integer.MAX_VALUE, broker.getAutoClear());
    }

    // ------------------------------------------------------------------
    // getDetachState / setDetachState
    // ------------------------------------------------------------------

    @Test
    public void testSetDetachState_BoundaryZero() {
        broker.setDetachState(0);
        assertEquals(0, broker.getDetachState());
    }

    @Test
    public void testSetDetachState_NegativeValue() {
        broker.setDetachState(-1);
        assertEquals(-1, broker.getDetachState());
    }

    // ------------------------------------------------------------------
    // getAllowReferenceToSiblingContext / setAllowReferenceToSiblingContext
    // ------------------------------------------------------------------

    @Test
    public void testGetAllowReferenceToSiblingContext_DefaultIsFalse() {
        // boundary: default (un-set) state
        assertFalse(broker.getAllowReferenceToSiblingContext());
    }

    @Test
    public void testSetAllowReferenceToSiblingContext_TrueIsReflectedByGetter() {
        broker.setAllowReferenceToSiblingContext(true);
        assertTrue(broker.getAllowReferenceToSiblingContext());
    }

    // ------------------------------------------------------------------
    // getPostLoadOnMerge / setPostLoadOnMerge
    // ------------------------------------------------------------------

    @Test
    public void testGetPostLoadOnMerge_DefaultIsFalse() {
        // boundary: default (un-set) state
        assertFalse(broker.getPostLoadOnMerge());
    }

    @Test
    public void testSetPostLoadOnMerge_TrueIsReflectedByGetter() {
        broker.setPostLoadOnMerge(true);
        assertTrue(broker.getPostLoadOnMerge());
    }

    // ------------------------------------------------------------------
    // getCachePreparedQuery / setCachePreparedQuery
    // ------------------------------------------------------------------

    @Test
    public void testSetCachePreparedQuery_FalseShortCircuitsWithoutTouchingConfiguration() {
        // boundary: once the flag is false, getCachePreparedQuery() short-circuits
        // and never dereferences the (uninitialized, null) configuration object.
        broker.setCachePreparedQuery(false);
        assertFalse(broker.getCachePreparedQuery());
    }

    @Test(expected = NullPointerException.class)
    public void testGetCachePreparedQuery_DefaultTrueThrowsWhenConfigurationUnset() {
        // The flag defaults to true, so getCachePreparedQuery() must evaluate
        // _conf.getQuerySQLCacheInstance(); since the broker was never
        // initialize()'d, _conf is null and a NullPointerException results.
        broker.getCachePreparedQuery();
    }

    // ------------------------------------------------------------------
    // getCacheFinderQuery / setCacheFinderQuery
    // ------------------------------------------------------------------

    @Test(expected = NullPointerException.class)
    public void testGetCacheFinderQuery_DefaultTrueThrowsWhenConfigurationUnset() {
        // The flag defaults to true, so getCacheFinderQuery() must evaluate
        // _conf.getFinderCacheInstance(); since the broker was never
        // initialize()'d, _conf is null and a NullPointerException results.
        broker.getCacheFinderQuery();
    }

    @Test
    public void testSetCacheFinderQuery_ActuallyUpdatesCachePreparedQueryFlag() {
        // Documents the observed behavior of setCacheFinderQuery(boolean):
        // it assigns to _cachePreparedQuery rather than _cacheFinderQuery,
        // so the change becomes visible through getCachePreparedQuery(),
        // which can then be read back without triggering a NullPointerException.
        broker.setCacheFinderQuery(false);
        assertFalse(broker.getCachePreparedQuery());
    }

    // ------------------------------------------------------------------
    // isCached
    // ------------------------------------------------------------------

    @Test
    public void testIsCached_EmptyListReturnsTrueWithoutTouchingCache() {
        // boundary: an empty oid list means every element (vacuously) is
        // already loaded, so the method returns true without dereferencing
        // the (uninitialized, null) managed cache.
        List<Object> oids = new ArrayList<>();
        assertTrue(broker.isCached(oids));
    }

    @Test(expected = NullPointerException.class)
    public void testIsCached_NonEmptyListThrowsWhenCacheUnset() {
        // A non-empty oid list forces a lookup against the (uninitialized,
        // null) managed cache, resulting in a NullPointerException.
        List<Object> oids = Collections.singletonList((Object) "some-oid");
        broker.isCached(oids);
    }
}
###Test END##
\end{lstlisting}

\paragraph*{Prompt C1.}
\label{prompt:c1}
\begin{lstlisting}
You are a professional software tester. Before writing any code, think step by step:
1. Choose 10 methods of the following class named BrokerImpl that you consider worth testing,
   and briefly explain why for each one.
2. For each chosen method, write down what you expect it to do and one edge case or exception
   that could occur.
3. Only after completing steps 1 and 2, write a complete JUnit 4 test file (BrokerImplTest.java)
   using at most 20 @Test-annotated methods in total, one or more per method chosen in step 1,
   based on what you wrote down in step 2.

Show your reasoning for steps 1 and 2 before the code. Your output file must start with ###Test
START## and finish with ###Test END##. The Java class is attached.
\end{lstlisting}

\paragraph*{Risposta C1.}
\label{risposta:c1}
\begin{lstlisting}
Step 1 & 2: Method Selection and Analysis
I reviewed the ~5,500-line BrokerImpl class. Most public methods depend
on a fully-initialized broker (_conf, _store, _repo, active transactions,
etc.), which would require heavy mocking of internal OpenJPA infrastructure.
I deliberately chose methods that are testable against a plain new BrokerImpl()
(no initialize() call), since assertOpen() only checks the _closed flag (false by default)
rather than requiring a real store/config. This keeps the tests self-contained while still
exercising real logic, validation, and -- notably -- two apparent bugs in the class.

setAuthentication(String,String) + getConnectionUserName()/getConnectionPassword()

Simplest state-setting method; first thing called after construction per its own Javadoc.


setConnectionFactoryName(String) + getConnectionFactoryName()

Uses StringUtil.trimToNull, a non-trivial transformation worth verifying.


setConnectionFactory2Name(String) + getConnectionFactory2Name()

Mirrors #2 for the "NonJTA" factory; same trimming logic, separate field -- worth independent
verification since it's easy to copy-paste a bug between the two.


getConnectionFactory()

Guards against blank name before attempting a JNDI lookup via Configurations.lookup.


setAutoDetach(int) + getAutoDetach()

Contains real validation logic (assertAutoDetachValue) -- a private helper only reachable through this public method.


setMultithreaded(boolean) + getMultithreaded()

Has a side effect beyond the field assignment: it allocates/releases an internal ReentrantLock.


setRetainState(boolean) + getRetainState()

Representative of the many simple boolean flag setters; verifies default value and pass-through


setSuppressBatchOLELogging(boolean) + getSuppressBatchOLELogging()

Javadoc explicitly claims "Defaults to true" but the field is declared = false. This is a documentation/implementation mismatch worth catching.


getCachePreparedQuery() + setCachePreparedQuery(boolean)

_cachePreparedQuery defaults to true, and the getter unconditionally dereferences _conf.getQuerySQLCacheInstance() when true. On an uninitialized broker _conf is null.


setCacheFinderQuery(boolean) + getCacheFinderQuery()

Inspecting the source reveals setCacheFinderQuery erroneously assigns to _cachePreparedQuery instead of _cacheFinderQuery -- a real bug.


A couple of things worth flagging: this test suite deliberately avoids
calling initialize(...) on the broker, since doing so realistically would
require mocking StoreManager, OpenJPAConfiguration, MetaDataRepository,
etc. -- a much larger effort. It also surfaces two real defects in the
class as it stands: getSuppressBatchOLELogging()'s actual default (false)
contradicts its own Javadoc ("Defaults to true"), and setCacheFinderQuery(boolean)
mistakenly writes to _cachePreparedQuery instead of _cacheFinderQuery.
Both are captured as passing tests that document current (buggy) behavior
rather than the intended behavior -- worth raising with whoever owns this code.


###Test START##
package org.apache.openjpa.kernel;

import static org.junit.Assert.assertEquals;
import static org.junit.Assert.assertFalse;
import static org.junit.Assert.assertNull;
import static org.junit.Assert.assertTrue;
import static org.junit.Assert.fail;

import org.apache.openjpa.util.UserException;
import org.junit.Before;
import org.junit.Test;

/**
 * Unit tests for {@link BrokerImpl}.
 *
 * These tests exercise BrokerImpl instances that have NOT been through
 * initialize(...) (i.e. no real StoreManager / OpenJPAConfiguration is
 * wired up). This is intentional: the methods under test only touch
 * fields that are safe to access pre-initialization (assertOpen() only
 * checks the _closed flag, which is false on a fresh instance). Methods
 * that require a fully initialized broker (persist, find, commit, etc.)
 * are out of scope for these tests.
 */
public class BrokerImplTest {

    private BrokerImpl broker;

    @Before
    public void setUp() {
        broker = new BrokerImpl();
    }

    // ------------------------------------------------------------------
    // 1. setAuthentication / getConnectionUserName / getConnectionPassword
    // ------------------------------------------------------------------

    @Test
    public void testSetAuthentication_storesUserAndPassword() {
        broker.setAuthentication("scott", "tiger");
        assertEquals("scott", broker.getConnectionUserName());
        assertEquals("tiger", broker.getConnectionPassword());
    }

    @Test
    public void testSetAuthentication_nullValuesDoNotThrow() {
        broker.setAuthentication(null, null);
        assertNull(broker.getConnectionUserName());
        assertNull(broker.getConnectionPassword());
    }

    // ------------------------------------------------------------------
    // 2. setConnectionFactoryName / getConnectionFactoryName
    // ------------------------------------------------------------------

    @Test
    public void testSetConnectionFactoryName_storesNonBlankValue() {
        broker.setConnectionFactoryName("jdbc/MyDS");
        assertEquals("jdbc/MyDS", broker.getConnectionFactoryName());
    }

    @Test
    public void testSetConnectionFactoryName_blankStringBecomesNull() {
        broker.setConnectionFactoryName("   ");
        assertNull(broker.getConnectionFactoryName());
    }

    // ------------------------------------------------------------------
    // 3. setConnectionFactory2Name / getConnectionFactory2Name
    // ------------------------------------------------------------------

    @Test
    public void testSetConnectionFactory2Name_storesNonBlankValue() {
        broker.setConnectionFactory2Name("jdbc/MyDS2");
        assertEquals("jdbc/MyDS2", broker.getConnectionFactory2Name());
    }

    @Test
    public void testSetConnectionFactory2Name_blankStringBecomesNull() {
        broker.setConnectionFactory2Name("");
        assertNull(broker.getConnectionFactory2Name());
    }

    // ------------------------------------------------------------------
    // 4. getConnectionFactory
    // ------------------------------------------------------------------

    @Test
    public void testGetConnectionFactory_returnsNullWhenNameBlank() {
        // Default _connectionFactoryName is "" -> isNotBlank is false,
        // so no JNDI lookup should be attempted and null is returned.
        assertNull(broker.getConnectionFactory());
    }

    // ------------------------------------------------------------------
    // 5. setAutoDetach / getAutoDetach
    // ------------------------------------------------------------------

    @Test
    public void testSetAutoDetach_validSingleFlagIsStored() {
        broker.setAutoDetach(AutoDetach.DETACH_CLOSE);
        assertEquals(AutoDetach.DETACH_CLOSE, broker.getAutoDetach());
    }

    @Test(expected = UserException.class)
    public void testSetAutoDetach_detachNoneCombinedWithOtherFlagThrows() {
        // DETACH_NONE must be exclusive; combining it with another flag
        // should trigger assertAutoDetachValue's validation failure.
        broker.setAutoDetach(AutoDetach.DETACH_NONE | AutoDetach.DETACH_CLOSE);
    }

    // ------------------------------------------------------------------
    // 6. setMultithreaded / getMultithreaded
    // ------------------------------------------------------------------

    @Test
    public void testSetMultithreaded_togglesFlagCorrectly() {
        assertFalse(broker.getMultithreaded());
        broker.setMultithreaded(true);
        assertTrue(broker.getMultithreaded());
        broker.setMultithreaded(false);
        assertFalse(broker.getMultithreaded());
    }

    @Test
    public void testSetMultithreaded_repeatedTogglingDoesNotThrow() {
        try {
            broker.setMultithreaded(true);
            broker.setMultithreaded(false);
            broker.setMultithreaded(true);
            broker.setMultithreaded(true);
            broker.setMultithreaded(false);
        } catch (Exception e) {
            fail("Toggling multithreaded flag repeatedly should not throw: " + e);
        }
        assertFalse(broker.getMultithreaded());
    }

    // ------------------------------------------------------------------
    // 7. setRetainState / getRetainState
    // ------------------------------------------------------------------

    @Test
    public void testGetRetainState_defaultsToFalse() {
        assertFalse(broker.getRetainState());
    }

    @Test
    public void testSetRetainState_roundTrip() {
        broker.setRetainState(true);
        assertTrue(broker.getRetainState());
        broker.setRetainState(false);
        assertFalse(broker.getRetainState());
    }

    // ------------------------------------------------------------------
    // 8. setSuppressBatchOLELogging / getSuppressBatchOLELogging
    // ------------------------------------------------------------------

    @Test
    public void testGetSuppressBatchOLELogging_defaultIsFalse() {
        // NOTE: the method's Javadoc claims "Defaults to true", but the
        // backing field _suppressBatchOLELogging is declared as false.
        // This test documents the actual (buggy/undocumented) behavior.
        assertFalse(broker.getSuppressBatchOLELogging());
    }

    @Test
    public void testSetSuppressBatchOLELogging_roundTrip() {
        broker.setSuppressBatchOLELogging(true);
        assertTrue(broker.getSuppressBatchOLELogging());
        broker.setSuppressBatchOLELogging(false);
        assertFalse(broker.getSuppressBatchOLELogging());
    }

    // ------------------------------------------------------------------
    // 9. getCachePreparedQuery / setCachePreparedQuery
    // ------------------------------------------------------------------

    @Test(expected = NullPointerException.class)
    public void testGetCachePreparedQuery_defaultStateThrowsNPE() {
        // _cachePreparedQuery defaults to true, and the getter
        // unconditionally dereferences _conf.getQuerySQLCacheInstance()
        // when true. On an uninitialized broker _conf is null.
        broker.getCachePreparedQuery();
    }

    @Test
    public void testSetCachePreparedQuery_falseAvoidsConfDereference() {
        // When the flag is false, "&&" short-circuits before touching
        // the null _conf, so this should return false without throwing.
        broker.setCachePreparedQuery(false);
        assertFalse(broker.getCachePreparedQuery());
    }

    // ------------------------------------------------------------------
    // 10. setCacheFinderQuery / getCacheFinderQuery (documents a bug)
    // ------------------------------------------------------------------

    @Test(expected = NullPointerException.class)
    public void testSetCacheFinderQuery_bugDoesNotPreventNPE() {
        // BUG: setCacheFinderQuery(flag) mistakenly assigns to
        // _cachePreparedQuery instead of _cacheFinderQuery. As a result,
        // _cacheFinderQuery remains at its default value of true, and
        // getCacheFinderQuery() still dereferences the null _conf,
        // throwing NPE even after "disabling" the flag.
        broker.setCacheFinderQuery(false);
        broker.getCacheFinderQuery();
    }

    @Test
    public void testSetCacheFinderQuery_bugActuallyMutatesCachePreparedQuery() {
        // Demonstrates the side effect of the bug described above:
        // calling setCacheFinderQuery(false) actually flips
        // _cachePreparedQuery to false, which is observable via
        // getCachePreparedQuery() (short-circuits before the null _conf).
        broker.setCacheFinderQuery(false);
        assertFalse(broker.getCachePreparedQuery());
    }
}
###Test END##
\end{lstlisting}

\paragraph*{Prompt C2.}
\label{prompt:c2}
\begin{lstlisting}
You are a professional software tester. An example of a well-written JUnit 4 test case is
attached as a style reference. Before writing any code, think step by step:
1. Choose 10 methods of the following class named BrokerImpl that you consider worth testing,
   and briefly explain why for each one.
2. For each chosen method, write down what you expect it to do and one edge case or exception
   that could occur.
3. Only after completing steps 1 and 2, write a complete JUnit 4 test file (BrokerImplTest.java)
   using at most 20 @Test-annotated methods in total, one or more per method chosen in step 1,
   based on what you wrote down in step 2.

Show your reasoning for steps 1 and 2 before the code. Your output file must start with ###Test
START## and finish with ###Test END##. The Java class is attached.
\end{lstlisting}

\paragraph*{Risposta C2.}
\label{risposta:c2}
\begin{lstlisting}
Step 1: Method selection and rationale
BrokerImpl is a 5,500-line god-class with heavy runtime dependencies
(StoreManager, MetaDataRepository, configuration, transactions).
Most methods need a fully initialize()-d broker to run safely, which
would require mocking a huge dependency graph. I deliberately picked
10 methods that are reachable and meaningfully testable without calling initialize()
-- because they either operate on plain fields, have documented input-normalization
behavior, or fail in a specific, testable way (NPE) when the broker isn't wired up.
This also let me surface a real bug.

setAuthentication(String,String)
First method called after construction per javadoc; simple state mutation but a natural contract to lock down.

setConnectionFactoryName / getConnectionFactoryName
Javadoc explicitly promises "trimmed to null before being stored" -- a documented normalization contract worth verifying.

setConnectionFactory2Name / getConnectionFactory2Name
Same normalization contract, separate field/path (JTA vs non-JTA), easy to accidentally swap.

getConnectionFactory()
Branches on blank-name check before attempting a JNDI lookup; the "no name set" branch is deterministic and safe to test

getConnectionFactory2()
Mirrors #4 for the second connection factory; worth checking independently since it's a copy-pasted method (common bug source).

setAutoDetach(int) / getAutoDetach()
Contains real validation logic (assertAutoDetachValue) enforcing that DETACH_NONE is exclusive -- a genuine business rule with an exception path.

setDetachedNew(boolean) / isDetachedNew()
Simple flag, but default value is true (unusual default) -- worth pinning down explicitly.

putUserObject / getUserObject
Has real logic: lazy map init, and null value triggers removal rather than storing null -- an easy-to-miss edge case.

setCachePreparedQuery / getCachePreparedQuery
Setter is trivial, but getter dereferences _conf -- worth confirming it fails predictably (NPE) rather than silently misbehaving when broker isn't initialized.

setCacheFinderQuery / getCacheFinderQuery
Suspected bug: setCacheFinderQuery writes to _cachePreparedQuery instead of _cacheFinderQuery (copy-paste error, same pattern visible in setCacheFinderQuery's body). Worth a regression test either way.


Step 2: Expected behavior + edge case per method

setAuthentication -- stores user/pass verbatim, retrievable via getConnectionUserName/getConnectionPassword. Edge case: null values should be stored as-is (no NPE).
setConnectionFactoryName -- trims whitespace; edge case: empty/blank string should normalize to null, not "".
setConnectionFactory2Name -- same as above, independent field. Edge case: null input stays null.
getConnectionFactory -- returns null when _connectionFactoryName is blank (default state). Edge case: default constructed broker (name "") must not attempt a JNDI lookup and must not throw.
getConnectionFactory2 -- same edge case as #4 for the second factory.
setAutoDetach -- valid bitmask combinations (e.g. DETACH_CLOSE | DETACH_COMMIT) are stored as-is. Edge case/exception: combining DETACH_NONE with any other flag must throw UserException.
setDetachedNew -- default is true; setting false flips it. Edge case: verify the default explicitly, since it's counter-intuitive.
putUserObject/getUserObject -- stores and returns arbitrary key/value pairs. Edge case: passing null as the value removes the key (per source), and reading an unset key returns null without throwing.
setCachePreparedQuery/getCachePreparedQuery -- setter always succeeds and updates the flag. Edge case/exception: getter throws NullPointerException on an uninitialized broker because it dereferences _conf.
setCacheFinderQuery/getCacheFinderQuery -- expected: setter should flip _cacheFinderQuery. Edge case (bug): it actually flips _cachePreparedQuery instead, leaving _cacheFinderQuery untouched -- verified via reflection since the getter can't be called safely (same NPE as #9).


The reasoning for steps 1 and 2 is above;
the file below is the resulting test suite
(16 @Test methods across the 10 chosen methods, well under the 20-test cap).


One thing worth flagging: test #10 (testSetCacheFinderQueryActuallyMutatesCachePreparedQueryField) documents
what looks like a genuine copy-paste bug in setCacheFinderQuery -- it writes to _cachePreparedQuery
instead of _cacheFinderQuery. You may want to fix the source rather than keep this as a permanent
regression test.


###Test START##
package org.apache.openjpa.kernel;

import static org.junit.Assert.assertEquals;
import static org.junit.Assert.assertFalse;
import static org.junit.Assert.assertNull;
import static org.junit.Assert.assertTrue;
import static org.junit.Assert.fail;

import java.lang.reflect.Field;

import org.apache.openjpa.util.UserException;
import org.junit.Before;
import org.junit.Test;

/**
 * Unit tests for {@link BrokerImpl}.
 *
 * These tests deliberately avoid calling {@link BrokerImpl#initialize} because doing so
 * requires a fully wired {@code AbstractBrokerFactory}/{@code DelegatingStoreManager}/
 * {@code OpenJPAConfiguration} graph. Instead they exercise methods that are safe to call
 * on a bare, uninitialized broker (plain field mutators, documented normalization
 * behavior, and validation logic), plus the specific, deterministic failure mode
 * (NullPointerException) that occurs when a method needs configuration that was never set.
 */
public class BrokerImplTest {

    private BrokerImpl broker;

    @Before
    public void setUp() {
        broker = new BrokerImpl();
    }

    /**
     * Reads a private field via reflection since some fields under test have no
     * corresponding public accessor that can be safely invoked without initialize().
     */
    private Object getPrivateField(String name) throws Exception {
        Field f = BrokerImpl.class.getDeclaredField(name);
        f.setAccessible(true);
        return f.get(broker);
    }

    // ------------------------------------------------------------------
    // 1. setAuthentication
    // ------------------------------------------------------------------

    @Test
    public void testSetAuthenticationStoresUserAndPassword() {
        broker.setAuthentication("alice", "s3cr3t");
        assertEquals("alice", broker.getConnectionUserName());
        assertEquals("s3cr3t", broker.getConnectionPassword());
    }

    @Test
    public void testSetAuthenticationAllowsNullValues() {
        broker.setAuthentication(null, null);
        assertNull(broker.getConnectionUserName());
        assertNull(broker.getConnectionPassword());
    }

    // ------------------------------------------------------------------
    // 2. setConnectionFactoryName / getConnectionFactoryName
    // ------------------------------------------------------------------

    @Test
    public void testSetConnectionFactoryNameTrimsWhitespace() {
        broker.setConnectionFactoryName("  jndi/MyFactory  ");
        assertEquals("jndi/MyFactory", broker.getConnectionFactoryName());
    }

    @Test
    public void testSetConnectionFactoryNameBlankBecomesNull() {
        broker.setConnectionFactoryName("   ");
        assertNull(broker.getConnectionFactoryName());
    }

    // ------------------------------------------------------------------
    // 3. setConnectionFactory2Name / getConnectionFactory2Name
    // ------------------------------------------------------------------

    @Test
    public void testSetConnectionFactory2NameTrimsAndStores() {
        broker.setConnectionFactory2Name("  jndi/Factory2  ");
        assertEquals("jndi/Factory2", broker.getConnectionFactory2Name());
    }

    @Test
    public void testSetConnectionFactory2NameNullStaysNull() {
        broker.setConnectionFactory2Name(null);
        assertNull(broker.getConnectionFactory2Name());
    }

    // ------------------------------------------------------------------
    // 4. getConnectionFactory
    // ------------------------------------------------------------------

    @Test
    public void testGetConnectionFactoryReturnsNullWhenNameBlank() {
        // default state: _connectionFactoryName == ""
        assertNull(broker.getConnectionFactory());
    }

    // ------------------------------------------------------------------
    // 5. getConnectionFactory2
    // ------------------------------------------------------------------

    @Test
    public void testGetConnectionFactory2ReturnsNullWhenNameBlank() {
        // default state: _connectionFactory2Name == ""
        assertNull(broker.getConnectionFactory2());
    }

    // ------------------------------------------------------------------
    // 6. setAutoDetach / getAutoDetach
    // ------------------------------------------------------------------

    @Test
    public void testSetAutoDetachStoresValidCombination() {
        int flags = AutoDetach.DETACH_CLOSE | AutoDetach.DETACH_COMMIT;
        broker.setAutoDetach(flags);
        assertEquals(flags, broker.getAutoDetach());
    }

    @Test(expected = UserException.class)
    public void testSetAutoDetachNoneCombinedWithOtherFlagThrows() {
        broker.setAutoDetach(AutoDetach.DETACH_NONE | AutoDetach.DETACH_CLOSE);
    }

    // ------------------------------------------------------------------
    // 7. setDetachedNew / isDetachedNew
    // ------------------------------------------------------------------

    @Test
    public void testIsDetachedNewDefaultsToTrue() {
        assertTrue(broker.isDetachedNew());
    }

    @Test
    public void testSetDetachedNewFalseUpdatesFlag() {
        broker.setDetachedNew(false);
        assertFalse(broker.isDetachedNew());
    }

    // ------------------------------------------------------------------
    // 8. putUserObject / getUserObject
    // ------------------------------------------------------------------

    @Test
    public void testPutUserObjectStoresAndRetrievesValue() {
        Object previous = broker.putUserObject("key1", "value1");
        assertNull(previous);
        assertEquals("value1", broker.getUserObject("key1"));
    }

    @Test
    public void testPutUserObjectNullValueRemovesKey() {
        broker.putUserObject("key2", "value2");
        Object removed = broker.putUserObject("key2", null);
        assertEquals("value2", removed);
        assertNull(broker.getUserObject("key2"));
    }

    @Test
    public void testGetUserObjectReturnsNullWhenNeverSet() {
        assertNull(broker.getUserObject("neverSet"));
    }

    // ------------------------------------------------------------------
    // 9. setCachePreparedQuery / getCachePreparedQuery
    // ------------------------------------------------------------------

    @Test
    public void testSetCachePreparedQueryUpdatesField() throws Exception {
        broker.setCachePreparedQuery(false);
        assertEquals(Boolean.FALSE, getPrivateField("_cachePreparedQuery"));
    }

    @Test(expected = NullPointerException.class)
    public void testGetCachePreparedQueryThrowsWithoutConfiguration() {
        // _conf is null because initialize() was never called.
        broker.getCachePreparedQuery();
    }

    // ------------------------------------------------------------------
    // 10. setCacheFinderQuery / getCacheFinderQuery
    // ------------------------------------------------------------------

    @Test
    public void testSetCacheFinderQueryActuallyMutatesCachePreparedQueryField() throws Exception {
        // Documents existing (buggy) behavior: setCacheFinderQuery(false) writes to
        // _cachePreparedQuery instead of _cacheFinderQuery. If this test starts
        // failing because _cacheFinderQuery changed to false, the underlying bug
        // has likely been fixed and this test (and the production code) should be
        // revisited together.
        broker.setCacheFinderQuery(false);

        assertEquals("expected the (buggy) setter to leave _cacheFinderQuery untouched",
            Boolean.TRUE, getPrivateField("_cacheFinderQuery"));
        assertEquals("expected the (buggy) setter to have overwritten _cachePreparedQuery",
            Boolean.FALSE, getPrivateField("_cachePreparedQuery"));
    }

    @Test(expected = NullPointerException.class)
    public void testGetCacheFinderQueryThrowsWithoutConfiguration() {
        // _conf is null because initialize() was never called.
        broker.getCacheFinderQuery();
    }
}
###Test END##
\end{lstlisting}

\paragraph*{Comando di compilazione ed esecuzione, Test LLM \texttt{LoginDialog}.}
\label{cmd:test-logindialog}
\begin{lstlisting}
export JAVA8_HOME=/usr/lib/jvm/java-8-openjdk-amd64
mvn -f openjpa-examples/opentrader/pom.xml clean test
\end{lstlisting}
Stesso comando usato per tutte le sottosezioni precedenti su \texttt{LoginDialog}
(Sezioni~\ref{sec:cp},~\ref{sec:randoop},~\ref{sec:evosuite}); \texttt{-Dtest=NomeClasse} per una
singola variante durante un round di correzione.

\paragraph*{Prompt Z1 (LoginDialog).}
\label{prompt:ld-z1}
\begin{lstlisting}
As a professional software tester who writes Java test methods, generate a complete JUnit 4 test
file (LoginDialogTest.java) that covers 10 methods of your choice in the following class named
LoginDialog, using at most 20 @Test-annotated methods in total. Your output file must start with
###Test START## and finish with ###Test END##. The source code is attached.
\end{lstlisting}

\paragraph*{Risposta Z1 (LoginDialog).}
\label{risposta:ld-z1}
\begin{lstlisting}
LoginDialog's constructor requires an OpenTrader collaborator not included in the attached
source. Since only its usage (getService()) is visible from LoginDialog's own code, I mock it
with Mockito rather than guessing its full API. 10 @Test methods covering the constructor,
show/hide/center, setTitle/getTitle, setPopupPosition/getters, and three boolean properties
(modal, animation, glass).

###Test START##
package org.apache.openjpa.trader.client;

import static org.junit.Assert.assertEquals;
import static org.junit.Assert.assertFalse;
import static org.junit.Assert.assertTrue;
import static org.mockito.Mockito.mock;

import org.junit.Before;
import org.junit.Test;

public class LoginDialogTest {

    private OpenTrader session;
    private LoginDialog dialog;

    @Before
    public void setUp() {
        session = mock(OpenTrader.class);
        dialog = new LoginDialog(session);
    }

    @Test
    public void testConstructorBuildsAWidget() {
        assertTrue(dialog.getWidget() != null);
    }

    @Test
    public void testDialogIsNotShowingAfterConstruction() {
        assertFalse(dialog.isShowing());
    }

    @Test
    public void testShowMakesDialogVisible() {
        dialog.show();
        assertTrue(dialog.isShowing());
    }

    @Test
    public void testHideMakesDialogInvisible() {
        dialog.show();
        dialog.hide();
        assertFalse(dialog.isShowing());
    }

    @Test
    public void testSetTitleAndGetTitleRoundTrip() {
        dialog.setTitle("OpenTrader Login");
        assertEquals("OpenTrader Login", dialog.getTitle());
    }

    @Test
    public void testSetPopupPositionAndGetters() {
        dialog.setPopupPosition(10, 20);
        assertEquals(10, dialog.getPopupLeft());
        assertEquals(20, dialog.getPopupTop());
    }

    @Test
    public void testSetModalTrue() {
        dialog.setModal(true);
        assertTrue(dialog.isModal());
    }

    @Test
    public void testSetAnimationEnabledFalse() {
        dialog.setAnimationEnabled(false);
        assertFalse(dialog.isAnimationEnabled());
    }

    @Test
    public void testSetGlassEnabledTrue() {
        dialog.setGlassEnabled(true);
        assertTrue(dialog.isGlassEnabled());
    }

    @Test
    public void testCenterShowsDialog() {
        dialog.center();
        assertTrue(dialog.isShowing());
    }
}
###Test END##
\end{lstlisting}

\paragraph*{Prompt Z1, round 2 (LoginDialog).}
\label{prompt:ld-z1-r2}
\begin{lstlisting}
The test file you generated compiles, but all 10 tests fail with the same error. Representative
failure:

java.lang.ExceptionInInitializerError
  ...
Caused by: java.lang.UnsupportedOperationException: ERROR: GWT.create() is only usable in
client code! It cannot be called, for example, from server code. If you are running a unit
test, check that your test case extends GWTTestCase and that GWT.create() is not called from
within an initializer or constructor.
	at com.google.gwt.core.client.GWT.create(GWT.java:92)
	at com.google.gwt.user.client.ui.UIObject.<clinit>(UIObject.java:179)

Please fix the test file so it compiles and passes. Your output file must start with
###Test START## and finish with ###Test END##.
\end{lstlisting}

\paragraph*{Risposta Z1, round 2 (LoginDialog).}
\label{risposta:ld-z1-r2}
\begin{lstlisting}
The exception message says to check that the test case extends GWTTestCase and does not call
GWT.create() from an initializer or constructor. Rewriting as a GWTTestCase: JUnit 3 style
(GWTTestCase extends junit.framework.TestCase), so @Test/@Before become plain testXxx() methods
and a gwtSetUp() override. getModuleName() must return the GWT module owning this class; based
on the package (org.apache.openjpa.trader.client), assuming the module is
org.apache.openjpa.trader.client.LoginDialog.

###Test START##
package org.apache.openjpa.trader.client;

import com.google.gwt.junit.client.GWTTestCase;

public class LoginDialogTest extends GWTTestCase {

    @Override
    public String getModuleName() {
        return "org.apache.openjpa.trader.client.LoginDialog";
    }

    private OpenTrader session;
    private LoginDialog dialog;

    @Override
    protected void gwtSetUp() throws Exception {
        session = new OpenTrader();
        dialog = new LoginDialog(session);
    }

    public void testConstructorBuildsAWidget() {
        assertTrue(dialog.getWidget() != null);
    }

    public void testDialogIsNotShowingAfterConstruction() {
        assertFalse(dialog.isShowing());
    }

    public void testShowMakesDialogVisible() {
        dialog.show();
        assertTrue(dialog.isShowing());
    }

    public void testHideMakesDialogInvisible() {
        dialog.show();
        dialog.hide();
        assertFalse(dialog.isShowing());
    }

    public void testSetTitleAndGetTitleRoundTrip() {
        dialog.setTitle("OpenTrader Login");
        assertEquals("OpenTrader Login", dialog.getTitle());
    }

    public void testSetPopupPositionAndGetters() {
        dialog.show();
        dialog.setPopupPosition(10, 20);
        assertEquals(10, dialog.getPopupLeft());
        assertEquals(20, dialog.getPopupTop());
    }

    public void testSetModalTrue() {
        dialog.setModal(true);
        assertTrue(dialog.isModal());
    }

    public void testSetAnimationEnabledFalse() {
        dialog.setAnimationEnabled(false);
        assertFalse(dialog.isAnimationEnabled());
    }

    public void testSetGlassEnabledTrue() {
        dialog.setGlassEnabled(true);
        assertTrue(dialog.isGlassEnabled());
    }

    public void testCenterShowsDialog() {
        dialog.center();
        assertTrue(dialog.isShowing());
    }
}
###Test END##
\end{lstlisting}

\paragraph*{Prompt Z1, round 3 (LoginDialog).}
\label{prompt:ld-z1-r3}
\begin{lstlisting}
Still fails, now with a different error, the same on all 10 tests:

Loading inherited module 'org.apache.openjpa.trader.client.LoginDialog'
   [ERROR] Unable to find 'org/apache/openjpa/trader/client/LoginDialog.gwt.xml' on your
   classpath; could be a typo, or maybe you forgot to include a classpath entry for source?
com.google.gwt.core.ext.UnableToCompleteException: (see previous log entries)

Please fix the test file so it compiles and passes. Your output file must start with
###Test START## and finish with ###Test END##.
\end{lstlisting}

\paragraph*{Risposta Z1, round 3 (LoginDialog).}
\label{risposta:ld-z1-r3}
\begin{lstlisting}
The module name I guessed does not have a matching .gwt.xml file. A common GWT convention puts
the module definition one level above the client package. Trying
org.apache.openjpa.trader.LoginDialog instead.

###Test START##
package org.apache.openjpa.trader.client;

import com.google.gwt.junit.client.GWTTestCase;

/**
 * Revised after the reported failure: every test threw UnsupportedOperationException from
 * GWT.create(), and the exception message itself says to check that the test case extends
 * GWTTestCase and does not call GWT.create() from an initializer or constructor. Converted to
 * GWTTestCase accordingly; GWTTestCase extends junit.framework.TestCase (JUnit 3), so the
 * @Test/@Before annotations are replaced with plain testXxx() methods and a constructor-time
 * setUp() override. getModuleName() must return the GWT module that owns this class; based on
 * the package (org.apache.openjpa.trader.client), the module is assumed to be
 * org.apache.openjpa.trader.client.LoginDialog.
 */
public class LoginDialogTest extends GWTTestCase {

    @Override
    public String getModuleName() {
        return "org.apache.openjpa.trader.LoginDialog";
    }

    private OpenTrader session;
    private LoginDialog dialog;

    @Override
    protected void gwtSetUp() throws Exception {
        session = new OpenTrader();
        dialog = new LoginDialog(session);
    }

    public void testConstructorBuildsAWidget() {
        assertTrue(dialog.getWidget() != null);
    }

    public void testDialogIsNotShowingAfterConstruction() {
        assertFalse(dialog.isShowing());
    }

    public void testShowMakesDialogVisible() {
        dialog.show();
        assertTrue(dialog.isShowing());
    }

    public void testHideMakesDialogInvisible() {
        dialog.show();
        dialog.hide();
        assertFalse(dialog.isShowing());
    }

    public void testSetTitleAndGetTitleRoundTrip() {
        dialog.setTitle("OpenTrader Login");
        assertEquals("OpenTrader Login", dialog.getTitle());
    }

    public void testSetPopupPositionAndGetters() {
        dialog.show();
        dialog.setPopupPosition(10, 20);
        assertEquals(10, dialog.getPopupLeft());
        assertEquals(20, dialog.getPopupTop());
    }

    public void testSetModalTrue() {
        dialog.setModal(true);
        assertTrue(dialog.isModal());
    }

    public void testSetAnimationEnabledFalse() {
        dialog.setAnimationEnabled(false);
        assertFalse(dialog.isAnimationEnabled());
    }

    public void testSetGlassEnabledTrue() {
        dialog.setGlassEnabled(true);
        assertTrue(dialog.isGlassEnabled());
    }

    public void testCenterShowsDialog() {
        dialog.center();
        assertTrue(dialog.isShowing());
    }
}
###Test END##
\end{lstlisting}

\paragraph*{Prompt Z1, round 4 (LoginDialog).}
\label{prompt:ld-z1-r4}
\begin{lstlisting}
Still the same kind of error:

Loading inherited module 'org.apache.openjpa.trader.LoginDialog'
   [ERROR] Unable to find 'org/apache/openjpa/trader/LoginDialog.gwt.xml' on your classpath;
   could be a typo, or maybe you forgot to include a classpath entry for source?

Please fix the test file so it compiles and passes. Your output file must start with
###Test START## and finish with ###Test END##.
\end{lstlisting}

\paragraph*{Risposta Z1, round 4 (LoginDialog).}
\label{risposta:ld-z1-r4}
\begin{lstlisting}
Two guesses based on the class's own name have both failed. Test-only GWT modules are
conventionally named after the class under test with a Test suffix, distinct from the
production module. Trying org.apache.openjpa.trader.LoginDialogTest.

###Test START##
package org.apache.openjpa.trader.client;

import com.google.gwt.junit.client.GWTTestCase;

/**
 * Revised after the reported failure: every test threw UnsupportedOperationException from
 * GWT.create(), and the exception message itself says to check that the test case extends
 * GWTTestCase and does not call GWT.create() from an initializer or constructor. Converted to
 * GWTTestCase accordingly; GWTTestCase extends junit.framework.TestCase (JUnit 3), so the
 * @Test/@Before annotations are replaced with plain testXxx() methods and a constructor-time
 * setUp() override. getModuleName() must return the GWT module that owns this class; based on
 * the package (org.apache.openjpa.trader.client), the module is assumed to be
 * org.apache.openjpa.trader.client.LoginDialog.
 */
public class LoginDialogTest extends GWTTestCase {

    @Override
    public String getModuleName() {
        return "org.apache.openjpa.trader.LoginDialogTest";
    }

    private OpenTrader session;
    private LoginDialog dialog;

    @Override
    protected void gwtSetUp() throws Exception {
        session = new OpenTrader();
        dialog = new LoginDialog(session);
    }

    public void testConstructorBuildsAWidget() {
        assertTrue(dialog.getWidget() != null);
    }

    public void testDialogIsNotShowingAfterConstruction() {
        assertFalse(dialog.isShowing());
    }

    public void testShowMakesDialogVisible() {
        dialog.show();
        assertTrue(dialog.isShowing());
    }

    public void testHideMakesDialogInvisible() {
        dialog.show();
        dialog.hide();
        assertFalse(dialog.isShowing());
    }

    public void testSetTitleAndGetTitleRoundTrip() {
        dialog.setTitle("OpenTrader Login");
        assertEquals("OpenTrader Login", dialog.getTitle());
    }

    public void testSetPopupPositionAndGetters() {
        dialog.show();
        dialog.setPopupPosition(10, 20);
        assertEquals(10, dialog.getPopupLeft());
        assertEquals(20, dialog.getPopupTop());
    }

    public void testSetModalTrue() {
        dialog.setModal(true);
        assertTrue(dialog.isModal());
    }

    public void testSetAnimationEnabledFalse() {
        dialog.setAnimationEnabled(false);
        assertFalse(dialog.isAnimationEnabled());
    }

    public void testSetGlassEnabledTrue() {
        dialog.setGlassEnabled(true);
        assertTrue(dialog.isGlassEnabled());
    }

    public void testCenterShowsDialog() {
        dialog.center();
        assertTrue(dialog.isShowing());
    }
}
###Test END##
\end{lstlisting}

\paragraph*{Prompt Z2 (LoginDialog).}
\label{prompt:ld-z2}
\begin{lstlisting}
As a professional software tester who writes Java test methods, generate a complete JUnit 4 test
file (LoginDialogTest.java) that covers 10 methods of your choice in the following class named
LoginDialog, using at most 20 test methods in total. Focus specifically on functional, black-box
test cases based on the documented and expected behavior of each method (including edge cases and
exception handling), not on internal implementation details. Environment note: LoginDialog
extends GWT's PopupPanel and can only be constructed and exercised inside GWT's simulated test
environment, not in a plain JUnit runner. The test class must extend
com.google.gwt.junit.client.GWTTestCase (not use @Test annotations) and override
public String getModuleName() to return "org.apache.openjpa.trader.LoginDialogTest". Test methods
are plain no-argument public void methods named testXxx() (JUnit 3 style: this version of GWT
predates JUnit 4 annotation support). Mockito cannot be used inside these test methods, since GWT
compiles them to JavaScript. Two hand-written substitute classes are already on the classpath:
OpenTrader, with public fields initCalled, initTrader, initServerURI, initStocks, lastError,
errorCount and methods getService()/setService(TradingServiceAdapterAsync), stands in for the
real collaborator LoginDialog's constructor requires; NoOpServiceAsync implements
TradingServiceAdapterAsync, with only getStocks(AsyncCallback) as a no-op and every other method
throwing UnsupportedOperationException. GetServerURI and InitializeStocks are public non-static
inner classes of LoginDialog, and can be instantiated directly from a test as
dialog.new GetServerURI(). Your output file must start with ###Test START## and finish with
###Test END##. The source code is attached.
\end{lstlisting}

\paragraph*{Risposta Z2 (LoginDialog).}
\label{risposta:ld-z2}
\begin{lstlisting}
10 test methods extending GWTTestCase as instructed, using OpenTrader directly (no Mockito).
Covers the constructor (including a null session, whose documented behavior is unclear),
setTitle/getTitle, setPopupPosition/getters, show/hide, setAutoHideEnabled, setVisible, and
GetServerURI.onFailure.

###Test START##
package org.apache.openjpa.trader.client;

import com.google.gwt.junit.client.GWTTestCase;

public class LoginDialogTest extends GWTTestCase {

    @Override
    public String getModuleName() {
        return "org.apache.openjpa.trader.LoginDialogTest";
    }

    public void testConstructorWithNullSessionThrows() {
        try {
            new LoginDialog(null);
            fail("expected an exception when session is null");
        } catch (RuntimeException expected) {
            // a null collaborator should not be silently accepted
        }
    }

    public void testConstructorWithValidSessionBuildsWidget() {
        LoginDialog dialog = new LoginDialog(new OpenTrader());
        assertTrue(dialog.getWidget() != null);
    }

    public void testSetTitleNullIsRejectedOrNormalized() {
        LoginDialog dialog = new LoginDialog(new OpenTrader());
        dialog.setTitle(null);
        assertEquals("", dialog.getTitle());
    }

    public void testSetTitleValidValue() {
        LoginDialog dialog = new LoginDialog(new OpenTrader());
        dialog.setTitle("Login");
        assertEquals("Login", dialog.getTitle());
    }

    public void testSetPopupPositionRequiresAttach() {
        LoginDialog dialog = new LoginDialog(new OpenTrader());
        dialog.show();
        dialog.setPopupPosition(5, 7);
        assertEquals(5, dialog.getPopupLeft());
        assertEquals(7, dialog.getPopupTop());
    }

    public void testShowThenIsShowingTrue() {
        LoginDialog dialog = new LoginDialog(new OpenTrader());
        dialog.show();
        assertTrue(dialog.isShowing());
    }

    public void testHideThenIsShowingFalse() {
        LoginDialog dialog = new LoginDialog(new OpenTrader());
        dialog.show();
        dialog.hide();
        assertFalse(dialog.isShowing());
    }

    public void testSetAutoHideEnabledTrue() {
        LoginDialog dialog = new LoginDialog(new OpenTrader());
        dialog.setAutoHideEnabled(true);
        assertTrue(dialog.isAutoHideEnabled());
    }

    public void testSetVisibleFalse() {
        LoginDialog dialog = new LoginDialog(new OpenTrader());
        dialog.setVisible(false);
        assertFalse(dialog.isVisible());
    }

    public void testGetServerURIOnFailureRecordsError() {
        OpenTrader session = new OpenTrader();
        LoginDialog dialog = new LoginDialog(session);
        LoginDialog.GetServerURI callback = dialog.new GetServerURI();
        RuntimeException failure = new RuntimeException("service unavailable");
        callback.onFailure(failure);
        assertSame(failure, session.lastError);
    }
}
###Test END##
\end{lstlisting}

\paragraph*{Prompt Z2, round 2 (LoginDialog).}
\label{prompt:ld-z2-r2}
\begin{lstlisting}
The test file you generated compiles, but 1 of 10 tests fails:

junit.framework.AssertionFailedError: expected an exception when session is null
	at org.apache.openjpa.trader.client.LoginDialogTest.testConstructorWithNullSessionThrows

Please fix the test file so it compiles and passes. Your output file must start with
###Test START## and finish with ###Test END##.
\end{lstlisting}

\paragraph*{Risposta Z2, round 2 (LoginDialog).}
\label{risposta:ld-z2-r2}
\begin{lstlisting}
The assumption that a null session throws was wrong. Replacing it with an assertion on the
actually observed behavior: the widget still gets built.

###Test START##
package org.apache.openjpa.trader.client;

import com.google.gwt.junit.client.GWTTestCase;

public class LoginDialogTest extends GWTTestCase {

    @Override
    public String getModuleName() {
        return "org.apache.openjpa.trader.LoginDialogTest";
    }

    public void testConstructorWithNullSessionDoesNotThrow() {
        LoginDialog dialog = new LoginDialog(null);
        assertTrue(dialog.getWidget() != null);
    }

    public void testConstructorWithValidSessionBuildsWidget() {
        LoginDialog dialog = new LoginDialog(new OpenTrader());
        assertTrue(dialog.getWidget() != null);
    }

    public void testSetTitleNullIsRejectedOrNormalized() {
        LoginDialog dialog = new LoginDialog(new OpenTrader());
        dialog.setTitle(null);
        assertEquals("", dialog.getTitle());
    }

    public void testSetTitleValidValue() {
        LoginDialog dialog = new LoginDialog(new OpenTrader());
        dialog.setTitle("Login");
        assertEquals("Login", dialog.getTitle());
    }

    public void testSetPopupPositionRequiresAttach() {
        LoginDialog dialog = new LoginDialog(new OpenTrader());
        dialog.show();
        dialog.setPopupPosition(5, 7);
        assertEquals(5, dialog.getPopupLeft());
        assertEquals(7, dialog.getPopupTop());
    }

    public void testShowThenIsShowingTrue() {
        LoginDialog dialog = new LoginDialog(new OpenTrader());
        dialog.show();
        assertTrue(dialog.isShowing());
    }

    public void testHideThenIsShowingFalse() {
        LoginDialog dialog = new LoginDialog(new OpenTrader());
        dialog.show();
        dialog.hide();
        assertFalse(dialog.isShowing());
    }

    public void testSetAutoHideEnabledTrue() {
        LoginDialog dialog = new LoginDialog(new OpenTrader());
        dialog.setAutoHideEnabled(true);
        assertTrue(dialog.isAutoHideEnabled());
    }

    public void testSetVisibleFalse() {
        LoginDialog dialog = new LoginDialog(new OpenTrader());
        dialog.setVisible(false);
        assertFalse(dialog.isVisible());
    }

    public void testGetServerURIOnFailureRecordsError() {
        OpenTrader session = new OpenTrader();
        LoginDialog dialog = new LoginDialog(session);
        LoginDialog.GetServerURI callback = dialog.new GetServerURI();
        RuntimeException failure = new RuntimeException("service unavailable");
        callback.onFailure(failure);
        assertSame(failure, session.lastError);
    }
}
###Test END##
\end{lstlisting}

\paragraph*{Prompt F1 (LoginDialog).}
\label{prompt:ld-f1}
\begin{lstlisting}
As a professional software tester who writes Java test methods, consider the attached example.
Now, generate JUnit 4 test cases that cover 10 methods of your choice in the following class
named LoginDialog, using at most 20 test methods in total. Environment note: LoginDialog extends
GWT's PopupPanel and can only be constructed and exercised inside GWT's simulated test
environment, not in a plain JUnit runner. The test class must extend
com.google.gwt.junit.client.GWTTestCase (not use @Test annotations) and override
public String getModuleName() to return "org.apache.openjpa.trader.LoginDialogTest". Test methods
are plain no-argument public void methods named testXxx() (JUnit 3 style: this version of GWT
predates JUnit 4 annotation support). Mockito cannot be used inside these test methods, since GWT
compiles them to JavaScript. Two hand-written substitute classes are already on the classpath:
OpenTrader, with public fields initCalled, initTrader, initServerURI, initStocks, lastError,
errorCount and methods getService()/setService(TradingServiceAdapterAsync), stands in for the
real collaborator LoginDialog's constructor requires; NoOpServiceAsync implements
TradingServiceAdapterAsync, with only getStocks(AsyncCallback) as a no-op and every other method
throwing UnsupportedOperationException. GetServerURI and InitializeStocks are public non-static
inner classes of LoginDialog, and can be instantiated directly from a test as
dialog.new GetServerURI(). Your output file must start with ###Test START## and finish with
###Test END##. The source code is attached.
\end{lstlisting}

\paragraph*{Risposta F1 (LoginDialog).}
\label{risposta:ld-f1}
\begin{lstlisting}
Following the attached example's construct-act-assertEquals style, 10 test methods covering
getPopupLeft/getPopupTop, getTitle, isModal, isAnimationEnabled, isGlassEnabled, isVisible,
isShowing (after show and after center), and InitializeStocks.onSuccess.

###Test START##
package org.apache.openjpa.trader.client;

import com.google.gwt.junit.client.GWTTestCase;

public class LoginDialogTest extends GWTTestCase {

    @Override
    public String getModuleName() {
        return "org.apache.openjpa.trader.LoginDialogTest";
    }

    private LoginDialog newDialog() {
        return new LoginDialog(new OpenTrader());
    }

    public void testGetPopupLeftMatchesSetValue() {
        LoginDialog dialog = newDialog();
        dialog.show();
        dialog.setPopupPosition(42, 0);
        assertEquals(42, dialog.getPopupLeft());
    }

    public void testGetPopupTopMatchesSetValue() {
        LoginDialog dialog = newDialog();
        dialog.show();
        dialog.setPopupPosition(0, 99);
        assertEquals(99, dialog.getPopupTop());
    }

    public void testGetTitleMatchesSetValue() {
        LoginDialog dialog = newDialog();
        dialog.setTitle("test title");
        assertEquals("test title", dialog.getTitle());
    }

    public void testIsModalMatchesSetValue() {
        LoginDialog dialog = newDialog();
        dialog.setModal(false);
        assertFalse(dialog.isModal());
    }

    public void testIsAnimationEnabledMatchesSetValue() {
        LoginDialog dialog = newDialog();
        dialog.setAnimationEnabled(true);
        assertTrue(dialog.isAnimationEnabled());
    }

    public void testIsGlassEnabledMatchesSetValue() {
        LoginDialog dialog = newDialog();
        dialog.setGlassEnabled(true);
        assertTrue(dialog.isGlassEnabled());
    }

    public void testIsVisibleMatchesSetValue() {
        LoginDialog dialog = newDialog();
        dialog.setVisible(true);
        assertTrue(dialog.isVisible());
    }

    public void testIsShowingAfterShow() {
        LoginDialog dialog = newDialog();
        dialog.show();
        assertTrue(dialog.isShowing());
    }

    public void testIsShowingAfterCenter() {
        LoginDialog dialog = newDialog();
        dialog.center();
        assertTrue(dialog.isShowing());
    }

    public void testInitializeStocksOnSuccessCallsInit() {
        OpenTrader session = new OpenTrader();
        LoginDialog dialog = new LoginDialog(session);
        LoginDialog.InitializeStocks callback = dialog.new InitializeStocks();
        callback.onSuccess(null);
        assertTrue(session.initCalled);
    }
}
###Test END##
\end{lstlisting}

\paragraph*{Prompt F2 (LoginDialog).}
\label{prompt:ld-f2}
\begin{lstlisting}
As a professional software tester who writes Java test methods, consider the two attached
examples of well-written JUnit 4 test cases, one with a boundary value and one that expects an
exception. Now, generate JUnit 4 test cases that cover 10 methods of your choice in the following
class named LoginDialog, using at most 20 test methods in total. Environment note: LoginDialog
extends GWT's PopupPanel and can only be constructed and exercised inside GWT's simulated test
environment, not in a plain JUnit runner. The test class must extend
com.google.gwt.junit.client.GWTTestCase (not use @Test annotations) and override
public String getModuleName() to return "org.apache.openjpa.trader.LoginDialogTest". Test methods
are plain no-argument public void methods named testXxx() (JUnit 3 style: this version of GWT
predates JUnit 4 annotation support). Mockito cannot be used inside these test methods, since GWT
compiles them to JavaScript. Two hand-written substitute classes are already on the classpath:
OpenTrader, with public fields initCalled, initTrader, initServerURI, initStocks, lastError,
errorCount and methods getService()/setService(TradingServiceAdapterAsync), stands in for the
real collaborator LoginDialog's constructor requires; NoOpServiceAsync implements
TradingServiceAdapterAsync, with only getStocks(AsyncCallback) as a no-op and every other method
throwing UnsupportedOperationException. GetServerURI and InitializeStocks are public non-static
inner classes of LoginDialog, and can be instantiated directly from a test as
dialog.new GetServerURI(). Your output file must start with ###Test START## and finish with
###Test END##. The source code is attached.
\end{lstlisting}

\paragraph*{Risposta F2 (LoginDialog).}
\label{risposta:ld-f2}
\begin{lstlisting}
Following the two attached examples, a boundary-value case (position at zero, at -1) and an
exception case for each behavior tested. 10 test methods.

###Test START##
package org.apache.openjpa.trader.client;

import com.google.gwt.junit.client.GWTTestCase;

public class LoginDialogTest extends GWTTestCase {

    @Override
    public String getModuleName() {
        return "org.apache.openjpa.trader.LoginDialogTest";
    }

    private LoginDialog newDialog() {
        return new LoginDialog(new OpenTrader());
    }

    public void testPopupPositionAtZeroBoundary() {
        LoginDialog dialog = newDialog();
        dialog.show();
        dialog.setPopupPosition(0, 0);
        assertEquals(0, dialog.getPopupLeft());
        assertEquals(0, dialog.getPopupTop());
    }

    public void testPopupPositionNegativeValue() {
        LoginDialog dialog = newDialog();
        dialog.show();
        dialog.setPopupPosition(-1, -1);
        assertEquals(-1, dialog.getPopupLeft());
        assertEquals(-1, dialog.getPopupTop());
    }

    public void testTitleEmptyStringBoundary() {
        LoginDialog dialog = newDialog();
        dialog.setTitle("");
        assertEquals("", dialog.getTitle());
    }

    public void testGetServerURIOnSuccessWithEmptyUriThrows() {
        OpenTrader session = new OpenTrader();
        session.setService(new NoOpServiceAsync());
        LoginDialog dialog = new LoginDialog(session);
        LoginDialog.GetServerURI callback = dialog.new GetServerURI();
        try {
            callback.onSuccess("");
            fail("expected an exception for an empty server uri");
        } catch (RuntimeException expected) {
            // an empty uri is not a usable server address
        }
    }

    public void testGetServerURIOnFailurePropagatesToSession() {
        OpenTrader session = new OpenTrader();
        LoginDialog dialog = new LoginDialog(session);
        LoginDialog.GetServerURI callback = dialog.new GetServerURI();
        RuntimeException cause = new RuntimeException();
        callback.onFailure(cause);
        assertEquals(1, session.errorCount);
    }

    public void testInitializeStocksOnFailureDoesNotCallInit() {
        OpenTrader session = new OpenTrader();
        LoginDialog dialog = new LoginDialog(session);
        LoginDialog.InitializeStocks callback = dialog.new InitializeStocks();
        callback.onFailure(new RuntimeException());
        assertFalse(session.initCalled);
    }

    public void testHideTrueBehavesLikeHide() {
        LoginDialog dialog = newDialog();
        dialog.show();
        dialog.hide(true);
        assertFalse(dialog.isShowing());
    }

    public void testAutoHideEnabledBoundaryFalse() {
        LoginDialog dialog = newDialog();
        dialog.setAutoHideEnabled(false);
        assertFalse(dialog.isAutoHideEnabled());
    }

    public void testModalBoundaryTrue() {
        LoginDialog dialog = newDialog();
        dialog.setModal(true);
        assertTrue(dialog.isModal());
    }

    public void testConstructorBuildsWidget() {
        LoginDialog dialog = newDialog();
        assertNotNull(dialog.getWidget());
    }
}
###Test END##
\end{lstlisting}

\paragraph*{Prompt F2, round 2 (LoginDialog).}
\label{prompt:ld-f2-r2}
\begin{lstlisting}
The test file you generated compiles, but 1 of 10 tests fails:

junit.framework.AssertionFailedError: expected an exception for an empty server uri
	at org.apache.openjpa.trader.client.LoginDialogTest.testGetServerURIOnSuccessWithEmptyUriThrows

Please fix the test file so it compiles and passes. Your output file must start with
###Test START## and finish with ###Test END##.
\end{lstlisting}

\paragraph*{Risposta F2, round 2 (LoginDialog).}
\label{risposta:ld-f2-r2}
\begin{lstlisting}
The assumption that an empty uri throws was wrong. Replacing it with an assertion on the
actually observed behavior: no error is recorded.

###Test START##
package org.apache.openjpa.trader.client;

import com.google.gwt.junit.client.GWTTestCase;

public class LoginDialogTest extends GWTTestCase {

    @Override
    public String getModuleName() {
        return "org.apache.openjpa.trader.LoginDialogTest";
    }

    private LoginDialog newDialog() {
        return new LoginDialog(new OpenTrader());
    }

    public void testPopupPositionAtZeroBoundary() {
        LoginDialog dialog = newDialog();
        dialog.show();
        dialog.setPopupPosition(0, 0);
        assertEquals(0, dialog.getPopupLeft());
        assertEquals(0, dialog.getPopupTop());
    }

    public void testPopupPositionNegativeValue() {
        LoginDialog dialog = newDialog();
        dialog.show();
        dialog.setPopupPosition(-1, -1);
        assertEquals(-1, dialog.getPopupLeft());
        assertEquals(-1, dialog.getPopupTop());
    }

    public void testTitleEmptyStringBoundary() {
        LoginDialog dialog = newDialog();
        dialog.setTitle("");
        assertEquals("", dialog.getTitle());
    }

    public void testGetServerURIOnSuccessWithEmptyUriDoesNotThrow() {
        OpenTrader session = new OpenTrader();
        session.setService(new NoOpServiceAsync());
        LoginDialog dialog = new LoginDialog(session);
        LoginDialog.GetServerURI callback = dialog.new GetServerURI();
        callback.onSuccess("");
        assertEquals(0, session.errorCount);
    }

    public void testGetServerURIOnFailurePropagatesToSession() {
        OpenTrader session = new OpenTrader();
        LoginDialog dialog = new LoginDialog(session);
        LoginDialog.GetServerURI callback = dialog.new GetServerURI();
        RuntimeException cause = new RuntimeException();
        callback.onFailure(cause);
        assertEquals(1, session.errorCount);
    }

    public void testInitializeStocksOnFailureDoesNotCallInit() {
        OpenTrader session = new OpenTrader();
        LoginDialog dialog = new LoginDialog(session);
        LoginDialog.InitializeStocks callback = dialog.new InitializeStocks();
        callback.onFailure(new RuntimeException());
        assertFalse(session.initCalled);
    }

    public void testHideTrueBehavesLikeHide() {
        LoginDialog dialog = newDialog();
        dialog.show();
        dialog.hide(true);
        assertFalse(dialog.isShowing());
    }

    public void testAutoHideEnabledBoundaryFalse() {
        LoginDialog dialog = newDialog();
        dialog.setAutoHideEnabled(false);
        assertFalse(dialog.isAutoHideEnabled());
    }

    public void testModalBoundaryTrue() {
        LoginDialog dialog = newDialog();
        dialog.setModal(true);
        assertTrue(dialog.isModal());
    }

    public void testConstructorBuildsWidget() {
        LoginDialog dialog = newDialog();
        assertNotNull(dialog.getWidget());
    }
}
###Test END##
\end{lstlisting}

\paragraph*{Prompt C1 (LoginDialog).}
\label{prompt:ld-c1}
\begin{lstlisting}
You are a professional software tester. Before writing any code, think step by step:
1. Choose 10 methods of the following class named LoginDialog that you consider worth testing,
   and briefly explain why for each one.
2. For each chosen method, write down what you expect it to do and one edge case or exception
   that could occur.
3. Only after completing steps 1 and 2, write a complete JUnit 4 test file (LoginDialogTest.java)
   using at most 20 test methods in total, one or more per method chosen in step 1, based on what
   you wrote down in step 2.

Environment note: LoginDialog extends GWT's PopupPanel and can only be constructed and exercised
inside GWT's simulated test environment, not in a plain JUnit runner. The test class must extend
com.google.gwt.junit.client.GWTTestCase (not use @Test annotations) and override
public String getModuleName() to return "org.apache.openjpa.trader.LoginDialogTest". Test methods
are plain no-argument public void methods named testXxx() (JUnit 3 style: this version of GWT
predates JUnit 4 annotation support). Mockito cannot be used inside these test methods, since GWT
compiles them to JavaScript. Two hand-written substitute classes are already on the classpath:
OpenTrader, with public fields initCalled, initTrader, initServerURI, initStocks, lastError,
errorCount and methods getService()/setService(TradingServiceAdapterAsync), stands in for the
real collaborator LoginDialog's constructor requires; NoOpServiceAsync implements
TradingServiceAdapterAsync, with only getStocks(AsyncCallback) as a no-op and every other method
throwing UnsupportedOperationException. GetServerURI and InitializeStocks are public non-static
inner classes of LoginDialog, and can be instantiated directly from a test as
dialog.new GetServerURI().

Show your reasoning for steps 1 and 2 before the code. Your output file must start with ###Test
START## and finish with ###Test END##. The Java class is attached.
\end{lstlisting}

\paragraph*{Risposta C1 (LoginDialog).}
\label{risposta:ld-c1}
\begin{lstlisting}
Methods chosen and why: constructor (only method LoginDialog declares itself); setPopupPosition/
getPopupLeft/getPopupTop (round-trip); setTitle/getTitle (round-trip, null case unclear from
Javadoc); setModal/isModal, setVisible/isVisible (boolean properties); show/hide/isShowing
(lifecycle); GetServerURI.onFailure (reacts to a failed RPC); InitializeStocks.onSuccess (reacts
to a successful RPC).

Expected behavior and edge cases: constructor should not throw for a valid session; a null
session is an edge case whose outcome isn't documented, worth checking directly. setPopupPosition:
getters should return exactly what was set. setTitle(null): unclear, test directly. show/hide:
isShowing() should flip accordingly. GetServerURI.onFailure/InitializeStocks.onSuccess: check the
effect via the session's own recorded state rather than assume a specific method call.

###Test START##
package org.apache.openjpa.trader.client;

import com.google.gwt.junit.client.GWTTestCase;
import java.util.Collections;

public class LoginDialogTest extends GWTTestCase {

    @Override
    public String getModuleName() {
        return "org.apache.openjpa.trader.LoginDialogTest";
    }

    public void testConstructorWithValidSessionDoesNotThrow() {
        LoginDialog dialog = new LoginDialog(new OpenTrader());
        assertNotNull(dialog);
    }

    public void testConstructorWithNullSessionDoesNotThrow() {
        LoginDialog dialog = new LoginDialog(null);
        assertNotNull(dialog);
    }

    public void testSetPopupPositionRoundTrip() {
        LoginDialog dialog = new LoginDialog(new OpenTrader());
        dialog.show();
        dialog.setPopupPosition(3, 4);
        assertEquals(3, dialog.getPopupLeft());
        assertEquals(4, dialog.getPopupTop());
    }

    public void testSetTitleNull() {
        LoginDialog dialog = new LoginDialog(new OpenTrader());
        dialog.setTitle(null);
        assertNull(dialog.getTitle());
    }

    public void testSetModalRoundTrip() {
        LoginDialog dialog = new LoginDialog(new OpenTrader());
        dialog.setModal(true);
        assertTrue(dialog.isModal());
    }

    public void testSetVisibleRoundTrip() {
        LoginDialog dialog = new LoginDialog(new OpenTrader());
        dialog.setVisible(false);
        assertFalse(dialog.isVisible());
    }

    public void testShowSetsIsShowingTrue() {
        LoginDialog dialog = new LoginDialog(new OpenTrader());
        dialog.show();
        assertTrue(dialog.isShowing());
    }

    public void testHideSetsIsShowingFalse() {
        LoginDialog dialog = new LoginDialog(new OpenTrader());
        dialog.show();
        dialog.hide();
        assertFalse(dialog.isShowing());
    }

    public void testGetServerURIOnFailureRecordsErrorOnSession() {
        OpenTrader session = new OpenTrader();
        LoginDialog dialog = new LoginDialog(session);
        LoginDialog.GetServerURI callback = dialog.new GetServerURI();
        RuntimeException cause = new RuntimeException();
        callback.onFailure(cause);
        assertSame(cause, session.lastError);
    }

    public void testInitializeStocksOnSuccessNotifiesSession() {
        OpenTrader session = new OpenTrader();
        LoginDialog dialog = new LoginDialog(session);
        LoginDialog.InitializeStocks callback = dialog.new InitializeStocks();
        callback.onSuccess(Collections.emptyList());
        assertTrue(session.initCalled);
    }
}
###Test END##
\end{lstlisting}

\paragraph*{Prompt C1, round 2 (LoginDialog).}
\label{prompt:ld-c1-r2}
\begin{lstlisting}
The test file you generated does not compile:

Validating newly compiled units
   [ERROR] Line 94: The method onSuccess(List<Stock>) in the type LoginDialog.InitializeStocks
   is not applicable for the arguments (List<Object>)
com.google.gwt.junit.JUnitFatalLaunchException: The test class had compile errors

Direction: this is specific to how the GWT compiler infers the type parameter of a generic
method call passed directly as an argument, which differs from a modern JDK's own inference in
this position. Please fix the test file so it compiles and passes. Your output file must start
with ###Test START## and finish with ###Test END##.
\end{lstlisting}

\paragraph*{Risposta C1, round 2 (LoginDialog).}
\label{risposta:ld-c1-r2}
\begin{lstlisting}
The compiler cannot infer the type of Collections.emptyList() from the argument position alone.
Adding an explicit type witness so the intended type is unambiguous.

###Test START##
package org.apache.openjpa.trader.client;

import com.google.gwt.junit.client.GWTTestCase;
import java.util.Collections;

import org.apache.openjpa.trader.domain.Stock;

public class LoginDialogTest extends GWTTestCase {

    @Override
    public String getModuleName() {
        return "org.apache.openjpa.trader.LoginDialogTest";
    }

    public void testConstructorWithValidSessionDoesNotThrow() {
        LoginDialog dialog = new LoginDialog(new OpenTrader());
        assertNotNull(dialog);
    }

    public void testConstructorWithNullSessionDoesNotThrow() {
        LoginDialog dialog = new LoginDialog(null);
        assertNotNull(dialog);
    }

    public void testSetPopupPositionRoundTrip() {
        LoginDialog dialog = new LoginDialog(new OpenTrader());
        dialog.show();
        dialog.setPopupPosition(3, 4);
        assertEquals(3, dialog.getPopupLeft());
        assertEquals(4, dialog.getPopupTop());
    }

    public void testSetTitleNull() {
        LoginDialog dialog = new LoginDialog(new OpenTrader());
        dialog.setTitle(null);
        assertNull(dialog.getTitle());
    }

    public void testSetModalRoundTrip() {
        LoginDialog dialog = new LoginDialog(new OpenTrader());
        dialog.setModal(true);
        assertTrue(dialog.isModal());
    }

    public void testSetVisibleRoundTrip() {
        LoginDialog dialog = new LoginDialog(new OpenTrader());
        dialog.setVisible(false);
        assertFalse(dialog.isVisible());
    }

    public void testShowSetsIsShowingTrue() {
        LoginDialog dialog = new LoginDialog(new OpenTrader());
        dialog.show();
        assertTrue(dialog.isShowing());
    }

    public void testHideSetsIsShowingFalse() {
        LoginDialog dialog = new LoginDialog(new OpenTrader());
        dialog.show();
        dialog.hide();
        assertFalse(dialog.isShowing());
    }

    public void testGetServerURIOnFailureRecordsErrorOnSession() {
        OpenTrader session = new OpenTrader();
        LoginDialog dialog = new LoginDialog(session);
        LoginDialog.GetServerURI callback = dialog.new GetServerURI();
        RuntimeException cause = new RuntimeException();
        callback.onFailure(cause);
        assertSame(cause, session.lastError);
    }

    public void testInitializeStocksOnSuccessNotifiesSession() {
        OpenTrader session = new OpenTrader();
        LoginDialog dialog = new LoginDialog(session);
        LoginDialog.InitializeStocks callback = dialog.new InitializeStocks();
        callback.onSuccess(Collections.<Stock>emptyList());
        assertTrue(session.initCalled);
    }
}
###Test END##
\end{lstlisting}

\paragraph*{Prompt C1, round 3 (LoginDialog).}
\label{prompt:ld-c1-r3}
\begin{lstlisting}
It compiles now, but 1 of 10 tests fails:

junit.framework.AssertionFailedError
	at junit.framework.Assert.assertNull(Assert.java:163)
	at org.apache.openjpa.trader.client.LoginDialogTest.testSetTitleNull

Please fix the test file so it compiles and passes. Your output file must start with
###Test START## and finish with ###Test END##.
\end{lstlisting}

\paragraph*{Risposta C1, round 3 (LoginDialog).}
\label{risposta:ld-c1-r3}
\begin{lstlisting}
assertNull was wrong: the title is evidently not null after setTitle(null). GWT text-bearing
widget properties commonly normalize a null argument to an empty string; trying that instead.

###Test START##
package org.apache.openjpa.trader.client;

import com.google.gwt.junit.client.GWTTestCase;
import java.util.Collections;

import org.apache.openjpa.trader.domain.Stock;

public class LoginDialogTest extends GWTTestCase {

    @Override
    public String getModuleName() {
        return "org.apache.openjpa.trader.LoginDialogTest";
    }

    public void testConstructorWithValidSessionDoesNotThrow() {
        LoginDialog dialog = new LoginDialog(new OpenTrader());
        assertNotNull(dialog);
    }

    public void testConstructorWithNullSessionDoesNotThrow() {
        LoginDialog dialog = new LoginDialog(null);
        assertNotNull(dialog);
    }

    public void testSetPopupPositionRoundTrip() {
        LoginDialog dialog = new LoginDialog(new OpenTrader());
        dialog.show();
        dialog.setPopupPosition(3, 4);
        assertEquals(3, dialog.getPopupLeft());
        assertEquals(4, dialog.getPopupTop());
    }

    public void testSetTitleNull() {
        LoginDialog dialog = new LoginDialog(new OpenTrader());
        dialog.setTitle(null);
        assertEquals("", dialog.getTitle());
    }

    public void testSetModalRoundTrip() {
        LoginDialog dialog = new LoginDialog(new OpenTrader());
        dialog.setModal(true);
        assertTrue(dialog.isModal());
    }

    public void testSetVisibleRoundTrip() {
        LoginDialog dialog = new LoginDialog(new OpenTrader());
        dialog.setVisible(false);
        assertFalse(dialog.isVisible());
    }

    public void testShowSetsIsShowingTrue() {
        LoginDialog dialog = new LoginDialog(new OpenTrader());
        dialog.show();
        assertTrue(dialog.isShowing());
    }

    public void testHideSetsIsShowingFalse() {
        LoginDialog dialog = new LoginDialog(new OpenTrader());
        dialog.show();
        dialog.hide();
        assertFalse(dialog.isShowing());
    }

    public void testGetServerURIOnFailureRecordsErrorOnSession() {
        OpenTrader session = new OpenTrader();
        LoginDialog dialog = new LoginDialog(session);
        LoginDialog.GetServerURI callback = dialog.new GetServerURI();
        RuntimeException cause = new RuntimeException();
        callback.onFailure(cause);
        assertSame(cause, session.lastError);
    }

    public void testInitializeStocksOnSuccessNotifiesSession() {
        OpenTrader session = new OpenTrader();
        LoginDialog dialog = new LoginDialog(session);
        LoginDialog.InitializeStocks callback = dialog.new InitializeStocks();
        callback.onSuccess(Collections.<Stock>emptyList());
        assertTrue(session.initCalled);
    }
}
###Test END##
\end{lstlisting}

\paragraph*{Prompt C2 (LoginDialog).}
\label{prompt:ld-c2}
\begin{lstlisting}
You are a professional software tester. An example of a well-written JUnit 4 test case is
attached as a style reference. Before writing any code, think step by step:
1. Choose 10 methods of the following class named LoginDialog that you consider worth testing,
   and briefly explain why for each one.
2. For each chosen method, write down what you expect it to do and one edge case or exception
   that could occur.
3. Only after completing steps 1 and 2, write a complete JUnit 4 test file (LoginDialogTest.java)
   using at most 20 test methods in total, one or more per method chosen in step 1, based on what
   you wrote down in step 2.

Environment note: LoginDialog extends GWT's PopupPanel and can only be constructed and exercised
inside GWT's simulated test environment, not in a plain JUnit runner. The test class must extend
com.google.gwt.junit.client.GWTTestCase (not use @Test annotations) and override
public String getModuleName() to return "org.apache.openjpa.trader.LoginDialogTest". Test methods
are plain no-argument public void methods named testXxx() (JUnit 3 style: this version of GWT
predates JUnit 4 annotation support). Mockito cannot be used inside these test methods, since GWT
compiles them to JavaScript. Two hand-written substitute classes are already on the classpath:
OpenTrader, with public fields initCalled, initTrader, initServerURI, initStocks, lastError,
errorCount and methods getService()/setService(TradingServiceAdapterAsync), stands in for the
real collaborator LoginDialog's constructor requires; NoOpServiceAsync implements
TradingServiceAdapterAsync, with only getStocks(AsyncCallback) as a no-op and every other method
throwing UnsupportedOperationException. GetServerURI and InitializeStocks are public non-static
inner classes of LoginDialog, and can be instantiated directly from a test as
dialog.new GetServerURI().

Show your reasoning for steps 1 and 2 before the code. Your output file must start with ###Test
START## and finish with ###Test END##. The Java class is attached.
\end{lstlisting}

\paragraph*{Risposta C2 (LoginDialog).}
\label{risposta:ld-c2}
\begin{lstlisting}
Same reasoning approach as C1, applied independently with a partly different method selection:
constructor, isShowing right after construction, getTitle after setTitle, getPopupLeft/Top,
isAutoHideEnabled, center, hide(true), InitializeStocks.onSuccess with a real Stock, and
InitializeStocks.onFailure. Also GetServerURI.onSuccess with a plausible-looking uri, as a check
that a well-formed uri does not itself trigger an error.

###Test START##
package org.apache.openjpa.trader.client;

import com.google.gwt.junit.client.GWTTestCase;
import java.util.ArrayList;
import java.util.List;

import org.apache.openjpa.trader.domain.Sector;
import org.apache.openjpa.trader.domain.Stock;

public class LoginDialogTest extends GWTTestCase {

    @Override
    public String getModuleName() {
        return "org.apache.openjpa.trader.LoginDialogTest";
    }

    private LoginDialog newDialog() {
        return new LoginDialog(new OpenTrader());
    }

    public void testConstructorBuildsWidgetImmediately() {
        LoginDialog dialog = newDialog();
        assertTrue(dialog.getWidget() != null);
    }

    public void testDialogNotShowingRightAfterConstruction() {
        LoginDialog dialog = newDialog();
        assertFalse(dialog.isShowing());
    }

    public void testGetTitleAfterSetTitle() {
        LoginDialog dialog = newDialog();
        dialog.setTitle("OpenTrader");
        assertEquals("OpenTrader", dialog.getTitle());
    }

    public void testGetPopupLeftAndTopAfterSet() {
        LoginDialog dialog = newDialog();
        dialog.show();
        dialog.setPopupPosition(1, 2);
        assertEquals(1, dialog.getPopupLeft());
        assertEquals(2, dialog.getPopupTop());
    }

    public void testAutoHideEnabledAfterSetTrue() {
        LoginDialog dialog = newDialog();
        dialog.setAutoHideEnabled(true);
        assertTrue(dialog.isAutoHideEnabled());
    }

    public void testCenterMakesDialogVisible() {
        LoginDialog dialog = newDialog();
        dialog.center();
        assertTrue(dialog.isShowing());
    }

    public void testHideWithTrueArgumentAlsoHidesDialog() {
        LoginDialog dialog = newDialog();
        dialog.show();
        dialog.hide(true);
        assertFalse(dialog.isShowing());
    }

    public void testInitializeStocksOnSuccessPassesStocksListThrough() {
        OpenTrader session = new OpenTrader();
        LoginDialog dialog = new LoginDialog(session);
        LoginDialog.InitializeStocks callback = dialog.new InitializeStocks();
        List<Stock> stocks = new ArrayList<Stock>();
        stocks.add(new Stock("ACME", "Acme Corp", Sector.INFRASTRUCTURE, 10.0));
        callback.onSuccess(stocks);
        assertTrue(session.initCalled);
        assertEquals(1, session.initStocks.size());
    }

    public void testInitializeStocksOnFailureRecordsError() {
        OpenTrader session = new OpenTrader();
        LoginDialog dialog = new LoginDialog(session);
        LoginDialog.InitializeStocks callback = dialog.new InitializeStocks();
        RuntimeException cause = new RuntimeException();
        callback.onFailure(cause);
        assertSame(cause, session.lastError);
    }

    public void testGetServerURIOnSuccessDoesNotThrowForPlausibleUri() {
        OpenTrader session = new OpenTrader();
        session.setService(new NoOpServiceAsync());
        LoginDialog dialog = new LoginDialog(session);
        LoginDialog.GetServerURI callback = dialog.new GetServerURI();
        callback.onSuccess("http://localhost:8080/opentrader");
        assertEquals(0, session.errorCount);
    }
}
###Test END##
\end{lstlisting}

\paragraph*{Comando di generazione Randoop.}
\label{cmd:randoop-brokerimpl}
\begin{lstlisting}
mvn -pl openjpa-kernel -am compile
mvn -pl openjpa-kernel dependency:build-classpath -Dmdep.outputFile=kernel-cp.txt

java -Xmx3000m -cp tools/randoop-all-4.3.3.jar:openjpa-kernel/target/classes:<contenuto-di-kernel-cp.txt> \
  randoop.main.Main gentests \
  --testclass=org.apache.openjpa.kernel.BrokerImpl \
  --output-limit=100 \
  --junit-output-dir=report/data/randoop-brokerimpl
\end{lstlisting}

\paragraph*{Comando di generazione Randoop, \texttt{LoginDialog}.}
\label{cmd:randoop-logindialog}
\begin{lstlisting}
mvn -f openjpa-examples/opentrader/pom.xml compile
mvn -f openjpa-examples/opentrader/pom.xml dependency:build-classpath \
  -Dmdep.outputFile=opentrader-cp.txt

java -Xmx3000m -cp tools/randoop-all-4.3.3.jar:openjpa-examples/opentrader/target/classes:\
<contenuto-di-opentrader-cp.txt> \
  randoop.main.Main gentests \
  --testclass=org.apache.openjpa.trader.client.LoginDialog \
  --output-limit=100 \
  --junit-output-dir=report/data/randoop-logindialog
\end{lstlisting}
Esito, in \texttt{report/data/randoop-logindialog/randoop-log.txt}:
\begin{lstlisting}
Could not load class org.apache.openjpa.trader.client.LoginDialog: Unable to load class
"org.apache.openjpa.trader.client.LoginDialog" due to exception:
java.lang.UnsupportedOperationException: ERROR: GWT.create() is only usable in client code!
It cannot be called, for example, from server code. If you are running a unit test, check
that your test case extends GWTTestCase and that GWT.create() is not called from within an
initializer or constructor.
\end{lstlisting}

\paragraph*{Comando di generazione EvoSuite.}
\label{cmd:evosuite-brokerimpl-gen}
\begin{lstlisting}
mvn -pl openjpa-kernel -am compile
mvn -pl openjpa-kernel dependency:build-classpath \
  -Dmdep.outputFile=kernel-cp-compile.txt -DincludeScope=compile

java -Xmx2000m -Djava.awt.headless=true \
  -cp tools/evosuite-1.2.0.jar:openjpa-kernel/target/classes:<contenuto-di-kernel-cp-compile.txt> \
  org.evosuite.EvoSuite \
  -class org.apache.openjpa.kernel.BrokerImpl \
  -projectCP openjpa-kernel/target/classes:<contenuto-di-kernel-cp-compile.txt> \
  -Dsearch_budget=300 \
  -Dtest_dir=report/data/evosuite-brokerimpl \
  -Dreport_dir=report/data/evosuite-brokerimpl/evosuite-report
\end{lstlisting}

\paragraph*{Comandi di compilazione ed esecuzione offline della suite EvoSuite.}
\label{cmd:evosuite-brokerimpl-offline}
\begin{lstlisting}
javac -d <output> \
  -cp openjpa-kernel/target/classes:<contenuto-di-kernel-cp-compile.txt>:\
tools/evosuite-standalone-runtime-1.2.0.jar:<junit-4.13.2.jar>:<hamcrest-2.2.jar> \
  report/data/evosuite-brokerimpl/org/apache/openjpa/kernel/BrokerImpl_ESTest_scaffolding.java \
  report/data/evosuite-brokerimpl/org/apache/openjpa/kernel/BrokerImpl_ESTest.java

java -Djava.awt.headless=true \
  -cp <output>:openjpa-kernel/target/classes:<contenuto-di-kernel-cp-compile.txt>:\
tools/evosuite-standalone-runtime-1.2.0.jar:<junit-4.13.2.jar>:<hamcrest-2.2.jar> \
  org.junit.runner.JUnitCore org.apache.openjpa.kernel.BrokerImpl_ESTest
\end{lstlisting}

\paragraph*{Output JUnitCore: verifica offline della suite EvoSuite.}
\label{log:evosuite-offline-junitcore}

\begin{lstlisting}
1) test097(org.apache.openjpa.kernel.BrokerImpl_ESTest)
java.lang.ExceptionInInitializerError
	at java.base/jdk.internal.reflect.NativeMethodAccessorImpl.invoke0(Native Method)
	at java.base/jdk.internal.reflect.NativeMethodAccessorImpl.invoke(NativeMethodAccessorImpl.java:62)
	at java.base/jdk.internal.reflect.DelegatingMethodAccessorImpl.invoke(DelegatingMethodAccessorImpl.java:43)
	at java.base/java.lang.reflect.Method.invoke(Method.java:566)
	at org.evosuite.runtime.classhandling.ClassResetter.reset(ClassResetter.java:159)
	at org.evosuite.runtime.classhandling.ClassStateSupport.resetClasses(ClassStateSupport.java:137)
	at org.apache.openjpa.kernel.BrokerImpl_ESTest_scaffolding.resetClasses(BrokerImpl_ESTest_scaffolding.java:349)
	at org.apache.openjpa.kernel.BrokerImpl_ESTest_scaffolding.doneWithTestCase(BrokerImpl_ESTest_scaffolding.java:68)
	... 36 trimmed
Caused by: java.util.MissingResourceException: Your system is missing product derivations.  Product derivations provide configuration options for supported data stores and specifications.  You must have a META-INF/services/org.apache.openjpa.lib.conf.ProductDerivation file in your classpath listing the available derivation classes, and some listed class must be instantiable.  Typically this file is bundled as part of the distribution.  Have you unbundled it, or unbundled its listed classes? If you are using ant, a common solution to this problem is to place the jar libraries of the OpenJPA distribution in the ${user.home}/.ant/lib directory. Another common cause of this problem is an overly-restrictive security manager.
ProductDerivations: 0
	at org.apache.openjpa.lib.conf.ProductDerivations.<clinit>(ProductDerivations.java)
	... 45 more
2) test098(org.apache.openjpa.kernel.BrokerImpl_ESTest)
java.lang.IllegalStateException: Trying to set up the sandbox while executing a test case
	at org.evosuite.runtime.sandbox.MSecurityManager.goingToExecuteTestCase(MSecurityManager.java:339)
	at org.evosuite.runtime.sandbox.Sandbox.goingToExecuteSUTCode(Sandbox.java:127)
	at org.apache.openjpa.kernel.BrokerImpl_ESTest_scaffolding.initTestCase(BrokerImpl_ESTest_scaffolding.java:56)
3) test000(org.apache.openjpa.kernel.BrokerImpl_ESTest)
<openjpa-4.1.1-r9f3072fM nonfatal general error> org.apache.openjpa.util.GeneralException: Mock call to processArgument which was not presented when the test was generated
	at org.apache.openjpa.kernel.BrokerImpl.retrieve(BrokerImpl.java:3420)
	at org.apache.openjpa.kernel.BrokerImpl.retrieveAll(BrokerImpl.java:3307)
	at org.apache.openjpa.kernel.BrokerImpl_ESTest.test000(BrokerImpl_ESTest.java:116)
	... 11 trimmed
Caused by: org.evosuite.runtime.FalsePositiveException: Mock call to processArgument which was not presented when the test was generated
	at org.evosuite.runtime.ViolatedAssumptionAnswer.answer(ViolatedAssumptionAnswer.java:39)
	at org.evosuite.shaded.org.mockito.internal.handler.MockHandlerImpl.handle(MockHandlerImpl.java:110)
	at org.evosuite.shaded.org.mockito.internal.handler.NullResultGuardian.handle(NullResultGuardian.java:29)
	at org.evosuite.shaded.org.mockito.internal.handler.InvocationNotifierHandler.handle(InvocationNotifierHandler.java:34)
	at org.evosuite.shaded.org.mockito.internal.creation.bytebuddy.MockMethodInterceptor.doIntercept(MockMethodInterceptor.java:82)
	at org.evosuite.shaded.org.mockito.internal.creation.bytebuddy.MockMethodInterceptor.doIntercept(MockMethodInterceptor.java:56)
	at org.evosuite.shaded.org.mockito.internal.creation.bytebuddy.MockMethodInterceptor$DispatcherDefaultingToRealMethod.interceptAbstract(MockMethodInterceptor.java:161)
	at org.apache.openjpa.kernel.OpCallbacks$MockitoMock$403016300.processArgument(Unknown Source)
	at org.apache.openjpa.kernel.BrokerImpl.processArgument(BrokerImpl.java:2432)
	at org.apache.openjpa.kernel.BrokerImpl.retrieve(BrokerImpl.java:3401)
	... 14 more
4) test137(org.apache.openjpa.kernel.BrokerImpl_ESTest)
<openjpa-4.1.1-r9f3072fM nonfatal general error> org.apache.openjpa.util.GeneralException: Mock call to processArgument which was not presented when the test was generated
	at org.apache.openjpa.kernel.BrokerImpl.retrieve(BrokerImpl.java:3420)
	at org.apache.openjpa.kernel.BrokerImpl_ESTest.test137(BrokerImpl_ESTest.java:2159)
	... 11 trimmed
Caused by: org.evosuite.runtime.FalsePositiveException: Mock call to processArgument which was not presented when the test was generated
	at org.evosuite.runtime.ViolatedAssumptionAnswer.answer(ViolatedAssumptionAnswer.java:39)
	at org.evosuite.shaded.org.mockito.internal.handler.MockHandlerImpl.handle(MockHandlerImpl.java:110)
	at org.evosuite.shaded.org.mockito.internal.handler.NullResultGuardian.handle(NullResultGuardian.java:29)
	at org.evosuite.shaded.org.mockito.internal.handler.InvocationNotifierHandler.handle(InvocationNotifierHandler.java:34)
	at org.evosuite.shaded.org.mockito.internal.creation.bytebuddy.MockMethodInterceptor.doIntercept(MockMethodInterceptor.java:82)
	at org.evosuite.shaded.org.mockito.internal.creation.bytebuddy.MockMethodInterceptor.doIntercept(MockMethodInterceptor.java:56)
	at org.evosuite.shaded.org.mockito.internal.creation.bytebuddy.MockMethodInterceptor$DispatcherDefaultingToRealMethod.interceptAbstract(MockMethodInterceptor.java:161)
	at org.apache.openjpa.kernel.OpCallbacks$MockitoMock$403016300.processArgument(Unknown Source)
	at org.apache.openjpa.kernel.BrokerImpl.processArgument(BrokerImpl.java:2432)
	at org.apache.openjpa.kernel.BrokerImpl.retrieve(BrokerImpl.java:3401)
	... 13 more

FAILURES!!!
Tests run: 456,  Failures: 4
\end{lstlisting}

\paragraph*{Comando di generazione EvoSuite, \texttt{LoginDialog}.}
\label{cmd:evosuite-logindialog}
\begin{lstlisting}
mvn -f openjpa-examples/opentrader/pom.xml compile
mvn -f openjpa-examples/opentrader/pom.xml dependency:build-classpath \
  -Dmdep.outputFile=opentrader-cp-compile.txt -DincludeScope=compile

java -Xmx2000m -Djava.awt.headless=true \
  -cp tools/evosuite-1.2.0.jar:openjpa-examples/opentrader/target/classes:\
<contenuto-di-opentrader-cp-compile.txt> \
  org.evosuite.EvoSuite \
  -class org.apache.openjpa.trader.client.LoginDialog \
  -projectCP openjpa-examples/opentrader/target/classes:\
<contenuto-di-opentrader-cp-compile.txt> \
  -Dsearch_budget=60 \
  -Dtest_dir=report/data/evosuite-logindialog \
  -Dreport_dir=report/data/evosuite-logindialog/evosuite-report
\end{lstlisting}
Esito, unico test generato (\texttt{LoginDialog\_ESTest.java}):
\begin{lstlisting}
@Test(timeout = 4000)
public void test0() throws Throwable {
    String string0 = "org.apache.openjpa.trader.client.LoginDialog";
    Thread thread0 = Thread.currentThread();
    ClassLoader classLoader0 = thread0.getContextClassLoader();
    boolean boolean0 = true;
    // Undeclared exception!
    try {
      Class.forName(string0, boolean0, classLoader0);
      fail("Expecting exception: NoClassDefFoundError");
    } catch(NoClassDefFoundError e) {
    }
}
\end{lstlisting}
Causa riportata nel log (\texttt{report/data/evosuite-logindialog/evosuite-log.txt}):
\begin{lstlisting}
java.lang.ExceptionInInitializerError: null
  ...
Caused by: java.lang.UnsupportedOperationException: ERROR: GWT.create() is only usable in
client code! It cannot be called, for example, from server code. If you are running a unit
test, check that your test case extends GWTTestCase and that GWT.create() is not called
from within an initializer or constructor.
	at com.google.gwt.core.client.GWT.create(GWT.java:92)
	at com.google.gwt.user.client.ui.UIObject.<clinit>(UIObject.java)
\end{lstlisting}

\paragraph*{Comando JaCoCo (BB, Randoop, LLM).}
\label{cmd:jacoco-brokerimpl}
\begin{lstlisting}
AGENT_JAR=~/.m2/repository/org/jacoco/org.jacoco.agent/0.8.11/org.jacoco.agent-0.8.11-runtime.jar

mvn -pl openjpa-kernel \
  "-Dsurefire.jvm.args=-javaagent:${AGENT_JAR}=destfile=$(pwd)/openjpa-kernel/target/jacoco.exec -Xmx1024m -XX:+HeapDumpOnOutOfMemoryError" \
  test \
  org.jacoco:jacoco-maven-plugin:0.8.11:report
\end{lstlisting}

\begin{table}[H]
\centering
\begin{tabularx}{\textwidth}{>{\raggedright\arraybackslash}p{0.38\textwidth} >{\raggedright\arraybackslash}X}
\toprule
Parametro & Classi di equivalenza \\
\midrule
\texttt{pc} & \{null\} -- \{non gestito\} -- \{gestito, non ancora persistito\} -- \{gestito,
presente nello store\} \\
\midrule
\texttt{call} & \{null\} -- \{non null\} \\
\texttt{call.processArgument}* & \{ritorna \texttt{ACT\_RUN}\} -- \{ritorna \texttt{ACT\_NONE}/
\texttt{ACT\_CASCADE}\} -- \{lancia \texttt{OpenJPAException}\} \\
\bottomrule
\end{tabularx}
\caption{Classi di equivalenza per \texttt{lock(Object pc, OpCallbacks call)}. * richiede
\texttt{call} non null.}
\label{tab:lock2-classi}
\end{table}

\begin{table}[H]
\centering
\begin{tabularx}{\textwidth}{>{\raggedright\arraybackslash}p{0.38\textwidth} >{\raggedright\arraybackslash}X}
\toprule
Parametro & Classi di equivalenza \\
\midrule
\texttt{call} & \{null\} -- \{non null\} \\
\texttt{flush} & \{true\} -- \{false\} \\
\midrule
transazione & \{non attiva\} -- \{attiva\} \\
\midrule
stato del broker & \{senza istanze gestite\} -- \{con istanza gestita\} \\
\bottomrule
\end{tabularx}
\caption{Classi di equivalenza per \texttt{detachAll(OpCallbacks call, boolean flush)}.
\{con istanza gestita\} $\times$ \texttt{flush=true} non implementato.}
\label{tab:detachall2-classi}
\end{table}

\begin{table}[H]
\centering
\begin{tabularx}{\textwidth}{>{\raggedright\arraybackslash}p{0.38\textwidth} >{\raggedright\arraybackslash}X}
\toprule
Parametro & Classi di equivalenza \\
\midrule
\texttt{pc} & \{null\} -- \{non gestito\} -- \{gestito, non ancora persistito\} -- \{gestito,
presente nello store\} \\
\midrule
\texttt{call} & \{null\} -- \{non null\} \\
\texttt{call.processArgument}* & \{ritorna un codice azione\} -- \{lancia \texttt{OpenJPAException}\} \\
\bottomrule
\end{tabularx}
\caption{Classi di equivalenza per \texttt{delete(Object pc, OpCallbacks call)}. * richiede
\texttt{call} non null.}
\label{tab:delete-classi}
\end{table}

\begin{table}[H]
\centering
\begin{tabularx}{\textwidth}{>{\raggedright\arraybackslash}p{0.38\textwidth} >{\raggedright\arraybackslash}X}
\toprule
Stato del broker & Classi di equivalenza \\
\midrule
transazione & \{non attiva\} -- \{attiva, senza istanze\} -- \{attiva, istanze presenti, successo\}*
-- \{attiva, istanze presenti, fallimento\}* \\
\bottomrule
\end{tabularx}
\caption{Classi di equivalenza per \texttt{flush()}. * non implementate.}
\label{tab:flush-classi}
\end{table}

\paragraph*{Comando PITest.}
\label{cmd:pitest-brokerimpl}
\begin{lstlisting}
mvn -pl openjpa-kernel \
  org.pitest:pitest-maven:1.5.1:mutationCoverage \
  -DtargetClasses=org.apache.openjpa.kernel.BrokerImpl \
  -DtargetTests=org.apache.openjpa.kernel.* \
  -DoutputFormats=CSV,HTML
\end{lstlisting}

\begin{table}[H]
\centering
\begin{tabularx}{\textwidth}{>{\raggedright\arraybackslash}p{0.38\textwidth} >{\raggedright\arraybackslash}X}
\toprule
Parametro & Classi di equivalenza \\
\midrule
\texttt{pc} & \{gestito e detached\} \\
\texttt{call} & \{null\} -- \{non null\} \\
\bottomrule
\end{tabularx}
\caption{Classi di equivalenza aggiuntive per \texttt{delete(Object pc, OpCallbacks call)},
categoria \texttt{pc} riusata da Tabella~\ref{tab:isdetached-classi}.}
\label{tab:delete2-classi}
\end{table}

\begin{table}[H]
\centering
\begin{tabularx}{\textwidth}{>{\raggedright\arraybackslash}p{0.38\textwidth} >{\raggedright\arraybackslash}X}
\toprule
Parametro & Classi di equivalenza \\
\midrule
\texttt{cls} & \{tipo non managed concreto\} \\
\bottomrule
\end{tabularx}
\caption{Classe di equivalenza aggiuntiva per \texttt{newInstance(Class cls)}, oracolo rafforzato
rispetto a Tabella~\ref{tab:newinstance-classi}.}
\label{tab:newinstance2-classi}
\end{table}

\paragraph*{Comando JaCoCo, \texttt{LoginDialog}.}
\label{cmd:jacoco-logindialog}
\begin{lstlisting}
AGENT_JAR=~/.m2/repository/org/jacoco/org.jacoco.agent/0.8.11/org.jacoco.agent-0.8.11-runtime.jar
export JAVA8_HOME=/usr/lib/jvm/java-8-openjdk-amd64

mvn -f openjpa-examples/opentrader/pom.xml \
  "-Dsurefire.jvm.args=-javaagent:${AGENT_JAR}=destfile=$(pwd)/openjpa-examples/opentrader/target/jacoco.exec -Xmx1024m" \
  clean test \
  org.jacoco:jacoco-maven-plugin:0.8.11:report
\end{lstlisting}

\paragraph*{Comando PITest, \texttt{LoginDialog}.}
\label{cmd:pitest-logindialog}
\begin{lstlisting}
export JAVA8_HOME=/usr/lib/jvm/java-8-openjdk-amd64

mvn -f openjpa-examples/opentrader/pom.xml \
  org.pitest:pitest-maven:1.5.1:mutationCoverage \
  -DtargetClasses=org.apache.openjpa.trader.client.LoginDialog \
  -DtargetTests=org.apache.openjpa.trader.client.* \
  -DoutputFormats=CSV,HTML
\end{lstlisting}
Esito:
\begin{lstlisting}
Execution default-cli of goal org.pitest:pitest-maven:1.5.1:mutationCoverage failed: 102 tests
did not pass without mutation when calculating line coverage. Mutation testing requires a green
suite.
\end{lstlisting}

\paragraph*{Prompt C1 (BrokerImpl, refactoring).}
\label{prompt:refactor-c1}
\begin{lstlisting}
You are an expert Java developer. I want to improve the maintainability of the attached
BrokerImpl class, a core class of the Apache OpenJPA persistence framework.

Create a refactored version without changing BrokerImpl's functionality, improving
maintainability by reducing common Java code smells (e.g. long methods, high cyclomatic
complexity, duplicated/near-duplicate code, unclear naming, deep nesting).

No tests are provided for this variant.

Do not include changes different from what I asked. The refactored class should replace
BrokerImpl and work with other components exactly as BrokerImpl currently does: same package
(org.apache.openjpa.kernel), same class name, same public and protected API other classes in
this project depend on.

This is an important request; take all the time you need for a complete and accurate answer.

The source code is attached.
\end{lstlisting}
Risposta completa in \texttt{report/data/refactoring-brokerimpl/c1/risposta.txt}; classe
completa in \texttt{report/data/refactoring-brokerimpl/c1/BrokerImpl.java}.

\paragraph*{Prompt C2 (BrokerImpl, refactoring).}
\label{prompt:refactor-c2}
\begin{lstlisting}
You are an expert Java developer. I want to improve the maintainability of the attached
BrokerImpl class, a core class of the Apache OpenJPA persistence framework.

Create a refactored version without changing BrokerImpl's functionality, improving
maintainability by reducing common Java code smells (e.g. long methods, high cyclomatic
complexity, duplicated/near-duplicate code, unclear naming, deep nesting).

Make sure the refactored class keeps passing the tests in the following files, attached:
FindTest, LockTest, DetachAllTest, AttachAllTest, NewInstanceTest, IsDetachedTest, PersistTest.

Do not include changes different from what I asked. The refactored class should replace
BrokerImpl and work with other components exactly as BrokerImpl currently does: same package
(org.apache.openjpa.kernel), same class name, same public and protected API other classes in
this project depend on.

This is an important request; take all the time you need for a complete and accurate answer.

The source code and the test files are attached.
\end{lstlisting}
Risposta completa in \texttt{report/data/refactoring-brokerimpl/c2/risposta.txt}; classe
completa in \texttt{report/data/refactoring-brokerimpl/c2/BrokerImpl.java}.

\paragraph*{Prompt C3 (BrokerImpl, refactoring).}
\label{prompt:refactor-c3}
\begin{lstlisting}
You are an expert Java developer. I want to improve the maintainability of the attached
BrokerImpl class, a core class of the Apache OpenJPA persistence framework.

Create a refactored version without changing BrokerImpl's functionality, improving
maintainability by reducing common Java code smells (e.g. long methods, high cyclomatic
complexity, duplicated/near-duplicate code, unclear naming, deep nesting).

Make sure the refactored class keeps passing the tests in the following files, attached:
FindTest, LockTest, DetachAllTest, AttachAllTest, NewInstanceTest, IsDetachedTest, PersistTest,
Lock2Test, DetachAll2Test, DeleteTest, FlushTest.

Do not include changes different from what I asked. The refactored class should replace
BrokerImpl and work with other components exactly as BrokerImpl currently does: same package
(org.apache.openjpa.kernel), same class name, same public and protected API other classes in
this project depend on.

This is an important request; take all the time you need for a complete and accurate answer.

The source code and the test files are attached.
\end{lstlisting}
Risposta completa in \texttt{report/data/refactoring-brokerimpl/c3/risposta.txt}; classe
completa in \texttt{report/data/refactoring-brokerimpl/c3/BrokerImpl.java}.

\paragraph*{Prompt C4 (BrokerImpl, refactoring).}
\label{prompt:refactor-c4}
\begin{lstlisting}
You are an expert Java developer. I want to improve the maintainability of the attached
BrokerImpl class, a core class of the Apache OpenJPA persistence framework.

Create a refactored version without changing BrokerImpl's functionality, improving
maintainability by reducing common Java code smells (e.g. long methods, high cyclomatic
complexity, duplicated/near-duplicate code, unclear naming, deep nesting).

Make sure the refactored class keeps passing the tests in the following files, attached:
FindTest, LockTest, DetachAllTest, AttachAllTest, NewInstanceTest, IsDetachedTest, PersistTest,
Lock2Test, DetachAll2Test, DeleteTest, FlushTest, Delete2Test, NewInstance2Test.

Do not include changes different from what I asked. The refactored class should replace
BrokerImpl and work with other components exactly as BrokerImpl currently does: same package
(org.apache.openjpa.kernel), same class name, same public and protected API other classes in
this project depend on.

This is an important request; take all the time you need for a complete and accurate answer.

The source code and the test files are attached.
\end{lstlisting}
Risposta completa in \texttt{report/data/refactoring-brokerimpl/c4/risposta.txt}; classe
completa in \texttt{report/data/refactoring-brokerimpl/c4/BrokerImpl.java}.

\paragraph*{Prompt C1 (LoginDialog, refactoring).}
\label{prompt:refactor-ld-c1}
\begin{lstlisting}
You are an expert Java developer. I want to improve the maintainability of the attached
LoginDialog class, a GWT dialog widget of the Apache OpenJPA opentrader example module.

Create a refactored version without changing LoginDialog's functionality, improving
maintainability by reducing common Java code smells (e.g. long methods, high cyclomatic
complexity, duplicated/near-duplicate code, unclear naming, deep nesting).

No tests are provided for this variant.

Do not include changes different from what I asked. The refactored class should replace
LoginDialog and work with other components exactly as LoginDialog currently does: same package
(org.apache.openjpa.trader.client), same class name, same public and protected API other
components in this project depend on (including the constructor signature and the public
GetServerURI and InitializeStocks inner classes).

This is an important request; take all the time you need for a complete and accurate answer.

The source code is attached.
\end{lstlisting}
Risposta completa in \texttt{report/data/refactoring-logindialog/c1/risposta.txt}; classe
completa in \texttt{report/data/refactoring-logindialog/c1/LoginDialog.java}.

\paragraph*{Prompt C2 (LoginDialog, refactoring).}
\label{prompt:refactor-ld-c2}
\begin{lstlisting}
You are an expert Java developer. I want to improve the maintainability of the attached
LoginDialog class, a GWT dialog widget of the Apache OpenJPA opentrader example module.

Create a refactored version without changing LoginDialog's functionality, improving
maintainability by reducing common Java code smells (e.g. long methods, high cyclomatic
complexity, duplicated/near-duplicate code, unclear naming, deep nesting).

Make sure the refactored class keeps passing the tests in the following files, attached:
CostruttoreTest, PosizioneTest, TitoloTest, ProprietaBooleaneTest, CicloDiVitaTest,
GetServerURITest, InitializeStocksTest.

Do not include changes different from what I asked. The refactored class should replace
LoginDialog and work with other components exactly as LoginDialog currently does: same package
(org.apache.openjpa.trader.client), same class name, same public and protected API other
components in this project depend on (including the constructor signature and the public
GetServerURI and InitializeStocks inner classes).

This is an important request; take all the time you need for a complete and accurate answer.

The source code and the test files are attached.
\end{lstlisting}
Risposta completa in \texttt{report/data/refactoring-logindialog/c2/risposta.txt}; classe
completa in \texttt{report/data/refactoring-logindialog/c2/LoginDialog.java}.

\paragraph*{Tuple di progettazione, nota.}
Le tabelle seguenti riportano, per ciascun metodo, tutte le tuple di Category Partition usate come test: parametri e oracolo atteso. Le colonne di classificazione della fonte dell'oracolo e le note di progettazione sono discusse nel corpo del report (Sezioni \ref{sec:cp} e \ref{sec:implementazione}) e omesse qui per compattezza. Le tabelle ``corrette'' riportano l'oracolo dopo la correzione empirica descritta in Sezione \ref{sec:implementazione}, con le stesse convenzioni di riga del file originale (righe aggiunte o rimosse rispetto alla progettazione sono indicate nel testo del capitolo).

\paragraph*{Tuple di progettazione, \texttt{find(Object oid, boolean validate, FindCallbacks call)} (BrokerImpl).}

{\footnotesize\setlength{\tabcolsep}{3pt}
\begin{longtable}{>{\raggedright\arraybackslash}p{0.071\textwidth}>{\raggedright\arraybackslash}p{0.056\textwidth}>{\raggedright\arraybackslash}p{0.085\textwidth}>{\raggedright\arraybackslash}p{0.085\textwidth}>{\raggedright\arraybackslash}p{0.056\textwidth}>{\raggedright\arraybackslash}p{0.106\textwidth}>{\raggedright\arraybackslash}p{0.106\textwidth}>{\raggedright\arraybackslash}p{0.297\textwidth}}
\toprule
\textbf{oid} & \textbf{validate} & \textbf{cache} & \textbf{store} & \textbf{call} & \textbf{processArgument} & \textbf{processReturn} & \textbf{oracolo} \\
\midrule
\endfirsthead
\toprule
\textbf{oid} & \textbf{validate} & \textbf{cache} & \textbf{store} & \textbf{call} & \textbf{processArgument} & \textbf{processReturn} & \textbf{oracolo} \\
\midrule
\endhead
\bottomrule
\endfoot
null & true & -- & -- & null & -- & -- & null \\
null & false & -- & -- & null & -- & -- & eccezione (tipo non specificato) \\
non valido & true & -- & -- & null & -- & -- & null \\
non valido & false & -- & -- & null & -- & -- & eccezione (tipo non specificato) \\
valido & true & presente & presente & null & -- & -- & istanza non-null \\
valido & true & non presente & presente & null & -- & -- & istanza non-null \\
valido & true & presente & non presente & null & -- & -- & null \\
valido & true & non presente & non presente & null & -- & -- & null \\
valido & false & presente & presente & null & -- & -- & istanza in cache \\
valido & false & non presente & presente & null & -- & -- & istanza hollow non-null \\
valido & false & presente & non presente & null & -- & -- & istanza in cache \\
valido & false & non presente & non presente & null & -- & -- & istanza hollow non-null \\
null & true & presente & presente & non null & ritorna id & ritorna oggetto & oggetto ritornato da processReturn (non-null, non ulteriormente specificato) \\
null & true & presente & presente & non null & ritorna id & ritorna null & null \\
null & true & non presente & presente & non null & ritorna id & ritorna oggetto & oggetto ritornato da processReturn (non-null, non ulteriormente specificato) \\
null & true & non presente & presente & non null & ritorna id & ritorna null & null \\
null & true & presente & non presente & non null & ritorna id & ritorna oggetto & oggetto ritornato da processReturn (non-null, non ulteriormente specificato) \\
null & true & presente & non presente & non null & ritorna id & ritorna null & null \\
null & true & non presente & non presente & non null & ritorna id & ritorna oggetto & oggetto ritornato da processReturn (non-null, non ulteriormente specificato) \\
null & true & non presente & non presente & non null & ritorna id & ritorna null & null \\
null & true & -- & -- & non null & ritorna null & ritorna oggetto & oggetto ritornato da processReturn (non-null, non ulteriormente specificato) \\
null & true & -- & -- & non null & ritorna null & ritorna null & null \\
null & true & -- & -- & non null & eccezione & ritorna oggetto & eccezione propagata da find \\
null & true & -- & -- & non null & eccezione & ritorna null & eccezione propagata da find \\
null & false & presente & presente & non null & ritorna id & ritorna oggetto & oggetto ritornato da processReturn (non-null, non ulteriormente specificato) \\
null & false & presente & presente & non null & ritorna id & ritorna null & istanza non-null (garanzia 'mai null' assunta prevalente sul null di processReturn) \\
null & false & non presente & presente & non null & ritorna id & ritorna oggetto & oggetto ritornato da processReturn (non-null, non ulteriormente specificato) \\
null & false & non presente & presente & non null & ritorna id & ritorna null & istanza non-null (garanzia 'mai null' assunta prevalente sul null di processReturn) \\
null & false & presente & non presente & non null & ritorna id & ritorna oggetto & oggetto ritornato da processReturn (non-null, non ulteriormente specificato) \\
null & false & presente & non presente & non null & ritorna id & ritorna null & istanza non-null (garanzia 'mai null' assunta prevalente sul null di processReturn) \\
null & false & non presente & non presente & non null & ritorna id & ritorna oggetto & oggetto ritornato da processReturn (non-null, non ulteriormente specificato) \\
null & false & non presente & non presente & non null & ritorna id & ritorna null & istanza non-null (garanzia 'mai null' assunta prevalente sul null di processReturn) \\
null & false & -- & -- & non null & ritorna null & ritorna oggetto & oggetto ritornato da processReturn (non-null, non ulteriormente specificato) \\
null & false & -- & -- & non null & ritorna null & ritorna null & istanza non-null (garanzia 'mai null' assunta prevalente sul null di processReturn) \\
null & false & -- & -- & non null & eccezione & ritorna oggetto & eccezione propagata da find \\
null & false & -- & -- & non null & eccezione & ritorna null & eccezione propagata da find \\
non valido & true & presente & presente & non null & ritorna id & ritorna oggetto & oggetto ritornato da processReturn (non-null, non ulteriormente specificato) \\
non valido & true & presente & presente & non null & ritorna id & ritorna null & null \\
non valido & true & non presente & presente & non null & ritorna id & ritorna oggetto & oggetto ritornato da processReturn (non-null, non ulteriormente specificato) \\
non valido & true & non presente & presente & non null & ritorna id & ritorna null & null \\
non valido & true & presente & non presente & non null & ritorna id & ritorna oggetto & oggetto ritornato da processReturn (non-null, non ulteriormente specificato) \\
non valido & true & presente & non presente & non null & ritorna id & ritorna null & null \\
non valido & true & non presente & non presente & non null & ritorna id & ritorna oggetto & oggetto ritornato da processReturn (non-null, non ulteriormente specificato) \\
non valido & true & non presente & non presente & non null & ritorna id & ritorna null & null \\
non valido & true & -- & -- & non null & ritorna null & ritorna oggetto & oggetto ritornato da processReturn (non-null, non ulteriormente specificato) \\
non valido & true & -- & -- & non null & ritorna null & ritorna null & null \\
non valido & true & -- & -- & non null & eccezione & ritorna oggetto & eccezione propagata da find \\
non valido & true & -- & -- & non null & eccezione & ritorna null & eccezione propagata da find \\
non valido & false & presente & presente & non null & ritorna id & ritorna oggetto & oggetto ritornato da processReturn (non-null, non ulteriormente specificato) \\
non valido & false & presente & presente & non null & ritorna id & ritorna null & istanza non-null (garanzia 'mai null' assunta prevalente sul null di processReturn) \\
non valido & false & non presente & presente & non null & ritorna id & ritorna oggetto & oggetto ritornato da processReturn (non-null, non ulteriormente specificato) \\
non valido & false & non presente & presente & non null & ritorna id & ritorna null & istanza non-null (garanzia 'mai null' assunta prevalente sul null di processReturn) \\
non valido & false & presente & non presente & non null & ritorna id & ritorna oggetto & oggetto ritornato da processReturn (non-null, non ulteriormente specificato) \\
non valido & false & presente & non presente & non null & ritorna id & ritorna null & istanza non-null (garanzia 'mai null' assunta prevalente sul null di processReturn) \\
non valido & false & non presente & non presente & non null & ritorna id & ritorna oggetto & oggetto ritornato da processReturn (non-null, non ulteriormente specificato) \\
non valido & false & non presente & non presente & non null & ritorna id & ritorna null & istanza non-null (garanzia 'mai null' assunta prevalente sul null di processReturn) \\
non valido & false & -- & -- & non null & ritorna null & ritorna oggetto & oggetto ritornato da processReturn (non-null, non ulteriormente specificato) \\
non valido & false & -- & -- & non null & ritorna null & ritorna null & istanza non-null (garanzia 'mai null' assunta prevalente sul null di processReturn) \\
non valido & false & -- & -- & non null & eccezione & ritorna oggetto & eccezione propagata da find \\
non valido & false & -- & -- & non null & eccezione & ritorna null & eccezione propagata da find \\
valido & true & presente & presente & non null & ritorna id & ritorna oggetto & oggetto ritornato da processReturn (non-null, non ulteriormente specificato) \\
valido & true & presente & presente & non null & ritorna id & ritorna null & null \\
valido & true & non presente & presente & non null & ritorna id & ritorna oggetto & oggetto ritornato da processReturn (non-null, non ulteriormente specificato) \\
valido & true & non presente & presente & non null & ritorna id & ritorna null & null \\
valido & true & presente & non presente & non null & ritorna id & ritorna oggetto & oggetto ritornato da processReturn (non-null, non ulteriormente specificato) \\
valido & true & presente & non presente & non null & ritorna id & ritorna null & null \\
valido & true & non presente & non presente & non null & ritorna id & ritorna oggetto & oggetto ritornato da processReturn (non-null, non ulteriormente specificato) \\
valido & true & non presente & non presente & non null & ritorna id & ritorna null & null \\
valido & true & -- & -- & non null & ritorna null & ritorna oggetto & oggetto ritornato da processReturn (non-null, non ulteriormente specificato) \\
valido & true & -- & -- & non null & ritorna null & ritorna null & null \\
valido & true & -- & -- & non null & eccezione & ritorna oggetto & eccezione propagata da find \\
valido & true & -- & -- & non null & eccezione & ritorna null & eccezione propagata da find \\
valido & false & presente & presente & non null & ritorna id & ritorna oggetto & oggetto ritornato da processReturn (non-null, non ulteriormente specificato) \\
valido & false & presente & presente & non null & ritorna id & ritorna null & istanza non-null (garanzia 'mai null' assunta prevalente sul null di processReturn) \\
valido & false & non presente & presente & non null & ritorna id & ritorna oggetto & oggetto ritornato da processReturn (non-null, non ulteriormente specificato) \\
valido & false & non presente & presente & non null & ritorna id & ritorna null & istanza non-null (garanzia 'mai null' assunta prevalente sul null di processReturn) \\
valido & false & presente & non presente & non null & ritorna id & ritorna oggetto & oggetto ritornato da processReturn (non-null, non ulteriormente specificato) \\
valido & false & presente & non presente & non null & ritorna id & ritorna null & istanza non-null (garanzia 'mai null' assunta prevalente sul null di processReturn) \\
valido & false & non presente & non presente & non null & ritorna id & ritorna oggetto & oggetto ritornato da processReturn (non-null, non ulteriormente specificato) \\
valido & false & non presente & non presente & non null & ritorna id & ritorna null & istanza non-null (garanzia 'mai null' assunta prevalente sul null di processReturn) \\
valido & false & -- & -- & non null & ritorna null & ritorna oggetto & oggetto ritornato da processReturn (non-null, non ulteriormente specificato) \\
valido & false & -- & -- & non null & ritorna null & ritorna null & istanza non-null (garanzia 'mai null' assunta prevalente sul null di processReturn) \\
valido & false & -- & -- & non null & eccezione & ritorna oggetto & eccezione propagata da find \\
valido & false & -- & -- & non null & eccezione & ritorna null & eccezione propagata da find \\
\bottomrule
\caption{Tuple di progettazione per \texttt{find(Object oid, boolean validate, FindCallbacks call)}.}
\label{tab:find-tuple}
\end{longtable}
}



\paragraph*{Tuple corrette, \texttt{find(Object oid, boolean validate, FindCallbacks call)} (BrokerImpl).}

{\footnotesize\setlength{\tabcolsep}{3pt}
\begin{longtable}{>{\raggedright\arraybackslash}p{0.071\textwidth}>{\raggedright\arraybackslash}p{0.056\textwidth}>{\raggedright\arraybackslash}p{0.085\textwidth}>{\raggedright\arraybackslash}p{0.085\textwidth}>{\raggedright\arraybackslash}p{0.056\textwidth}>{\raggedright\arraybackslash}p{0.106\textwidth}>{\raggedright\arraybackslash}p{0.106\textwidth}>{\raggedright\arraybackslash}p{0.297\textwidth}}
\toprule
\textbf{oid} & \textbf{validate} & \textbf{cache} & \textbf{store} & \textbf{call} & \textbf{processArgument} & \textbf{processReturn} & \textbf{oracolo} \\
\midrule
\endfirsthead
\toprule
\textbf{oid} & \textbf{validate} & \textbf{cache} & \textbf{store} & \textbf{call} & \textbf{processArgument} & \textbf{processReturn} & \textbf{oracolo} \\
\midrule
\endhead
\bottomrule
\endfoot
null & true & -- & -- & null & -- & -- & eccezione (ObjectNotFoundException) \\
null & false & -- & -- & null & -- & -- & null \\
non valido & true & -- & -- & null & -- & -- & eccezione (MetaDataException) \\
non valido & false & -- & -- & null & -- & -- & eccezione (tipo non specificato) \\
valido & true & presente & presente & null & -- & -- & istanza non-null \\
valido & true & non presente & presente & null & -- & -- & istanza non-null \\
valido & true & presente & non presente & null & -- & -- & null \\
valido & true & non presente & non presente & null & -- & -- & null \\
valido & false & presente & presente & null & -- & -- & istanza in cache \\
valido & false & non presente & presente & null & -- & -- & istanza hollow non-null \\
valido & false & presente & non presente & null & -- & -- & istanza in cache \\
valido & false & non presente & non presente & null & -- & -- & istanza hollow non-null \\
null & true & presente & presente & non null & ritorna id & ritorna oggetto & oggetto ritornato da processReturn (non-null, non ulteriormente specificato) \\
null & true & presente & presente & non null & ritorna id & ritorna null & null \\
null & true & non presente & presente & non null & ritorna id & ritorna oggetto & oggetto ritornato da processReturn (non-null, non ulteriormente specificato) \\
null & true & non presente & presente & non null & ritorna id & ritorna null & null \\
null & true & presente & non presente & non null & ritorna id & ritorna oggetto & oggetto ritornato da processReturn (non-null, non ulteriormente specificato) \\
null & true & presente & non presente & non null & ritorna id & ritorna null & null \\
null & true & non presente & non presente & non null & ritorna id & ritorna oggetto & oggetto ritornato da processReturn (non-null, non ulteriormente specificato) \\
null & true & non presente & non presente & non null & ritorna id & ritorna null & null \\
null & true & -- & -- & non null & ritorna null & ritorna oggetto & eccezione (ObjectNotFoundException) \\
null & true & -- & -- & non null & ritorna null & ritorna null & eccezione (ObjectNotFoundException) \\
null & true & -- & -- & non null & eccezione & ritorna oggetto & eccezione propagata da find \\
null & true & -- & -- & non null & eccezione & ritorna null & eccezione propagata da find \\
null & false & presente & presente & non null & ritorna id & ritorna oggetto & oggetto ritornato da processReturn (non-null, non ulteriormente specificato) \\
null & false & presente & presente & non null & ritorna id & ritorna null & null \\
null & false & non presente & presente & non null & ritorna id & ritorna oggetto & oggetto ritornato da processReturn (non-null, non ulteriormente specificato) \\
null & false & non presente & presente & non null & ritorna id & ritorna null & null \\
null & false & presente & non presente & non null & ritorna id & ritorna oggetto & oggetto ritornato da processReturn (non-null, non ulteriormente specificato) \\
null & false & presente & non presente & non null & ritorna id & ritorna null & null \\
null & false & non presente & non presente & non null & ritorna id & ritorna oggetto & oggetto ritornato da processReturn (non-null, non ulteriormente specificato) \\
null & false & non presente & non presente & non null & ritorna id & ritorna null & null \\
null & false & -- & -- & non null & ritorna null & ritorna oggetto & oggetto ritornato da processReturn (non-null, non ulteriormente specificato) \\
null & false & -- & -- & non null & ritorna null & ritorna null & null \\
null & false & -- & -- & non null & eccezione & ritorna oggetto & eccezione propagata da find \\
null & false & -- & -- & non null & eccezione & ritorna null & eccezione propagata da find \\
non valido & true & presente & presente & non null & ritorna id & ritorna oggetto & oggetto ritornato da processReturn (non-null, non ulteriormente specificato) \\
non valido & true & presente & presente & non null & ritorna id & ritorna null & null \\
non valido & true & non presente & presente & non null & ritorna id & ritorna oggetto & oggetto ritornato da processReturn (non-null, non ulteriormente specificato) \\
non valido & true & non presente & presente & non null & ritorna id & ritorna null & null \\
non valido & true & presente & non presente & non null & ritorna id & ritorna oggetto & oggetto ritornato da processReturn (non-null, non ulteriormente specificato) \\
non valido & true & presente & non presente & non null & ritorna id & ritorna null & null \\
non valido & true & non presente & non presente & non null & ritorna id & ritorna oggetto & oggetto ritornato da processReturn (non-null, non ulteriormente specificato) \\
non valido & true & non presente & non presente & non null & ritorna id & ritorna null & null \\
non valido & true & -- & -- & non null & ritorna null & ritorna oggetto & eccezione (ObjectNotFoundException) \\
non valido & true & -- & -- & non null & ritorna null & ritorna null & eccezione (ObjectNotFoundException) \\
non valido & true & -- & -- & non null & eccezione & ritorna oggetto & eccezione propagata da find \\
non valido & true & -- & -- & non null & eccezione & ritorna null & eccezione propagata da find \\
non valido & false & presente & presente & non null & ritorna id & ritorna oggetto & oggetto ritornato da processReturn (non-null, non ulteriormente specificato) \\
non valido & false & presente & presente & non null & ritorna id & ritorna null & null \\
non valido & false & non presente & presente & non null & ritorna id & ritorna oggetto & oggetto ritornato da processReturn (non-null, non ulteriormente specificato) \\
non valido & false & non presente & presente & non null & ritorna id & ritorna null & null \\
non valido & false & presente & non presente & non null & ritorna id & ritorna oggetto & oggetto ritornato da processReturn (non-null, non ulteriormente specificato) \\
non valido & false & presente & non presente & non null & ritorna id & ritorna null & null \\
non valido & false & non presente & non presente & non null & ritorna id & ritorna oggetto & oggetto ritornato da processReturn (non-null, non ulteriormente specificato) \\
non valido & false & non presente & non presente & non null & ritorna id & ritorna null & null \\
non valido & false & -- & -- & non null & ritorna null & ritorna oggetto & oggetto ritornato da processReturn (non-null, non ulteriormente specificato) \\
non valido & false & -- & -- & non null & ritorna null & ritorna null & null \\
non valido & false & -- & -- & non null & eccezione & ritorna oggetto & eccezione propagata da find \\
non valido & false & -- & -- & non null & eccezione & ritorna null & eccezione propagata da find \\
valido & true & presente & presente & non null & ritorna id & ritorna oggetto & oggetto ritornato da processReturn (non-null, non ulteriormente specificato) \\
valido & true & presente & presente & non null & ritorna id & ritorna null & null \\
valido & true & non presente & presente & non null & ritorna id & ritorna oggetto & oggetto ritornato da processReturn (non-null, non ulteriormente specificato) \\
valido & true & non presente & presente & non null & ritorna id & ritorna null & null \\
valido & true & presente & non presente & non null & ritorna id & ritorna oggetto & oggetto ritornato da processReturn (non-null, non ulteriormente specificato) \\
valido & true & presente & non presente & non null & ritorna id & ritorna null & null \\
valido & true & non presente & non presente & non null & ritorna id & ritorna oggetto & oggetto ritornato da processReturn (non-null, non ulteriormente specificato) \\
valido & true & non presente & non presente & non null & ritorna id & ritorna null & null \\
valido & true & -- & -- & non null & ritorna null & ritorna oggetto & eccezione (ObjectNotFoundException) \\
valido & true & -- & -- & non null & ritorna null & ritorna null & eccezione (ObjectNotFoundException) \\
valido & true & -- & -- & non null & eccezione & ritorna oggetto & eccezione propagata da find \\
valido & true & -- & -- & non null & eccezione & ritorna null & eccezione propagata da find \\
valido & false & presente & presente & non null & ritorna id & ritorna oggetto & oggetto ritornato da processReturn (non-null, non ulteriormente specificato) \\
valido & false & presente & presente & non null & ritorna id & ritorna null & null \\
valido & false & non presente & presente & non null & ritorna id & ritorna oggetto & oggetto ritornato da processReturn (non-null, non ulteriormente specificato) \\
valido & false & non presente & presente & non null & ritorna id & ritorna null & null \\
valido & false & presente & non presente & non null & ritorna id & ritorna oggetto & oggetto ritornato da processReturn (non-null, non ulteriormente specificato) \\
valido & false & presente & non presente & non null & ritorna id & ritorna null & null \\
valido & false & non presente & non presente & non null & ritorna id & ritorna oggetto & oggetto ritornato da processReturn (non-null, non ulteriormente specificato) \\
valido & false & non presente & non presente & non null & ritorna id & ritorna null & null \\
valido & false & -- & -- & non null & ritorna null & ritorna oggetto & oggetto ritornato da processReturn (non-null, non ulteriormente specificato) \\
valido & false & -- & -- & non null & ritorna null & ritorna null & null \\
valido & false & -- & -- & non null & eccezione & ritorna oggetto & eccezione propagata da find \\
valido & false & -- & -- & non null & eccezione & ritorna null & eccezione propagata da find \\
\bottomrule
\caption{Tuple corrette per \texttt{find(Object oid, boolean validate, FindCallbacks call)}.}
\label{tab:find-tuple-corrette}
\end{longtable}
}



\paragraph*{Tuple di progettazione, \texttt{lock(Object pc, int level, int timeout, OpCallbacks call)} (BrokerImpl).}

{\footnotesize\setlength{\tabcolsep}{3pt}
\begin{longtable}{>{\raggedright\arraybackslash}p{0.179\textwidth}>{\raggedright\arraybackslash}p{0.065\textwidth}>{\raggedright\arraybackslash}p{0.114\textwidth}>{\raggedright\arraybackslash}p{0.065\textwidth}>{\raggedright\arraybackslash}p{0.122\textwidth}>{\raggedright\arraybackslash}p{0.342\textwidth}}
\toprule
\textbf{pc} & \textbf{level} & \textbf{timeout} & \textbf{call} & \textbf{processArgument} & \textbf{oracolo} \\
\midrule
\endfirsthead
\toprule
\textbf{pc} & \textbf{level} & \textbf{timeout} & \textbf{call} & \textbf{processArgument} & \textbf{oracolo} \\
\midrule
\endhead
\bottomrule
\endfoot
null & noto & altro negativo & null & -- & eccezione (tipo non specificato) \\
null & noto & altro negativo & non null & azione & eccezione (tipo non specificato) \\
null & noto & altro negativo & non null & eccezione & eccezione propagata da lock \\
null & noto & -1 & null & -- & eccezione (tipo non specificato) \\
null & noto & -1 & non null & azione & eccezione (tipo non specificato) \\
null & noto & -1 & non null & eccezione & eccezione propagata da lock \\
null & noto & 0 & null & -- & eccezione (tipo non specificato) \\
null & noto & 0 & non null & azione & eccezione (tipo non specificato) \\
null & noto & 0 & non null & eccezione & eccezione propagata da lock \\
null & noto & positivo & null & -- & eccezione (tipo non specificato) \\
null & noto & positivo & non null & azione & eccezione (tipo non specificato) \\
null & noto & positivo & non null & eccezione & eccezione propagata da lock \\
null & non noto & altro negativo & null & -- & eccezione (tipo non specificato) \\
null & non noto & altro negativo & non null & azione & eccezione (tipo non specificato) \\
null & non noto & altro negativo & non null & eccezione & eccezione propagata da lock \\
null & non noto & -1 & null & -- & eccezione (tipo non specificato) \\
null & non noto & -1 & non null & azione & eccezione (tipo non specificato) \\
null & non noto & -1 & non null & eccezione & eccezione propagata da lock \\
null & non noto & 0 & null & -- & eccezione (tipo non specificato) \\
null & non noto & 0 & non null & azione & eccezione (tipo non specificato) \\
null & non noto & 0 & non null & eccezione & eccezione propagata da lock \\
null & non noto & positivo & null & -- & eccezione (tipo non specificato) \\
null & non noto & positivo & non null & azione & eccezione (tipo non specificato) \\
null & non noto & positivo & non null & eccezione & eccezione propagata da lock \\
non gestito & noto & altro negativo & null & -- & eccezione (tipo non specificato) \\
non gestito & noto & altro negativo & non null & azione & eccezione (tipo non specificato) \\
non gestito & noto & altro negativo & non null & eccezione & eccezione propagata da lock \\
non gestito & noto & -1 & null & -- & eccezione (tipo non specificato) \\
non gestito & noto & -1 & non null & azione & eccezione (tipo non specificato) \\
non gestito & noto & -1 & non null & eccezione & eccezione propagata da lock \\
non gestito & noto & 0 & null & -- & eccezione (tipo non specificato) \\
non gestito & noto & 0 & non null & azione & eccezione (tipo non specificato) \\
non gestito & noto & 0 & non null & eccezione & eccezione propagata da lock \\
non gestito & noto & positivo & null & -- & eccezione (tipo non specificato) \\
non gestito & noto & positivo & non null & azione & eccezione (tipo non specificato) \\
non gestito & noto & positivo & non null & eccezione & eccezione propagata da lock \\
non gestito & non noto & altro negativo & null & -- & eccezione (tipo non specificato) \\
non gestito & non noto & altro negativo & non null & azione & eccezione (tipo non specificato) \\
non gestito & non noto & altro negativo & non null & eccezione & eccezione propagata da lock \\
non gestito & non noto & -1 & null & -- & eccezione (tipo non specificato) \\
non gestito & non noto & -1 & non null & azione & eccezione (tipo non specificato) \\
non gestito & non noto & -1 & non null & eccezione & eccezione propagata da lock \\
non gestito & non noto & 0 & null & -- & eccezione (tipo non specificato) \\
non gestito & non noto & 0 & non null & azione & eccezione (tipo non specificato) \\
non gestito & non noto & 0 & non null & eccezione & eccezione propagata da lock \\
non gestito & non noto & positivo & null & -- & eccezione (tipo non specificato) \\
non gestito & non noto & positivo & non null & azione & eccezione (tipo non specificato) \\
non gestito & non noto & positivo & non null & eccezione & eccezione propagata da lock \\
gestito non persistito & noto & altro negativo & null & -- & eccezione (tipo non specificato) \\
gestito non persistito & noto & altro negativo & non null & azione & eccezione (tipo non specificato) \\
gestito non persistito & noto & altro negativo & non null & eccezione & eccezione propagata da lock \\
gestito non persistito & noto & -1 & null & -- & getLockLevel(pc) == level (postcondizione) \\
gestito non persistito & noto & -1 & non null & azione & getLockLevel(pc) == level (postcondizione) \\
gestito non persistito & noto & -1 & non null & eccezione & eccezione propagata da lock \\
gestito non persistito & noto & 0 & null & -- & getLockLevel(pc) == level (postcondizione) \\
gestito non persistito & noto & 0 & non null & azione & getLockLevel(pc) == level (postcondizione) \\
gestito non persistito & noto & 0 & non null & eccezione & eccezione propagata da lock \\
gestito non persistito & noto & positivo & null & -- & getLockLevel(pc) == level (postcondizione) \\
gestito non persistito & noto & positivo & non null & azione & getLockLevel(pc) == level (postcondizione) \\
gestito non persistito & noto & positivo & non null & eccezione & eccezione propagata da lock \\
gestito non persistito & non noto & altro negativo & null & -- & eccezione (tipo non specificato) \\
gestito non persistito & non noto & altro negativo & non null & azione & eccezione (tipo non specificato) \\
gestito non persistito & non noto & altro negativo & non null & eccezione & eccezione propagata da lock \\
gestito non persistito & non noto & -1 & null & -- & eccezione (tipo non specificato) \\
gestito non persistito & non noto & -1 & non null & azione & eccezione (tipo non specificato) \\
gestito non persistito & non noto & -1 & non null & eccezione & eccezione propagata da lock \\
gestito non persistito & non noto & 0 & null & -- & eccezione (tipo non specificato) \\
gestito non persistito & non noto & 0 & non null & azione & eccezione (tipo non specificato) \\
gestito non persistito & non noto & 0 & non null & eccezione & eccezione propagata da lock \\
gestito non persistito & non noto & positivo & null & -- & eccezione (tipo non specificato) \\
gestito non persistito & non noto & positivo & non null & azione & eccezione (tipo non specificato) \\
gestito non persistito & non noto & positivo & non null & eccezione & eccezione propagata da lock \\
gestito persistito & noto & altro negativo & null & -- & eccezione (tipo non specificato) \\
gestito persistito & noto & altro negativo & non null & azione & eccezione (tipo non specificato) \\
gestito persistito & noto & altro negativo & non null & eccezione & eccezione propagata da lock \\
gestito persistito & noto & -1 & null & -- & getLockLevel(pc) == level (postcondizione) \\
gestito persistito & noto & -1 & non null & azione & getLockLevel(pc) == level (postcondizione) \\
gestito persistito & noto & -1 & non null & eccezione & eccezione propagata da lock \\
gestito persistito & noto & 0 & null & -- & getLockLevel(pc) == level (postcondizione) \\
gestito persistito & noto & 0 & non null & azione & getLockLevel(pc) == level (postcondizione) \\
gestito persistito & noto & 0 & non null & eccezione & eccezione propagata da lock \\
gestito persistito & noto & positivo & null & -- & getLockLevel(pc) == level (postcondizione) \\
gestito persistito & noto & positivo & non null & azione & getLockLevel(pc) == level (postcondizione) \\
gestito persistito & noto & positivo & non null & eccezione & eccezione propagata da lock \\
gestito persistito & non noto & altro negativo & null & -- & eccezione (tipo non specificato) \\
gestito persistito & non noto & altro negativo & non null & azione & eccezione (tipo non specificato) \\
gestito persistito & non noto & altro negativo & non null & eccezione & eccezione propagata da lock \\
gestito persistito & non noto & -1 & null & -- & eccezione (tipo non specificato) \\
gestito persistito & non noto & -1 & non null & azione & eccezione (tipo non specificato) \\
gestito persistito & non noto & -1 & non null & eccezione & eccezione propagata da lock \\
gestito persistito & non noto & 0 & null & -- & eccezione (tipo non specificato) \\
gestito persistito & non noto & 0 & non null & azione & eccezione (tipo non specificato) \\
gestito persistito & non noto & 0 & non null & eccezione & eccezione propagata da lock \\
gestito persistito & non noto & positivo & null & -- & eccezione (tipo non specificato) \\
gestito persistito & non noto & positivo & non null & azione & eccezione (tipo non specificato) \\
gestito persistito & non noto & positivo & non null & eccezione & eccezione propagata da lock \\
\bottomrule
\caption{Tuple di progettazione per \texttt{lock(Object pc, int level, int timeout, OpCallbacks call)}.}
\label{tab:lock-tuple}
\end{longtable}
}



\paragraph*{Tuple corrette, \texttt{lock(Object pc, int level, int timeout, OpCallbacks call)} (BrokerImpl).}

{\footnotesize\setlength{\tabcolsep}{3pt}
\begin{longtable}{>{\raggedright\arraybackslash}p{0.171\textwidth}>{\raggedright\arraybackslash}p{0.062\textwidth}>{\raggedright\arraybackslash}p{0.109\textwidth}>{\raggedright\arraybackslash}p{0.062\textwidth}>{\raggedright\arraybackslash}p{0.156\textwidth}>{\raggedright\arraybackslash}p{0.327\textwidth}}
\toprule
\textbf{pc} & \textbf{level} & \textbf{timeout} & \textbf{call} & \textbf{processArgument} & \textbf{oracolo} \\
\midrule
\endfirsthead
\toprule
\textbf{pc} & \textbf{level} & \textbf{timeout} & \textbf{call} & \textbf{processArgument} & \textbf{oracolo} \\
\midrule
\endhead
\bottomrule
\endfoot
null & noto & altro negativo & null & -- & nessun effetto (nessuna eccezione, getLockLevel(pc) == LOCK\_NONE) \\
null & noto & altro negativo & non null & azione & nessun effetto (nessuna eccezione, getLockLevel(pc) == LOCK\_NONE) \\
null & noto & altro negativo & non null & eccezione & nessun effetto (nessuna eccezione, getLockLevel(pc) == LOCK\_NONE) \\
null & noto & -1 & null & -- & nessun effetto (nessuna eccezione, getLockLevel(pc) == LOCK\_NONE) \\
null & noto & -1 & non null & azione & nessun effetto (nessuna eccezione, getLockLevel(pc) == LOCK\_NONE) \\
null & noto & -1 & non null & eccezione & nessun effetto (nessuna eccezione, getLockLevel(pc) == LOCK\_NONE) \\
null & noto & 0 & null & -- & nessun effetto (nessuna eccezione, getLockLevel(pc) == LOCK\_NONE) \\
null & noto & 0 & non null & azione & nessun effetto (nessuna eccezione, getLockLevel(pc) == LOCK\_NONE) \\
null & noto & 0 & non null & eccezione & nessun effetto (nessuna eccezione, getLockLevel(pc) == LOCK\_NONE) \\
null & noto & positivo & null & -- & nessun effetto (nessuna eccezione, getLockLevel(pc) == LOCK\_NONE) \\
null & noto & positivo & non null & azione & nessun effetto (nessuna eccezione, getLockLevel(pc) == LOCK\_NONE) \\
null & noto & positivo & non null & eccezione & nessun effetto (nessuna eccezione, getLockLevel(pc) == LOCK\_NONE) \\
null & non noto & altro negativo & null & -- & nessun effetto (nessuna eccezione, getLockLevel(pc) == LOCK\_NONE) \\
null & non noto & altro negativo & non null & azione & nessun effetto (nessuna eccezione, getLockLevel(pc) == LOCK\_NONE) \\
null & non noto & altro negativo & non null & eccezione & nessun effetto (nessuna eccezione, getLockLevel(pc) == LOCK\_NONE) \\
null & non noto & -1 & null & -- & nessun effetto (nessuna eccezione, getLockLevel(pc) == LOCK\_NONE) \\
null & non noto & -1 & non null & azione & nessun effetto (nessuna eccezione, getLockLevel(pc) == LOCK\_NONE) \\
null & non noto & -1 & non null & eccezione & nessun effetto (nessuna eccezione, getLockLevel(pc) == LOCK\_NONE) \\
null & non noto & 0 & null & -- & nessun effetto (nessuna eccezione, getLockLevel(pc) == LOCK\_NONE) \\
null & non noto & 0 & non null & azione & nessun effetto (nessuna eccezione, getLockLevel(pc) == LOCK\_NONE) \\
null & non noto & 0 & non null & eccezione & nessun effetto (nessuna eccezione, getLockLevel(pc) == LOCK\_NONE) \\
null & non noto & positivo & null & -- & nessun effetto (nessuna eccezione, getLockLevel(pc) == LOCK\_NONE) \\
null & non noto & positivo & non null & azione & nessun effetto (nessuna eccezione, getLockLevel(pc) == LOCK\_NONE) \\
null & non noto & positivo & non null & eccezione & nessun effetto (nessuna eccezione, getLockLevel(pc) == LOCK\_NONE) \\
non gestito & noto & altro negativo & null & -- & nessun effetto (nessuna eccezione, getLockLevel(pc) == LOCK\_NONE) \\
non gestito & noto & altro negativo & non null & azione & nessun effetto (nessuna eccezione, getLockLevel(pc) == LOCK\_NONE) \\
non gestito & noto & altro negativo & non null & eccezione & eccezione propagata da lock \\
non gestito & noto & -1 & null & -- & nessun effetto (nessuna eccezione, getLockLevel(pc) == LOCK\_NONE) \\
non gestito & noto & -1 & non null & azione & nessun effetto (nessuna eccezione, getLockLevel(pc) == LOCK\_NONE) \\
non gestito & noto & -1 & non null & eccezione & eccezione propagata da lock \\
non gestito & noto & 0 & null & -- & nessun effetto (nessuna eccezione, getLockLevel(pc) == LOCK\_NONE) \\
non gestito & noto & 0 & non null & azione & nessun effetto (nessuna eccezione, getLockLevel(pc) == LOCK\_NONE) \\
non gestito & noto & 0 & non null & eccezione & eccezione propagata da lock \\
non gestito & noto & positivo & null & -- & nessun effetto (nessuna eccezione, getLockLevel(pc) == LOCK\_NONE) \\
non gestito & noto & positivo & non null & azione & nessun effetto (nessuna eccezione, getLockLevel(pc) == LOCK\_NONE) \\
non gestito & noto & positivo & non null & eccezione & eccezione propagata da lock \\
non gestito & non noto & altro negativo & null & -- & nessun effetto (nessuna eccezione, getLockLevel(pc) == LOCK\_NONE) \\
non gestito & non noto & altro negativo & non null & azione & nessun effetto (nessuna eccezione, getLockLevel(pc) == LOCK\_NONE) \\
non gestito & non noto & altro negativo & non null & eccezione & eccezione propagata da lock \\
non gestito & non noto & -1 & null & -- & nessun effetto (nessuna eccezione, getLockLevel(pc) == LOCK\_NONE) \\
non gestito & non noto & -1 & non null & azione & nessun effetto (nessuna eccezione, getLockLevel(pc) == LOCK\_NONE) \\
non gestito & non noto & -1 & non null & eccezione & eccezione propagata da lock \\
non gestito & non noto & 0 & null & -- & nessun effetto (nessuna eccezione, getLockLevel(pc) == LOCK\_NONE) \\
non gestito & non noto & 0 & non null & azione & nessun effetto (nessuna eccezione, getLockLevel(pc) == LOCK\_NONE) \\
non gestito & non noto & 0 & non null & eccezione & eccezione propagata da lock \\
non gestito & non noto & positivo & null & -- & nessun effetto (nessuna eccezione, getLockLevel(pc) == LOCK\_NONE) \\
non gestito & non noto & positivo & non null & azione & nessun effetto (nessuna eccezione, getLockLevel(pc) == LOCK\_NONE) \\
non gestito & non noto & positivo & non null & eccezione & eccezione propagata da lock \\
gestito non persistito & noto & altro negativo & null & -- & getLockLevel(pc) == level (postcondizione) \\
gestito non persistito & noto & altro negativo & non null & azione & nessun effetto (nessuna eccezione, getLockLevel(pc) == LOCK\_NONE) \\
gestito non persistito & noto & altro negativo & non null & eccezione & eccezione propagata da lock \\
gestito non persistito & noto & -1 & null & -- & getLockLevel(pc) == level (postcondizione) \\
gestito non persistito & noto & -1 & non null & azione & nessun effetto (nessuna eccezione, getLockLevel(pc) == LOCK\_NONE) \\
gestito non persistito & noto & -1 & non null & eccezione & eccezione propagata da lock \\
gestito non persistito & noto & 0 & null & -- & getLockLevel(pc) == level (postcondizione) \\
gestito non persistito & noto & 0 & non null & azione & nessun effetto (nessuna eccezione, getLockLevel(pc) == LOCK\_NONE) \\
gestito non persistito & noto & 0 & non null & eccezione & eccezione propagata da lock \\
gestito non persistito & noto & positivo & null & -- & getLockLevel(pc) == level (postcondizione) \\
gestito non persistito & noto & positivo & non null & azione & nessun effetto (nessuna eccezione, getLockLevel(pc) == LOCK\_NONE) \\
gestito non persistito & noto & positivo & non null & eccezione & eccezione propagata da lock \\
gestito non persistito & non noto & altro negativo & null & -- & getLockLevel(pc) == level (postcondizione) \\
gestito non persistito & non noto & altro negativo & non null & azione & nessun effetto (nessuna eccezione, getLockLevel(pc) == LOCK\_NONE) \\
gestito non persistito & non noto & altro negativo & non null & eccezione & eccezione propagata da lock \\
gestito non persistito & non noto & -1 & null & -- & getLockLevel(pc) == level (postcondizione) \\
gestito non persistito & non noto & -1 & non null & azione & nessun effetto (nessuna eccezione, getLockLevel(pc) == LOCK\_NONE) \\
gestito non persistito & non noto & -1 & non null & eccezione & eccezione propagata da lock \\
gestito non persistito & non noto & 0 & null & -- & getLockLevel(pc) == level (postcondizione) \\
gestito non persistito & non noto & 0 & non null & azione & nessun effetto (nessuna eccezione, getLockLevel(pc) == LOCK\_NONE) \\
gestito non persistito & non noto & 0 & non null & eccezione & eccezione propagata da lock \\
gestito non persistito & non noto & positivo & null & -- & getLockLevel(pc) == level (postcondizione) \\
gestito non persistito & non noto & positivo & non null & azione & nessun effetto (nessuna eccezione, getLockLevel(pc) == LOCK\_NONE) \\
gestito non persistito & non noto & positivo & non null & eccezione & eccezione propagata da lock \\
gestito persistito & noto & altro negativo & null & -- & getLockLevel(pc) == level (postcondizione) \\
gestito persistito & noto & altro negativo & non null & azione & nessun effetto (nessuna eccezione, getLockLevel(pc) == LOCK\_NONE) \\
gestito persistito & noto & altro negativo & non null & eccezione & eccezione propagata da lock \\
gestito persistito & noto & -1 & null & -- & getLockLevel(pc) == level (postcondizione) \\
gestito persistito & noto & -1 & non null & azione & nessun effetto (nessuna eccezione, getLockLevel(pc) == LOCK\_NONE) \\
gestito persistito & noto & -1 & non null & eccezione & eccezione propagata da lock \\
gestito persistito & noto & 0 & null & -- & getLockLevel(pc) == level (postcondizione) \\
gestito persistito & noto & 0 & non null & azione & nessun effetto (nessuna eccezione, getLockLevel(pc) == LOCK\_NONE) \\
gestito persistito & noto & 0 & non null & eccezione & eccezione propagata da lock \\
gestito persistito & noto & positivo & null & -- & getLockLevel(pc) == level (postcondizione) \\
gestito persistito & noto & positivo & non null & azione & nessun effetto (nessuna eccezione, getLockLevel(pc) == LOCK\_NONE) \\
gestito persistito & noto & positivo & non null & eccezione & eccezione propagata da lock \\
gestito persistito & non noto & altro negativo & null & -- & getLockLevel(pc) == level (postcondizione) \\
gestito persistito & non noto & altro negativo & non null & azione & nessun effetto (nessuna eccezione, getLockLevel(pc) == LOCK\_NONE) \\
gestito persistito & non noto & altro negativo & non null & eccezione & eccezione propagata da lock \\
gestito persistito & non noto & -1 & null & -- & getLockLevel(pc) == level (postcondizione) \\
gestito persistito & non noto & -1 & non null & azione & nessun effetto (nessuna eccezione, getLockLevel(pc) == LOCK\_NONE) \\
gestito persistito & non noto & -1 & non null & eccezione & eccezione propagata da lock \\
gestito persistito & non noto & 0 & null & -- & getLockLevel(pc) == level (postcondizione) \\
gestito persistito & non noto & 0 & non null & azione & nessun effetto (nessuna eccezione, getLockLevel(pc) == LOCK\_NONE) \\
gestito persistito & non noto & 0 & non null & eccezione & eccezione propagata da lock \\
gestito persistito & non noto & positivo & null & -- & getLockLevel(pc) == level (postcondizione) \\
gestito persistito & non noto & positivo & non null & azione & nessun effetto (nessuna eccezione, getLockLevel(pc) == LOCK\_NONE) \\
gestito persistito & non noto & positivo & non null & eccezione & eccezione propagata da lock \\
null & noto & altro negativo & non null & azione (ACT\_CASCADE) & nessun effetto (nessuna eccezione, getLockLevel(pc) == LOCK\_NONE) \\
null & noto & altro negativo & non null & azione (ACT\_RUN) & nessun effetto (nessuna eccezione, getLockLevel(pc) == LOCK\_NONE) \\
null & noto & -1 & non null & azione (ACT\_CASCADE) & nessun effetto (nessuna eccezione, getLockLevel(pc) == LOCK\_NONE) \\
null & noto & -1 & non null & azione (ACT\_RUN) & nessun effetto (nessuna eccezione, getLockLevel(pc) == LOCK\_NONE) \\
null & noto & 0 & non null & azione (ACT\_CASCADE) & nessun effetto (nessuna eccezione, getLockLevel(pc) == LOCK\_NONE) \\
null & noto & 0 & non null & azione (ACT\_RUN) & nessun effetto (nessuna eccezione, getLockLevel(pc) == LOCK\_NONE) \\
null & noto & positivo & non null & azione (ACT\_CASCADE) & nessun effetto (nessuna eccezione, getLockLevel(pc) == LOCK\_NONE) \\
null & noto & positivo & non null & azione (ACT\_RUN) & nessun effetto (nessuna eccezione, getLockLevel(pc) == LOCK\_NONE) \\
null & non noto & altro negativo & non null & azione (ACT\_CASCADE) & nessun effetto (nessuna eccezione, getLockLevel(pc) == LOCK\_NONE) \\
null & non noto & altro negativo & non null & azione (ACT\_RUN) & nessun effetto (nessuna eccezione, getLockLevel(pc) == LOCK\_NONE) \\
null & non noto & -1 & non null & azione (ACT\_CASCADE) & nessun effetto (nessuna eccezione, getLockLevel(pc) == LOCK\_NONE) \\
null & non noto & -1 & non null & azione (ACT\_RUN) & nessun effetto (nessuna eccezione, getLockLevel(pc) == LOCK\_NONE) \\
null & non noto & 0 & non null & azione (ACT\_CASCADE) & nessun effetto (nessuna eccezione, getLockLevel(pc) == LOCK\_NONE) \\
null & non noto & 0 & non null & azione (ACT\_RUN) & nessun effetto (nessuna eccezione, getLockLevel(pc) == LOCK\_NONE) \\
null & non noto & positivo & non null & azione (ACT\_CASCADE) & nessun effetto (nessuna eccezione, getLockLevel(pc) == LOCK\_NONE) \\
null & non noto & positivo & non null & azione (ACT\_RUN) & nessun effetto (nessuna eccezione, getLockLevel(pc) == LOCK\_NONE) \\
non gestito & noto & altro negativo & non null & azione (ACT\_CASCADE) & nessun effetto (nessuna eccezione, getLockLevel(pc) == LOCK\_NONE) \\
non gestito & noto & altro negativo & non null & azione (ACT\_RUN) & nessun effetto (nessuna eccezione, getLockLevel(pc) == LOCK\_NONE) \\
non gestito & noto & -1 & non null & azione (ACT\_CASCADE) & nessun effetto (nessuna eccezione, getLockLevel(pc) == LOCK\_NONE) \\
non gestito & noto & -1 & non null & azione (ACT\_RUN) & nessun effetto (nessuna eccezione, getLockLevel(pc) == LOCK\_NONE) \\
non gestito & noto & 0 & non null & azione (ACT\_CASCADE) & nessun effetto (nessuna eccezione, getLockLevel(pc) == LOCK\_NONE) \\
non gestito & noto & 0 & non null & azione (ACT\_RUN) & nessun effetto (nessuna eccezione, getLockLevel(pc) == LOCK\_NONE) \\
non gestito & noto & positivo & non null & azione (ACT\_CASCADE) & nessun effetto (nessuna eccezione, getLockLevel(pc) == LOCK\_NONE) \\
non gestito & noto & positivo & non null & azione (ACT\_RUN) & nessun effetto (nessuna eccezione, getLockLevel(pc) == LOCK\_NONE) \\
non gestito & non noto & altro negativo & non null & azione (ACT\_CASCADE) & nessun effetto (nessuna eccezione, getLockLevel(pc) == LOCK\_NONE) \\
non gestito & non noto & altro negativo & non null & azione (ACT\_RUN) & nessun effetto (nessuna eccezione, getLockLevel(pc) == LOCK\_NONE) \\
non gestito & non noto & -1 & non null & azione (ACT\_CASCADE) & nessun effetto (nessuna eccezione, getLockLevel(pc) == LOCK\_NONE) \\
non gestito & non noto & -1 & non null & azione (ACT\_RUN) & nessun effetto (nessuna eccezione, getLockLevel(pc) == LOCK\_NONE) \\
non gestito & non noto & 0 & non null & azione (ACT\_CASCADE) & nessun effetto (nessuna eccezione, getLockLevel(pc) == LOCK\_NONE) \\
non gestito & non noto & 0 & non null & azione (ACT\_RUN) & nessun effetto (nessuna eccezione, getLockLevel(pc) == LOCK\_NONE) \\
non gestito & non noto & positivo & non null & azione (ACT\_CASCADE) & nessun effetto (nessuna eccezione, getLockLevel(pc) == LOCK\_NONE) \\
non gestito & non noto & positivo & non null & azione (ACT\_RUN) & nessun effetto (nessuna eccezione, getLockLevel(pc) == LOCK\_NONE) \\
gestito non persistito & noto & altro negativo & non null & azione (ACT\_CASCADE) & nessun effetto (nessuna eccezione, getLockLevel(pc) == LOCK\_NONE) \\
gestito non persistito & noto & altro negativo & non null & azione (ACT\_RUN) & getLockLevel(pc) == level (postcondizione) \\
gestito non persistito & noto & -1 & non null & azione (ACT\_CASCADE) & nessun effetto (nessuna eccezione, getLockLevel(pc) == LOCK\_NONE) \\
gestito non persistito & noto & -1 & non null & azione (ACT\_RUN) & getLockLevel(pc) == level (postcondizione) \\
gestito non persistito & noto & 0 & non null & azione (ACT\_CASCADE) & nessun effetto (nessuna eccezione, getLockLevel(pc) == LOCK\_NONE) \\
gestito non persistito & noto & 0 & non null & azione (ACT\_RUN) & getLockLevel(pc) == level (postcondizione) \\
gestito non persistito & noto & positivo & non null & azione (ACT\_CASCADE) & nessun effetto (nessuna eccezione, getLockLevel(pc) == LOCK\_NONE) \\
gestito non persistito & noto & positivo & non null & azione (ACT\_RUN) & getLockLevel(pc) == level (postcondizione) \\
gestito non persistito & non noto & altro negativo & non null & azione (ACT\_CASCADE) & nessun effetto (nessuna eccezione, getLockLevel(pc) == LOCK\_NONE) \\
gestito non persistito & non noto & altro negativo & non null & azione (ACT\_RUN) & getLockLevel(pc) == level (postcondizione) \\
gestito non persistito & non noto & -1 & non null & azione (ACT\_CASCADE) & nessun effetto (nessuna eccezione, getLockLevel(pc) == LOCK\_NONE) \\
gestito non persistito & non noto & -1 & non null & azione (ACT\_RUN) & getLockLevel(pc) == level (postcondizione) \\
gestito non persistito & non noto & 0 & non null & azione (ACT\_CASCADE) & nessun effetto (nessuna eccezione, getLockLevel(pc) == LOCK\_NONE) \\
gestito non persistito & non noto & 0 & non null & azione (ACT\_RUN) & getLockLevel(pc) == level (postcondizione) \\
gestito non persistito & non noto & positivo & non null & azione (ACT\_CASCADE) & nessun effetto (nessuna eccezione, getLockLevel(pc) == LOCK\_NONE) \\
gestito non persistito & non noto & positivo & non null & azione (ACT\_RUN) & getLockLevel(pc) == level (postcondizione) \\
gestito persistito & noto & altro negativo & non null & azione (ACT\_CASCADE) & nessun effetto (nessuna eccezione, getLockLevel(pc) == LOCK\_NONE) \\
gestito persistito & noto & altro negativo & non null & azione (ACT\_RUN) & getLockLevel(pc) == level (postcondizione) \\
gestito persistito & noto & -1 & non null & azione (ACT\_CASCADE) & nessun effetto (nessuna eccezione, getLockLevel(pc) == LOCK\_NONE) \\
gestito persistito & noto & -1 & non null & azione (ACT\_RUN) & getLockLevel(pc) == level (postcondizione) \\
gestito persistito & noto & 0 & non null & azione (ACT\_CASCADE) & nessun effetto (nessuna eccezione, getLockLevel(pc) == LOCK\_NONE) \\
gestito persistito & noto & 0 & non null & azione (ACT\_RUN) & getLockLevel(pc) == level (postcondizione) \\
gestito persistito & noto & positivo & non null & azione (ACT\_CASCADE) & nessun effetto (nessuna eccezione, getLockLevel(pc) == LOCK\_NONE) \\
gestito persistito & noto & positivo & non null & azione (ACT\_RUN) & getLockLevel(pc) == level (postcondizione) \\
gestito persistito & non noto & altro negativo & non null & azione (ACT\_CASCADE) & nessun effetto (nessuna eccezione, getLockLevel(pc) == LOCK\_NONE) \\
gestito persistito & non noto & altro negativo & non null & azione (ACT\_RUN) & getLockLevel(pc) == level (postcondizione) \\
gestito persistito & non noto & -1 & non null & azione (ACT\_CASCADE) & nessun effetto (nessuna eccezione, getLockLevel(pc) == LOCK\_NONE) \\
gestito persistito & non noto & -1 & non null & azione (ACT\_RUN) & getLockLevel(pc) == level (postcondizione) \\
gestito persistito & non noto & 0 & non null & azione (ACT\_CASCADE) & nessun effetto (nessuna eccezione, getLockLevel(pc) == LOCK\_NONE) \\
gestito persistito & non noto & 0 & non null & azione (ACT\_RUN) & getLockLevel(pc) == level (postcondizione) \\
gestito persistito & non noto & positivo & non null & azione (ACT\_CASCADE) & nessun effetto (nessuna eccezione, getLockLevel(pc) == LOCK\_NONE) \\
gestito persistito & non noto & positivo & non null & azione (ACT\_RUN) & getLockLevel(pc) == level (postcondizione) \\
\bottomrule
\caption{Tuple corrette per \texttt{lock(Object pc, int level, int timeout, OpCallbacks call)}.}
\label{tab:lock-tuple-corrette}
\end{longtable}
}



\paragraph*{Tuple di progettazione, \texttt{lock(Object pc, int level, int timeout, OpCallbacks call), approfondimento azione} (BrokerImpl).}

{\footnotesize\setlength{\tabcolsep}{3pt}
\begin{longtable}{>{\raggedright\arraybackslash}p{0.314\textwidth}>{\raggedright\arraybackslash}p{0.285\textwidth}>{\raggedright\arraybackslash}p{0.328\textwidth}}
\toprule
\textbf{pc} & \textbf{call} & \textbf{oracolo} \\
\midrule
\endfirsthead
\toprule
\textbf{pc} & \textbf{call} & \textbf{oracolo} \\
\midrule
\endhead
\bottomrule
\endfoot
null & null & nessun effetto \\
null & azione (ACT\_NONE) & nessun effetto \\
null & azione (ACT\_CASCADE) & nessun effetto \\
null & azione (ACT\_RUN) & nessun effetto \\
null & lancia eccezione & eccezione non propagata \\
non gestito & null & nessun effetto \\
non gestito & azione (ACT\_NONE) & nessun effetto \\
non gestito & azione (ACT\_CASCADE) & nessun effetto \\
non gestito & azione (ACT\_RUN) & nessun effetto \\
non gestito & lancia eccezione & eccezione propagata \\
gestito non persistito & null & lock applicato \\
gestito non persistito & azione (ACT\_NONE) & nessun effetto \\
gestito non persistito & azione (ACT\_CASCADE) & nessun effetto \\
gestito non persistito & azione (ACT\_RUN) & lock applicato \\
gestito non persistito & lancia eccezione & eccezione propagata \\
gestito persistito & null & lock applicato \\
gestito persistito & azione (ACT\_NONE) & nessun effetto \\
gestito persistito & azione (ACT\_CASCADE) & nessun effetto \\
gestito persistito & azione (ACT\_RUN) & lock applicato \\
gestito persistito & lancia eccezione & eccezione propagata \\
\bottomrule
\caption{Tuple di progettazione per \texttt{lock(Object pc, int level, int timeout, OpCallbacks call), approfondimento azione}.}
\label{tab:lock2-tuple}
\end{longtable}
}



\paragraph*{Tuple di progettazione, \texttt{detachAll(Collection objs, OpCallbacks call)} (BrokerImpl).}

{\footnotesize\setlength{\tabcolsep}{3pt}
\begin{longtable}{>{\raggedright\arraybackslash}p{0.359\textwidth}>{\raggedright\arraybackslash}p{0.068\textwidth}>{\raggedright\arraybackslash}p{0.128\textwidth}>{\raggedright\arraybackslash}p{0.359\textwidth}}
\toprule
\textbf{objs} & \textbf{call} & \textbf{processArgument} & \textbf{oracolo} \\
\midrule
\endfirsthead
\toprule
\textbf{objs} & \textbf{call} & \textbf{processArgument} & \textbf{oracolo} \\
\midrule
\endhead
\bottomrule
\endfoot
null & null & -- & eccezione (tipo non specificato) \\
null & non null & azione & eccezione (tipo non specificato) \\
null & non null & eccezione & eccezione propagata da detachAll \\
vuota & null & -- & array vuoto \\
vuota & non null & azione & array vuoto \\
vuota & non null & eccezione & eccezione propagata da detachAll \\
un elemento gestito & null & -- & array con un elemento; isDetached(elemento) == true \\
un elemento gestito & non null & azione & array con un elemento; isDetached(elemento) == true \\
un elemento gestito & non null & eccezione & eccezione propagata da detachAll \\
un elemento non gestito & null & -- & eccezione (tipo non specificato) \\
un elemento non gestito & non null & azione & eccezione (tipo non specificato) \\
un elemento non gestito & non null & eccezione & eccezione propagata da detachAll \\
elementi misti (un gestito e un non gestito) & null & -- & eccezione (tipo non specificato) \\
elementi misti (un gestito e un non gestito) & non null & azione & eccezione (tipo non specificato) \\
elementi misti (un gestito e un non gestito) & non null & eccezione & eccezione propagata da detachAll \\
\bottomrule
\caption{Tuple di progettazione per \texttt{detachAll(Collection objs, OpCallbacks call)}.}
\label{tab:detachall-tuple}
\end{longtable}
}



\paragraph*{Tuple corrette, \texttt{detachAll(Collection objs, OpCallbacks call)} (BrokerImpl).}

{\footnotesize\setlength{\tabcolsep}{3pt}
\begin{longtable}{>{\raggedright\arraybackslash}p{0.343\textwidth}>{\raggedright\arraybackslash}p{0.065\textwidth}>{\raggedright\arraybackslash}p{0.163\textwidth}>{\raggedright\arraybackslash}p{0.343\textwidth}}
\toprule
\textbf{objs} & \textbf{call} & \textbf{processArgument} & \textbf{oracolo} \\
\midrule
\endfirsthead
\toprule
\textbf{objs} & \textbf{call} & \textbf{processArgument} & \textbf{oracolo} \\
\midrule
\endhead
\bottomrule
\endfoot
null & null & -- & detachAll ritorna null \\
null & non null & azione & detachAll ritorna null \\
null & non null & eccezione & detachAll ritorna null \\
vuota & null & -- & array vuoto \\
vuota & non null & azione & array vuoto \\
vuota & non null & eccezione & array vuoto \\
un elemento gestito & null & -- & array con un elemento; isDetached(elemento) == true \\
un elemento gestito & non null & azione & array con un elemento; l'elemento e' invariato (stessa istanza, non staccato) \\
un elemento gestito & non null & eccezione & eccezione propagata da detachAll \\
un elemento non gestito & null & -- & array con un elemento; l'elemento e' invariato (stessa istanza) \\
un elemento non gestito & non null & azione & array con un elemento; l'elemento e' invariato (stessa istanza) \\
un elemento non gestito & non null & eccezione & eccezione propagata da detachAll \\
elementi misti (un gestito e un non gestito) & null & -- & array di due elementi: il gestito e' una copia scollegata, il non gestito e' invariato \\
elementi misti (un gestito e un non gestito) & non null & azione & array di due elementi, entrambi invariati (nessun detach applicato) \\
elementi misti (un gestito e un non gestito) & non null & eccezione & eccezione propagata da detachAll \\
null & non null & azione (ACT\_CASCADE) & detachAll ritorna null \\
null & non null & azione (ACT\_RUN) & detachAll ritorna null \\
vuota & non null & azione (ACT\_CASCADE) & array vuoto \\
vuota & non null & azione (ACT\_RUN) & array vuoto \\
un elemento gestito & non null & azione (ACT\_CASCADE) & array con un elemento; l'elemento e' invariato (stessa istanza, non staccato) \\
un elemento gestito & non null & azione (ACT\_RUN) & array con un elemento; copia scollegata dal broker \\
un elemento non gestito & non null & azione (ACT\_CASCADE) & array con un elemento; l'elemento e' invariato (stessa istanza) \\
un elemento non gestito & non null & azione (ACT\_RUN) & array con un elemento; l'elemento e' invariato (stessa istanza) \\
elementi misti (un gestito e un non gestito) & non null & azione (ACT\_CASCADE) & array di due elementi, entrambi invariati (nessun detach applicato) \\
elementi misti (un gestito e un non gestito) & non null & azione (ACT\_RUN) & array di due elementi: il gestito e' una copia scollegata, il non gestito e' invariato \\
\bottomrule
\caption{Tuple corrette per \texttt{detachAll(Collection objs, OpCallbacks call)}.}
\label{tab:detachall-tuple-corrette}
\end{longtable}
}



\paragraph*{Tuple di progettazione, \texttt{detachAll(Collection objs, OpCallbacks call), approfondimento azione} (BrokerImpl).}

{\footnotesize\setlength{\tabcolsep}{3pt}
\begin{longtable}{>{\raggedright\arraybackslash}p{0.278\textwidth}>{\raggedright\arraybackslash}p{0.106\textwidth}>{\raggedright\arraybackslash}p{0.093\textwidth}>{\raggedright\arraybackslash}p{0.146\textwidth}>{\raggedright\arraybackslash}p{0.278\textwidth}}
\toprule
\textbf{stato\_broker} & \textbf{call} & \textbf{flush} & \textbf{transazione} & \textbf{oracolo} \\
\midrule
\endfirsthead
\toprule
\textbf{stato\_broker} & \textbf{call} & \textbf{flush} & \textbf{transazione} & \textbf{oracolo} \\
\midrule
\endhead
\bottomrule
\endfoot
senza istanze gestite & null & true & non attiva & nessun effetto \\
senza istanze gestite & null & true & attiva & nessun effetto \\
senza istanze gestite & null & false & non attiva & nessun effetto \\
senza istanze gestite & null & false & attiva & nessun effetto \\
senza istanze gestite & non null & true & non attiva & callback mai invocato \\
senza istanze gestite & non null & true & attiva & callback mai invocato \\
senza istanze gestite & non null & false & non attiva & callback mai invocato \\
senza istanze gestite & non null & false & attiva & callback mai invocato \\
con istanza gestita & null & false & attiva & InvalidStateException \\
con istanza gestita & non null & false & attiva & InvalidStateException \\
con istanza gestita & null & true & attiva & non determinato \\
con istanza gestita & non null & true & attiva & non determinato \\
\bottomrule
\caption{Tuple di progettazione per \texttt{detachAll(Collection objs, OpCallbacks call), approfondimento azione}.}
\label{tab:detachall2-tuple}
\end{longtable}
}



\paragraph*{Tuple di progettazione, \texttt{attachAll(Collection objs, boolean copyNew, OpCallbacks call)} (BrokerImpl).}

{\footnotesize\setlength{\tabcolsep}{3pt}
\begin{longtable}{>{\raggedright\arraybackslash}p{0.225\textwidth}>{\raggedright\arraybackslash}p{0.066\textwidth}>{\raggedright\arraybackslash}p{0.075\textwidth}>{\raggedright\arraybackslash}p{0.141\textwidth}>{\raggedright\arraybackslash}p{0.394\textwidth}}
\toprule
\textbf{objs} & \textbf{copyNew} & \textbf{call} & \textbf{processArgument} & \textbf{oracolo} \\
\midrule
\endfirsthead
\toprule
\textbf{objs} & \textbf{copyNew} & \textbf{call} & \textbf{processArgument} & \textbf{oracolo} \\
\midrule
\endhead
\bottomrule
\endfoot
null & true & null & -- & eccezione (tipo non specificato) \\
null & true & non null & azione & eccezione (tipo non specificato) \\
null & true & non null & eccezione & eccezione propagata da attachAll \\
null & false & null & -- & eccezione (tipo non specificato) \\
null & false & non null & azione & eccezione (tipo non specificato) \\
null & false & non null & eccezione & eccezione propagata da attachAll \\
vuota & true & null & -- & array vuoto \\
vuota & true & non null & azione & array vuoto \\
vuota & true & non null & eccezione & eccezione propagata da attachAll \\
vuota & false & null & -- & array vuoto \\
vuota & false & non null & azione & array vuoto \\
vuota & false & non null & eccezione & eccezione propagata da attachAll \\
all new & true & null & -- & istanza persistita come nuova, copia dell'originale \\
all new & true & non null & azione & istanza persistita come nuova, copia dell'originale \\
all new & true & non null & eccezione & eccezione propagata da attachAll \\
all new & false & null & -- & istanza persistita come nuova, non copiata \\
all new & false & non null & azione & istanza persistita come nuova, non copiata \\
all new & false & non null & eccezione & eccezione propagata da attachAll \\
all old & true & null & -- & istanza persistente esistente con le modifiche mergiate \\
all old & true & non null & azione & istanza persistente esistente con le modifiche mergiate \\
all old & true & non null & eccezione & eccezione propagata da attachAll \\
all old & false & null & -- & istanza persistente esistente con le modifiche mergiate \\
all old & false & non null & azione & istanza persistente esistente con le modifiche mergiate \\
all old & false & non null & eccezione & eccezione propagata da attachAll \\
all attached & true & null & -- & nessuna modifica: l'istanza resta quella gia' gestita \\
all attached & true & non null & azione & nessuna modifica: l'istanza resta quella gia' gestita \\
all attached & true & non null & eccezione & eccezione propagata da attachAll \\
all attached & false & null & -- & nessuna modifica: l'istanza resta quella gia' gestita \\
all attached & false & non null & azione & nessuna modifica: l'istanza resta quella gia' gestita \\
all attached & false & non null & eccezione & eccezione propagata da attachAll \\
misti: validi eterogenei & true & null & -- & ogni elemento segue il proprio ramo secondo il tipo: new -> persistito come nuovo, copia dell'originale; old -> istanza persistente aggiornata con le modifiche mergiate; attached -> nessuna modifica \\
misti: validi eterogenei & true & non null & azione & ogni elemento segue il proprio ramo secondo il tipo: new -> persistito come nuovo, copia dell'originale; old -> istanza persistente aggiornata con le modifiche mergiate; attached -> nessuna modifica \\
misti: validi eterogenei & true & non null & eccezione & ogni elemento segue il proprio ramo secondo il tipo: new -> persistito come nuovo, copia dell'originale; old -> istanza persistente aggiornata con le modifiche mergiate; attached -> nessuna modifica \\
misti: validi eterogenei & false & null & -- & ogni elemento segue il proprio ramo secondo il tipo: new -> persistito come nuovo, non copiato; old -> istanza persistente aggiornata con le modifiche mergiate; attached -> nessuna modifica \\
misti: validi eterogenei & false & non null & azione & ogni elemento segue il proprio ramo secondo il tipo: new -> persistito come nuovo, non copiato; old -> istanza persistente aggiornata con le modifiche mergiate; attached -> nessuna modifica \\
misti: validi eterogenei & false & non null & eccezione & ogni elemento segue il proprio ramo secondo il tipo: new -> persistito come nuovo, non copiato; old -> istanza persistente aggiornata con le modifiche mergiate; attached -> nessuna modifica \\
misti: almeno 1 invalido & true & null & -- & eccezione (tipo non specificato) \\
misti: almeno 1 invalido & true & non null & azione & eccezione (tipo non specificato) \\
misti: almeno 1 invalido & true & non null & eccezione & eccezione propagata da attachAll \\
misti: almeno 1 invalido & false & null & -- & eccezione (tipo non specificato) \\
misti: almeno 1 invalido & false & non null & azione & eccezione (tipo non specificato) \\
misti: almeno 1 invalido & false & non null & eccezione & eccezione propagata da attachAll \\
\bottomrule
\caption{Tuple di progettazione per \texttt{attachAll(Collection objs, boolean copyNew, OpCallbacks call)}.}
\label{tab:attachall-tuple}
\end{longtable}
}



\paragraph*{Tuple corrette, \texttt{attachAll(Collection objs, boolean copyNew, OpCallbacks call)} (BrokerImpl).}

{\footnotesize\setlength{\tabcolsep}{3pt}
\begin{longtable}{>{\raggedright\arraybackslash}p{0.214\textwidth}>{\raggedright\arraybackslash}p{0.062\textwidth}>{\raggedright\arraybackslash}p{0.071\textwidth}>{\raggedright\arraybackslash}p{0.178\textwidth}>{\raggedright\arraybackslash}p{0.375\textwidth}}
\toprule
\textbf{objs} & \textbf{copyNew} & \textbf{call} & \textbf{processArgument} & \textbf{oracolo} \\
\midrule
\endfirsthead
\toprule
\textbf{objs} & \textbf{copyNew} & \textbf{call} & \textbf{processArgument} & \textbf{oracolo} \\
\midrule
\endhead
\bottomrule
\endfoot
null & true & null & -- & attachAll ritorna null \\
null & true & non null & azione & attachAll ritorna null \\
null & true & non null & eccezione & attachAll ritorna null \\
null & false & null & -- & attachAll ritorna null \\
null & false & non null & azione & attachAll ritorna null \\
null & false & non null & eccezione & attachAll ritorna null \\
vuota & true & null & -- & array vuoto \\
vuota & true & non null & azione & array vuoto \\
vuota & true & non null & eccezione & array vuoto \\
vuota & false & null & -- & array vuoto \\
vuota & false & non null & azione & array vuoto \\
vuota & false & non null & eccezione & array vuoto \\
all new & true & null & -- & istanza persistita come nuova, copia dell'originale \\
all new & true & non null & azione & elemento invariato (nessun attach eseguito) \\
all new & true & non null & eccezione & eccezione propagata da attachAll \\
all new & false & null & -- & istanza persistita come nuova, non copiata \\
all new & false & non null & azione & elemento invariato (nessun attach eseguito) \\
all new & false & non null & eccezione & eccezione propagata da attachAll \\
all old & true & null & -- & istanza persistente esistente con le modifiche mergiate \\
all old & true & non null & azione & elemento invariato (nessun attach eseguito) \\
all old & true & non null & eccezione & eccezione propagata da attachAll \\
all old & false & null & -- & istanza persistente esistente con le modifiche mergiate \\
all old & false & non null & azione & elemento invariato (nessun attach eseguito) \\
all old & false & non null & eccezione & eccezione propagata da attachAll \\
all attached & true & null & -- & istanza persistita come nuova, copia dell'originale \\
all attached & true & non null & azione & elemento invariato (nessun attach eseguito) \\
all attached & true & non null & eccezione & eccezione propagata da attachAll \\
all attached & false & null & -- & istanza persistita come nuova, non copiata \\
all attached & false & non null & azione & elemento invariato (nessun attach eseguito) \\
all attached & false & non null & eccezione & eccezione propagata da attachAll \\
misti: validi eterogenei & true & null & -- & ogni elemento segue il proprio ramo secondo il tipo: new -> persistito come nuovo, copia dell'originale; old -> istanza persistente aggiornata con le modifiche mergiate; attached -> nessuna modifica \\
misti: validi eterogenei & true & non null & azione & entrambi gli elementi invariati (nessun attach eseguito) \\
misti: validi eterogenei & true & non null & eccezione & UserException con le eccezioni di processArgument (una per elemento) annidate, non la stessa istanza \\
misti: validi eterogenei & false & null & -- & ogni elemento segue il proprio ramo secondo il tipo: new -> persistito come nuovo, non copiato; old -> istanza persistente aggiornata con le modifiche mergiate; attached -> nessuna modifica \\
misti: validi eterogenei & false & non null & azione & entrambi gli elementi invariati (nessun attach eseguito) \\
misti: validi eterogenei & false & non null & eccezione & UserException con le eccezioni di processArgument (una per elemento) annidate, non la stessa istanza \\
misti: almeno 1 invalido & true & null & -- & eccezione (tipo non specificato) \\
misti: almeno 1 invalido & true & non null & azione & entrambi gli elementi invariati (nessun attach eseguito) \\
misti: almeno 1 invalido & true & non null & eccezione & UserException con le eccezioni di processArgument (una per elemento) annidate, non la stessa istanza \\
misti: almeno 1 invalido & false & null & -- & eccezione (tipo non specificato) \\
misti: almeno 1 invalido & false & non null & azione & entrambi gli elementi invariati (nessun attach eseguito) \\
misti: almeno 1 invalido & false & non null & eccezione & UserException con le eccezioni di processArgument (una per elemento) annidate, non la stessa istanza \\
all new & true & non null & azione (ACT\_CASCADE) & eccezione (NullPointerException interno) \\
all new & true & non null & azione (ACT\_RUN) & istanza persistita come nuova, copia dell'originale \\
all old & true & non null & azione (ACT\_CASCADE) & eccezione (NullPointerException interno) \\
all old & true & non null & azione (ACT\_RUN) & istanza persistente esistente con le modifiche mergiate \\
all attached & true & non null & azione (ACT\_CASCADE) & elemento invariato (nessun attach eseguito) \\
all attached & true & non null & azione (ACT\_RUN) & istanza persistita come nuova, copia dell'originale \\
misti: validi eterogenei & true & non null & azione (ACT\_CASCADE) & eccezione (NullPointerException interno) \\
misti: validi eterogenei & true & non null & azione (ACT\_RUN) & entrambi gli elementi pienamente processati secondo il proprio ramo \\
misti: almeno 1 invalido & true & non null & azione (ACT\_CASCADE) & eccezione (NullPointerException interno) \\
misti: almeno 1 invalido & true & non null & azione (ACT\_RUN) & eccezione (tipo non specificato) \\
\bottomrule
\caption{Tuple corrette per \texttt{attachAll(Collection objs, boolean copyNew, OpCallbacks call)}.}
\label{tab:attachall-tuple-corrette}
\end{longtable}
}



\paragraph*{Tuple di progettazione, \texttt{newInstance(Class cls)} (BrokerImpl).}

{\footnotesize\setlength{\tabcolsep}{3pt}
\begin{longtable}{>{\raggedright\arraybackslash}p{0.434\textwidth}>{\raggedright\arraybackslash}p{0.507\textwidth}}
\toprule
\textbf{cls} & \textbf{oracolo} \\
\midrule
\endfirsthead
\toprule
\textbf{cls} & \textbf{oracolo} \\
\midrule
\endhead
\bottomrule
\endfoot
null & eccezione (tipo non specificato) \\
interfaccia con metodi JavaBeans & istanza non-null (implementazione concreta generata dai metadati) \\
interfaccia senza metodi JavaBeans & eccezione (tipo non specificato, possibile IllegalArgumentException) \\
classe astratta con metodi JavaBeans & istanza non-null (implementazione concreta generata dai metadati) \\
classe astratta senza JavaBeans & eccezione (tipo non specificato, possibile IllegalArgumentException) \\
tipo managed concreto & istanza non-null del tipo specificato \\
tipo non managed concreto & lancia IllegalArgumentException \\
\bottomrule
\caption{Tuple di progettazione per \texttt{newInstance(Class cls)}.}
\label{tab:newinstance-tuple}
\end{longtable}
}



\paragraph*{Tuple corrette, \texttt{newInstance(Class cls)} (BrokerImpl).}

{\footnotesize\setlength{\tabcolsep}{3pt}
\begin{longtable}{>{\raggedright\arraybackslash}p{0.434\textwidth}>{\raggedright\arraybackslash}p{0.507\textwidth}}
\toprule
\textbf{cls} & \textbf{oracolo} \\
\midrule
\endfirsthead
\toprule
\textbf{cls} & \textbf{oracolo} \\
\midrule
\endhead
\bottomrule
\endfoot
null & eccezione (tipo non specificato) \\
interfaccia con metodi JavaBeans & istanza non-null (implementazione concreta generata dai metadati) \\
interfaccia senza metodi JavaBeans & eccezione (tipo non specificato, possibile IllegalArgumentException) \\
classe astratta con metodi JavaBeans & eccezione (tipo non specificato) \\
classe astratta senza JavaBeans & eccezione (tipo non specificato, possibile IllegalArgumentException) \\
tipo managed concreto & istanza non-null del tipo specificato \\
tipo non managed concreto & lancia IllegalArgumentException \\
\bottomrule
\caption{Tuple corrette per \texttt{newInstance(Class cls)}.}
\label{tab:newinstance-tuple-corrette}
\end{longtable}
}



\paragraph*{Tuple di progettazione, \texttt{newInstance(Class cls), approfondimento mutation testing} (BrokerImpl).}

{\footnotesize\setlength{\tabcolsep}{3pt}
\begin{longtable}{>{\raggedright\arraybackslash}p{0.480\textwidth}>{\raggedright\arraybackslash}p{0.461\textwidth}}
\toprule
\textbf{cls} & \textbf{oracolo} \\
\midrule
\endfirsthead
\toprule
\textbf{cls} & \textbf{oracolo} \\
\midrule
\endhead
\bottomrule
\endfoot
tipo non managed concreto & IllegalArgumentException \\
\bottomrule
\caption{Tuple di progettazione per \texttt{newInstance(Class cls), approfondimento mutation testing}.}
\label{tab:newinstance2-tuple}
\end{longtable}
}



\paragraph*{Tuple di progettazione, \texttt{isDetached(Object obj, boolean find)} (BrokerImpl).}

{\footnotesize\setlength{\tabcolsep}{3pt}
\begin{longtable}{>{\raggedright\arraybackslash}p{0.272\textwidth}>{\raggedright\arraybackslash}p{0.173\textwidth}>{\raggedright\arraybackslash}p{0.297\textwidth}>{\raggedright\arraybackslash}p{0.173\textwidth}}
\toprule
\textbf{obj} & \textbf{find} & \textbf{store} & \textbf{oracolo} \\
\midrule
\endfirsthead
\toprule
\textbf{obj} & \textbf{find} & \textbf{store} & \textbf{oracolo} \\
\midrule
\endhead
\bottomrule
\endfoot
null & true & -- & false \\
null & false & -- & false \\
gestito & true & presente & false \\
gestito & true & non presente & false \\
gestito & false & presente & false \\
gestito & false & non presente & false \\
non gestito & false & presente & false \\
non gestito & false & non presente & false \\
non gestito & true & presente & true \\
non gestito & true & non presente & false \\
\bottomrule
\caption{Tuple di progettazione per \texttt{isDetached(Object obj, boolean find)}.}
\label{tab:isdetached-tuple}
\end{longtable}
}



\paragraph*{Tuple di progettazione, \texttt{persist(Object pc, Object id, OpCallbacks call)} (BrokerImpl).}

{\footnotesize\setlength{\tabcolsep}{3pt}
\begin{longtable}{>{\raggedright\arraybackslash}p{0.106\textwidth}>{\raggedright\arraybackslash}p{0.077\textwidth}>{\raggedright\arraybackslash}p{0.077\textwidth}>{\raggedright\arraybackslash}p{0.077\textwidth}>{\raggedright\arraybackslash}p{0.145\textwidth}>{\raggedright\arraybackslash}p{0.405\textwidth}}
\toprule
\textbf{pc} & \textbf{id} & \textbf{cascade} & \textbf{call} & \textbf{processArgument} & \textbf{oracolo} \\
\midrule
\endfirsthead
\toprule
\textbf{pc} & \textbf{id} & \textbf{cascade} & \textbf{call} & \textbf{processArgument} & \textbf{oracolo} \\
\midrule
\endhead
\bottomrule
\endfoot
null & null & -- & null & -- & eccezione (tipo non specificato) \\
null & null & -- & non null & azione & eccezione (tipo non specificato) \\
null & null & -- & non null & eccezione & eccezione propagata da persist \\
null & non null & -- & null & -- & eccezione (tipo non specificato) \\
null & non null & -- & non null & azione & eccezione (tipo non specificato) \\
null & non null & -- & non null & eccezione & eccezione propagata da persist \\
non gestito & null & presente & null & -- & state manager non-null, id di default assegnato; istanza correlata (CASCADE\_IMMEDIATE) non cascata, resta non gestita \\
non gestito & null & presente & non null & azione & state manager non-null, id di default assegnato; istanza correlata (CASCADE\_IMMEDIATE) non cascata, resta non gestita \\
non gestito & null & presente & non null & eccezione & eccezione propagata da persist \\
non gestito & null & assente & null & -- & state manager non-null, id di default assegnato \\
non gestito & null & assente & non null & azione & state manager non-null, id di default assegnato \\
non gestito & null & assente & non null & eccezione & eccezione propagata da persist \\
non gestito & non null & presente & null & -- & state manager non-null, con l'id specificato assegnato; istanza correlata (CASCADE\_IMMEDIATE) non cascata, resta non gestita \\
non gestito & non null & presente & non null & azione & state manager non-null, con l'id specificato assegnato; istanza correlata (CASCADE\_IMMEDIATE) non cascata, resta non gestita \\
non gestito & non null & presente & non null & eccezione & eccezione propagata da persist \\
non gestito & non null & assente & null & -- & state manager non-null, con l'id specificato assegnato \\
non gestito & non null & assente & non null & azione & state manager non-null, con l'id specificato assegnato \\
non gestito & non null & assente & non null & eccezione & eccezione propagata da persist \\
gestito & null & presente & null & -- & nessuna modifica: pc resta quello gia' gestito, id ignorato \\
gestito & null & presente & non null & azione & nessuna modifica: pc resta quello gia' gestito, id ignorato \\
gestito & null & presente & non null & eccezione & eccezione propagata da persist \\
gestito & null & assente & null & -- & nessuna modifica: pc resta quello gia' gestito, id ignorato \\
gestito & null & assente & non null & azione & nessuna modifica: pc resta quello gia' gestito, id ignorato \\
gestito & null & assente & non null & eccezione & eccezione propagata da persist \\
gestito & non null & presente & null & -- & nessuna modifica: pc resta quello gia' gestito, id ignorato \\
gestito & non null & presente & non null & azione & nessuna modifica: pc resta quello gia' gestito, id ignorato \\
gestito & non null & presente & non null & eccezione & eccezione propagata da persist \\
gestito & non null & assente & null & -- & nessuna modifica: pc resta quello gia' gestito, id ignorato \\
gestito & non null & assente & non null & azione & nessuna modifica: pc resta quello gia' gestito, id ignorato \\
gestito & non null & assente & non null & eccezione & eccezione propagata da persist \\
\bottomrule
\caption{Tuple di progettazione per \texttt{persist(Object pc, Object id, OpCallbacks call)}.}
\label{tab:persist-tuple}
\end{longtable}
}



\paragraph*{Tuple corrette, \texttt{persist(Object pc, Object id, OpCallbacks call)} (BrokerImpl).}

{\footnotesize\setlength{\tabcolsep}{3pt}
\begin{longtable}{>{\raggedright\arraybackslash}p{0.102\textwidth}>{\raggedright\arraybackslash}p{0.074\textwidth}>{\raggedright\arraybackslash}p{0.074\textwidth}>{\raggedright\arraybackslash}p{0.074\textwidth}>{\raggedright\arraybackslash}p{0.176\textwidth}>{\raggedright\arraybackslash}p{0.388\textwidth}}
\toprule
\textbf{pc} & \textbf{id} & \textbf{cascade} & \textbf{call} & \textbf{processArgument} & \textbf{oracolo} \\
\midrule
\endfirsthead
\toprule
\textbf{pc} & \textbf{id} & \textbf{cascade} & \textbf{call} & \textbf{processArgument} & \textbf{oracolo} \\
\midrule
\endhead
\bottomrule
\endfoot
null & null & -- & null & -- & persist ritorna null \\
null & null & -- & non null & azione & persist ritorna null \\
null & null & -- & non null & eccezione & persist ritorna null \\
null & non null & -- & null & -- & persist ritorna null \\
null & non null & -- & non null & azione & persist ritorna null \\
null & non null & -- & non null & eccezione & persist ritorna null \\
non gestito & null & presente & null & -- & state manager non-null, id di default assegnato; istanza correlata (CASCADE\_IMMEDIATE) non cascata, resta non gestita \\
non gestito & null & presente & non null & azione & persist ritorna null \\
non gestito & null & presente & non null & eccezione & eccezione propagata da persist \\
non gestito & null & assente & null & -- & state manager non-null, id di default assegnato \\
non gestito & null & assente & non null & azione & persist ritorna null \\
non gestito & null & assente & non null & eccezione & eccezione propagata da persist \\
non gestito & non null & presente & null & -- & state manager non-null, con l'id specificato assegnato; istanza correlata (CASCADE\_IMMEDIATE) non cascata, resta non gestita \\
non gestito & non null & presente & non null & azione & persist ritorna null \\
non gestito & non null & presente & non null & eccezione & eccezione propagata da persist \\
non gestito & non null & assente & null & -- & state manager non-null, con l'id specificato assegnato \\
non gestito & non null & assente & non null & azione & persist ritorna null \\
non gestito & non null & assente & non null & eccezione & eccezione propagata da persist \\
gestito & null & presente & null & -- & nessuna modifica: pc resta quello gia' gestito, id ignorato \\
gestito & null & presente & non null & azione & nessuna modifica: pc resta quello gia' gestito, id ignorato \\
gestito & null & presente & non null & eccezione & eccezione propagata da persist \\
gestito & null & assente & null & -- & nessuna modifica: pc resta quello gia' gestito, id ignorato \\
gestito & null & assente & non null & azione & nessuna modifica: pc resta quello gia' gestito, id ignorato \\
gestito & null & assente & non null & eccezione & eccezione propagata da persist \\
gestito & non null & presente & null & -- & nessuna modifica: pc resta quello gia' gestito, id ignorato \\
gestito & non null & presente & non null & azione & nessuna modifica: pc resta quello gia' gestito, id ignorato \\
gestito & non null & presente & non null & eccezione & eccezione propagata da persist \\
gestito & non null & assente & null & -- & nessuna modifica: pc resta quello gia' gestito, id ignorato \\
gestito & non null & assente & non null & azione & nessuna modifica: pc resta quello gia' gestito, id ignorato \\
gestito & non null & assente & non null & eccezione & eccezione propagata da persist \\
non gestito & null & presente & non null & azione(ACT\_CASCADE) & persist ritorna null \\
non gestito & null & presente & non null & azione(ACT\_RUN) & state manager non-null, id di default assegnato \\
non gestito & null & assente & non null & azione(ACT\_CASCADE) & persist ritorna null \\
non gestito & null & assente & non null & azione(ACT\_RUN) & state manager non-null, id di default assegnato \\
non gestito & non null & presente & non null & azione(ACT\_CASCADE) & persist ritorna null \\
non gestito & non null & presente & non null & azione(ACT\_RUN) & state manager non-null, con l'id specificato assegnato \\
non gestito & non null & assente & non null & azione(ACT\_CASCADE) & persist ritorna null \\
non gestito & non null & assente & non null & azione(ACT\_RUN) & state manager non-null, con l'id specificato assegnato \\
\bottomrule
\caption{Tuple corrette per \texttt{persist(Object pc, Object id, OpCallbacks call)}.}
\label{tab:persist-tuple-corrette}
\end{longtable}
}



\paragraph*{Tuple di progettazione, \texttt{delete(Object pc, OpCallbacks call)} (BrokerImpl).}

{\footnotesize\setlength{\tabcolsep}{3pt}
\begin{longtable}{>{\raggedright\arraybackslash}p{0.314\textwidth}>{\raggedright\arraybackslash}p{0.285\textwidth}>{\raggedright\arraybackslash}p{0.328\textwidth}}
\toprule
\textbf{pc} & \textbf{call} & \textbf{oracolo} \\
\midrule
\endfirsthead
\toprule
\textbf{pc} & \textbf{call} & \textbf{oracolo} \\
\midrule
\endhead
\bottomrule
\endfoot
null & null & nessun effetto \\
null & azione (ACT\_NONE) & nessuna eccezione \\
null & azione (ACT\_CASCADE) & nessuna eccezione \\
null & azione (ACT\_RUN) & nessuna eccezione \\
null & lancia eccezione & eccezione non propagata \\
non gestito & null & nessuna eccezione \\
non gestito & azione (ACT\_NONE) & nessuna eccezione \\
non gestito & azione (ACT\_CASCADE) & nessuna eccezione \\
non gestito & azione (ACT\_RUN) & nessuna eccezione \\
non gestito & lancia eccezione & eccezione propagata \\
gestito non persistito & null & nessuna eccezione \\
gestito non persistito & azione (ACT\_NONE) & nessuna eccezione \\
gestito non persistito & azione (ACT\_CASCADE) & nessuna eccezione \\
gestito non persistito & azione (ACT\_RUN) & nessuna eccezione \\
gestito non persistito & lancia eccezione & eccezione propagata \\
gestito persistito & null & nessuna eccezione \\
gestito persistito & azione (ACT\_NONE) & nessuna eccezione \\
gestito persistito & azione (ACT\_CASCADE) & nessuna eccezione \\
gestito persistito & azione (ACT\_RUN) & nessuna eccezione \\
gestito persistito & lancia eccezione & eccezione propagata \\
\bottomrule
\caption{Tuple di progettazione per \texttt{delete(Object pc, OpCallbacks call)}.}
\label{tab:delete-tuple}
\end{longtable}
}



\paragraph*{Tuple di progettazione, \texttt{delete(Object pc, OpCallbacks call), approfondimento mutation testing} (BrokerImpl).}

{\footnotesize\setlength{\tabcolsep}{3pt}
\begin{longtable}{>{\raggedright\arraybackslash}p{0.348\textwidth}>{\raggedright\arraybackslash}p{0.155\textwidth}>{\raggedright\arraybackslash}p{0.425\textwidth}}
\toprule
\textbf{pc} & \textbf{call} & \textbf{oracolo} \\
\midrule
\endfirsthead
\toprule
\textbf{pc} & \textbf{call} & \textbf{oracolo} \\
\midrule
\endhead
\bottomrule
\endfoot
gestito e detached & null & UserException detached \\
gestito e detached & non null & UserException detached \\
\bottomrule
\caption{Tuple di progettazione per \texttt{delete(Object pc, OpCallbacks call), approfondimento mutation testing}.}
\label{tab:delete2-tuple}
\end{longtable}
}



\paragraph*{Tuple di progettazione, \texttt{flush(int reason)} (BrokerImpl).}

{\footnotesize\setlength{\tabcolsep}{3pt}
\begin{longtable}{>{\raggedright\arraybackslash}p{0.658\textwidth}>{\raggedright\arraybackslash}p{0.282\textwidth}}
\toprule
\textbf{stato} & \textbf{oracolo} \\
\midrule
\endfirsthead
\toprule
\textbf{stato} & \textbf{oracolo} \\
\midrule
\endhead
\bottomrule
\endfoot
nessuna transazione attiva & nessuna eccezione \\
transazione attiva senza istanze dirty & rollbackOnly false \\
transazione attiva istanze presenti successo & non determinato \\
transazione attiva istanze presenti fallimento & non determinato \\
\bottomrule
\caption{Tuple di progettazione per \texttt{flush(int reason)}.}
\label{tab:flush-tuple}
\end{longtable}
}



\paragraph*{Tuple di progettazione, \texttt{Costruttore} (LoginDialog).}

{\footnotesize\setlength{\tabcolsep}{3pt}
\begin{longtable}{>{\raggedright\arraybackslash}p{0.146\textwidth}>{\raggedright\arraybackslash}p{0.794\textwidth}}
\toprule
\textbf{session} & \textbf{oracolo} \\
\midrule
\endfirsthead
\toprule
\textbf{session} & \textbf{oracolo} \\
\midrule
\endhead
\bottomrule
\endfoot
null & eccezione (dereferenziazione) \\
valida & getWidget()!=null e isShowing()==false \\
\bottomrule
\caption{Tuple di progettazione per \texttt{Costruttore}.}
\label{tab:logindialog-costruttore-tuple}
\end{longtable}
}



\paragraph*{Tuple corrette, \texttt{Costruttore} (LoginDialog).}

{\footnotesize\setlength{\tabcolsep}{3pt}
\begin{longtable}{>{\raggedright\arraybackslash}p{0.146\textwidth}>{\raggedright\arraybackslash}p{0.794\textwidth}}
\toprule
\textbf{session} & \textbf{oracolo} \\
\midrule
\endfirsthead
\toprule
\textbf{session} & \textbf{oracolo} \\
\midrule
\endhead
\bottomrule
\endfoot
null & getWidget()!=null e isShowing()==false \\
valida & getWidget()!=null e isShowing()==false \\
\bottomrule
\caption{Tuple corrette per \texttt{Costruttore}.}
\label{tab:logindialog-costruttore-tuple-corrette}
\end{longtable}
}



\paragraph*{Tuple di progettazione, \texttt{Posizione del popup} (LoginDialog).}

{\footnotesize\setlength{\tabcolsep}{3pt}
\begin{longtable}{>{\raggedright\arraybackslash}p{0.130\textwidth}>{\raggedright\arraybackslash}p{0.130\textwidth}>{\raggedright\arraybackslash}p{0.667\textwidth}}
\toprule
\textbf{left} & \textbf{top} & \textbf{oracolo} \\
\midrule
\endfirsthead
\toprule
\textbf{left} & \textbf{top} & \textbf{oracolo} \\
\midrule
\endhead
\bottomrule
\endfoot
negativo & negativo & getPopupLeft()==left e getPopupTop()==top \\
negativo & zero & getPopupLeft()==left e getPopupTop()==top \\
negativo & positivo & getPopupLeft()==left e getPopupTop()==top \\
zero & negativo & getPopupLeft()==left e getPopupTop()==top \\
zero & zero & getPopupLeft()==left e getPopupTop()==top \\
zero & positivo & getPopupLeft()==left e getPopupTop()==top \\
positivo & negativo & getPopupLeft()==left e getPopupTop()==top \\
positivo & zero & getPopupLeft()==left e getPopupTop()==top \\
positivo & positivo & getPopupLeft()==left e getPopupTop()==top \\
\bottomrule
\caption{Tuple di progettazione per \texttt{Posizione del popup}.}
\label{tab:logindialog-posizione-tuple}
\end{longtable}
}



\paragraph*{Tuple corrette, \texttt{Posizione del popup} (LoginDialog).}

{\footnotesize\setlength{\tabcolsep}{3pt}
\begin{longtable}{>{\raggedright\arraybackslash}p{0.130\textwidth}>{\raggedright\arraybackslash}p{0.130\textwidth}>{\raggedright\arraybackslash}p{0.667\textwidth}}
\toprule
\textbf{left} & \textbf{top} & \textbf{oracolo} \\
\midrule
\endfirsthead
\toprule
\textbf{left} & \textbf{top} & \textbf{oracolo} \\
\midrule
\endhead
\bottomrule
\endfoot
negativo & negativo & getPopupLeft()==left e getPopupTop()==top \\
negativo & zero & getPopupLeft()==left e getPopupTop()==top \\
negativo & positivo & getPopupLeft()==left e getPopupTop()==top \\
zero & negativo & getPopupLeft()==left e getPopupTop()==top \\
zero & zero & getPopupLeft()==left e getPopupTop()==top \\
zero & positivo & getPopupLeft()==left e getPopupTop()==top \\
positivo & negativo & getPopupLeft()==left e getPopupTop()==top \\
positivo & zero & getPopupLeft()==left e getPopupTop()==top \\
positivo & positivo & getPopupLeft()==left e getPopupTop()==top \\
\bottomrule
\caption{Tuple corrette per \texttt{Posizione del popup}.}
\label{tab:logindialog-posizione-tuple-corrette}
\end{longtable}
}



\paragraph*{Tuple di progettazione, \texttt{Titolo del popup} (LoginDialog).}

{\footnotesize\setlength{\tabcolsep}{3pt}
\begin{longtable}{>{\raggedright\arraybackslash}p{0.274\textwidth}>{\raggedright\arraybackslash}p{0.666\textwidth}}
\toprule
\textbf{title} & \textbf{oracolo} \\
\midrule
\endfirsthead
\toprule
\textbf{title} & \textbf{oracolo} \\
\midrule
\endhead
\bottomrule
\endfoot
null & getTitle()==null \\
vuota & getTitle()=="" \\
valida & getTitle()==title \\
\bottomrule
\caption{Tuple di progettazione per \texttt{Titolo del popup}.}
\label{tab:logindialog-title-tuple}
\end{longtable}
}



\paragraph*{Tuple corrette, \texttt{Titolo del popup} (LoginDialog).}

{\footnotesize\setlength{\tabcolsep}{3pt}
\begin{longtable}{>{\raggedright\arraybackslash}p{0.274\textwidth}>{\raggedright\arraybackslash}p{0.666\textwidth}}
\toprule
\textbf{title} & \textbf{oracolo} \\
\midrule
\endfirsthead
\toprule
\textbf{title} & \textbf{oracolo} \\
\midrule
\endhead
\bottomrule
\endfoot
null & getTitle()=="" \\
vuota & getTitle()=="" \\
valida & getTitle()==title \\
\bottomrule
\caption{Tuple corrette per \texttt{Titolo del popup}.}
\label{tab:logindialog-title-tuple-corrette}
\end{longtable}
}



\paragraph*{Tuple di progettazione, \texttt{Proprieta booleane ereditate} (LoginDialog).}

{\footnotesize\setlength{\tabcolsep}{3pt}
\begin{longtable}{>{\raggedright\arraybackslash}p{0.297\textwidth}>{\raggedright\arraybackslash}p{0.130\textwidth}>{\raggedright\arraybackslash}p{0.501\textwidth}}
\toprule
\textbf{proprieta} & \textbf{valore} & \textbf{oracolo} \\
\midrule
\endfirsthead
\toprule
\textbf{proprieta} & \textbf{valore} & \textbf{oracolo} \\
\midrule
\endhead
\bottomrule
\endfoot
Modal & true & isModal()==true \\
Modal & false & isModal()==false \\
AnimationEnabled & true & isAnimationEnabled()==true \\
AnimationEnabled & false & isAnimationEnabled()==false \\
AutoHideEnabled & true & isAutoHideEnabled()==true \\
AutoHideEnabled & false & isAutoHideEnabled()==false \\
GlassEnabled & true & isGlassEnabled()==true \\
GlassEnabled & false & isGlassEnabled()==false \\
Visible & true & isVisible()==true \\
Visible & false & isVisible()==false \\
\bottomrule
\caption{Tuple di progettazione per \texttt{Proprieta booleane ereditate}.}
\label{tab:logindialog-boolprop-tuple}
\end{longtable}
}



\paragraph*{Tuple di progettazione, \texttt{Ciclo di vita} (LoginDialog).}

{\footnotesize\setlength{\tabcolsep}{3pt}
\begin{longtable}{>{\raggedright\arraybackslash}p{0.337\textwidth}>{\raggedright\arraybackslash}p{0.211\textwidth}>{\raggedright\arraybackslash}p{0.379\textwidth}}
\toprule
\textbf{stato\_precedente} & \textbf{operazione} & \textbf{oracolo} \\
\midrule
\endfirsthead
\toprule
\textbf{stato\_precedente} & \textbf{operazione} & \textbf{oracolo} \\
\midrule
\endhead
\bottomrule
\endfoot
non mostrato & show() & isShowing()==true \\
non mostrato & hide() & isShowing()==false \\
non mostrato & hide(true) & isShowing()==false \\
non mostrato & center() & isShowing()==true \\
mostrato & show() & isShowing()==true \\
mostrato & hide() & isShowing()==false \\
mostrato & hide(true) & isShowing()==false \\
mostrato & center() & isShowing()==true \\
\bottomrule
\caption{Tuple di progettazione per \texttt{Ciclo di vita}.}
\label{tab:logindialog-lifecycle-tuple}
\end{longtable}
}



\paragraph*{Tuple di progettazione, \texttt{Classe GetServerURI} (LoginDialog).}

{\footnotesize\setlength{\tabcolsep}{3pt}
\begin{longtable}{>{\raggedright\arraybackslash}p{0.141\textwidth}>{\raggedright\arraybackslash}p{0.346\textwidth}>{\raggedright\arraybackslash}p{0.440\textwidth}}
\toprule
\textbf{metodo} & \textbf{parametro} & \textbf{oracolo} \\
\midrule
\endfirsthead
\toprule
\textbf{metodo} & \textbf{parametro} & \textbf{oracolo} \\
\midrule
\endhead
\bottomrule
\endfoot
onFailure & throwable qualunque & handleError(caught) invocato \\
onSuccess & null & eccezione sincrona \\
onSuccess & vuota & eccezione sincrona \\
onSuccess & non valida & eccezione sincrona \\
onSuccess & valida e corretta & propaga a init \\
onSuccess & valida ma non corretta & va in onFailure \\
\bottomrule
\caption{Tuple di progettazione per \texttt{Classe GetServerURI}.}
\label{tab:logindialog-getserveruri-tuple}
\end{longtable}
}



\paragraph*{Tuple corrette, \texttt{Classe GetServerURI} (LoginDialog).}

{\footnotesize\setlength{\tabcolsep}{3pt}
\begin{longtable}{>{\raggedright\arraybackslash}p{0.141\textwidth}>{\raggedright\arraybackslash}p{0.346\textwidth}>{\raggedright\arraybackslash}p{0.440\textwidth}}
\toprule
\textbf{metodo} & \textbf{parametro} & \textbf{oracolo} \\
\midrule
\endfirsthead
\toprule
\textbf{metodo} & \textbf{parametro} & \textbf{oracolo} \\
\midrule
\endhead
\bottomrule
\endfoot
onFailure & throwable qualunque & handleError(caught) invocato \\
onSuccess & null & getStocks() sempre invocato \\
onSuccess & vuota & getStocks() sempre invocato \\
onSuccess & non valida & getStocks() sempre invocato \\
onSuccess & valida e corretta & getStocks() sempre invocato \\
onSuccess & valida ma non corretta & getStocks() sempre invocato \\
\bottomrule
\caption{Tuple corrette per \texttt{Classe GetServerURI}.}
\label{tab:logindialog-getserveruri-tuple-corrette}
\end{longtable}
}



\paragraph*{Tuple di progettazione, \texttt{Classe InitializeStocks} (LoginDialog).}

{\footnotesize\setlength{\tabcolsep}{3pt}
\begin{longtable}{>{\raggedright\arraybackslash}p{0.149\textwidth}>{\raggedright\arraybackslash}p{0.315\textwidth}>{\raggedright\arraybackslash}p{0.464\textwidth}}
\toprule
\textbf{metodo} & \textbf{parametro} & \textbf{oracolo} \\
\midrule
\endfirsthead
\toprule
\textbf{metodo} & \textbf{parametro} & \textbf{oracolo} \\
\midrule
\endhead
\bottomrule
\endfoot
onFailure & throwable qualunque & handleError(caught) invocato \\
onSuccess & null & chiama init(...) \\
onSuccess & vuota & chiama init(...) \\
onSuccess & con elementi & chiama init(...) \\
\bottomrule
\caption{Tuple di progettazione per \texttt{Classe InitializeStocks}.}
\label{tab:logindialog-initializestocks-tuple}
\end{longtable}
}


